\chapter{Induction and the Proof Kernel}\label{ch:kernel}
\begingroup\sloppy

This chapter specifies how ArchytasZX proves an equation between two diagrams for every value of a !-box multiplicity, and how it turns every proof it finds into a theorem of its \term{proof kernel}. Three modules of \module{archytaszx.semantics} implement it:
\begin{itemize}
\item \module{archytaszx.semantics.induction}: induction over one multiplicity: the choice of the index, the obligation of six diagrams, the base case, and the ladder of discharge tiers that settles the step case (\cref{sec:kernel-induction,sec:kernel-tiers,sec:kernel-regions,sec:kernel-expose,sec:kernel-chain,sec:kernel-tree});
\item \module{archytaszx.semantics.induction_kernel}: the kernel theorem of an induction-rung \texttt{EQUAL}, rebuilt from the data the settling tier recorded (\cref{sec:kernel-indthm});
\item \module{archytaszx.semantics.kernel}: the proof kernel itself: the \py!Theorem! value, its inference rules, the reading of certificates and decisions as theorems, the refusals, and the top-level dispatch of the export to Lean~4 (\cref{sec:kernel-kernel,sec:kernel-rules,sec:kernel-check,sec:kernel-refusals,sec:kernel-lean}).
\end{itemize}

\paragraph{Position in the architecture.} The decision procedure \py!decide_equal! of \cref{ch:proofs} runs a ladder of rungs. Its induction rung decides the base case by a nested ladder, calls \py!prove_by_induction! (\cref{sec:kernel-ladder}) with the step ladder for the step case, and, when the step case is proved from base~$0$, calls \py!induction_theorem! (\cref{sec:kernel-indthm}) to build the kernel theorem that the decision carries as its \py!induction_proof!. The certification step of the decision procedure and the shell read every other proof (step-sequence certificates, proof certificates of the goal-directed search, induction proofs, restricted proofs, and decisions) through \py!kernel.check! and \py!kernel.from_decision! (\cref{sec:kernel-check}), and export the resulting theorem with \py!kernel.lean! (\cref{sec:kernel-lean}). The detailed Lean emission over tensor networks is specified in \cref{ch:lean-export-net}; the emission of path-sum, symbolic-sum, refutation, restriction, and linked-diagram proofs in \cref{ch:lean-export}; the Lean library they target in \cref{ch:lean-core,ch:lean-boxes,ch:lean-pathsum,ch:lean-sym}.

\paragraph{Design.} The kernel follows the design of Edinburgh LCF~\cite{gordon1979lcf}: a theorem is a value of an abstract type whose constructor only the inference rules of the kernel can call, so every theorem in the system is the conclusion of a finite tree of inference-rule applications, each of which re-checked its premises and side conditions when it ran. Induction over !-box multiplicities follows the first-order logic for string diagrams of Kissinger, Merry, and Soloviev~\cite{kissinger2015firstorder}; !-boxes are those of pattern graphs~\cite{dixon2009graphical,kissinger2014pattern}; the diagrams are those of the qufinite ZXW calculus~\cite{poor2023zxw,wang2023qufinitezxw}, and a diagram with a discard denotes a completely positive map~\cite{selinger2007cpm,carette2019mixed}.

\section{Notation and Shared Notions}\label{sec:kernel-notation}

The diagram model (nodes, ports, wires, boundary lists, !-boxes and their extensions) is that of \cref{ch:diagrams}. This section fixes the relations and operations the three modules share.

\begin{definition}[Boundary entry]\label{def:kernel-entry}
A \term{boundary entry} (in the code, a \emph{slot}) is one element of a diagram's boundary input list or boundary output list: a \texttt{PortRef(\allowbreak{}node\_\allowbreak{}id,\allowbreak{} direction,\allowbreak{} index)} naming one port of one node. The \emph{place} of an entry is the pair (list, position), list $0$ for inputs and $1$ for outputs.
\end{definition}

\begin{definition}[Identity, isomorphism, view isomorphism]\label{def:kernel-same}
For diagrams $a$ and $b$:
\begin{enumerate}
\item $a \equiv b$ (\emph{identical}, \py!kernel._same(a, b)!) when \texttt{compare\_\allowbreak{}structure(\allowbreak{}a,\allowbreak{} b).\allowbreak{}identical} holds: the two diagrams agree node for node and wire for wire under the identity map on node ids, with equal boundary lists, scalars, and !-boxes (\cref{ch:proofs}, certificates).
\item $a \cong b$ (\emph{isomorphic}) when \py!isomorphic(a, b)! of \module{archytaszx.diagram.compare} finds a node bijection preserving generators, port dimensions, phases, wires, boundary lists, and !-boxes; with \py!symmetric_legs=True! the legs of Z and X spiders may be permuted (\cref{ch:diagrams}).
\item $a \cong_v b$ (\emph{view-isomorphic}) when \texttt{views\_\allowbreak{}isomorphic(\allowbreak{}comparison\_\allowbreak{}view(\allowbreak{}a),\allowbreak{} comparison\_\allowbreak{}view(\allowbreak{}b))} holds: the comparison view puts Z and X spider legs in canonical order, empties the parameter environment, and simplifies the scalar (for a mixed diagram, the scalar $c\bar c$); the two views are compared by \py!isomorphic! with symmetric legs (\cref{ch:rewriting}).
\end{enumerate}
\end{definition}

\begin{definition}[Bare diagram, footprint, interface]\label{def:kernel-bare}
\py!_bare(D)! is a copy of $D$ with scalar one and an empty parameter environment. The \term{footprint} of a !-box $b$ (\py!_box_footprint!) is its node scope together with the nodes of its port-scope entries. The \term{interface} of $D$ (\py!_interface!) is the pair of tuples of the dimensions of its boundary inputs and of its boundary outputs, in order; \py!_ref_dim(D, r)! is the dimension of the port \py!r!.
\end{definition}

\begin{definition}[Substitution and peeling]\label{def:kernel-peel}
For a diagram $D$, a multiplicity symbol $n$, and a multiplicity expression $e$, $D[n:=e]$ is \py!rename_multiplicity(D, n, e)! (\cref{sec:kernel-rename}). $D|_{n=0}$ is \py!induction_base(D, n)!: $D$ with $n$ at $0$ in multiplicities, scalar, copy scalars, phases, and port dimensions (\cref{ch:proofs}). Peeling one copy off a !-box of multiplicity $m$ (\py!peel_one!, \cref{ch:diagrams}) leaves the !-box at $m-1$ beside one bare copy laid out after the !-box's own boundary block; it requires a multiplicity whose constant term is at least one, and with \py!last=True! a chain that peels its first copy (\py!peels_first!: every link runs from an input $p$ to an output $q$) loses its last copy instead. The \term{joint peel} \texttt{peel\_\allowbreak{}owners(\allowbreak{}D,\allowbreak{} n,\allowbreak{} last=.\allowbreak{}.\allowbreak{}.\allowbreak{},\allowbreak{} partial=.\allowbreak{}.\allowbreak{}.\allowbreak{})} peels one copy off every !-box of \py!index_owners(D, n)! (the !-boxes whose multiplicity reads $n$, by id; with more than one, the pair !-boxes of another such !-box left out) in id order and joins the copies; with \py!partial=True! only the !-boxes whose multiplicity has a positive constant term peel. It raises \py!BangBoxError! when no !-box peels or one lies inside another. Peeling preserves the denotation at every value of the multiplicity at which both sides instantiate.
\end{definition}

\paragraph{Helpers of the prover.} The kernel and the induction-theorem builder use six functions of \module{archytaszx.semantics.prove} (\cref{ch:proofs}). With $n$ the index and $k$ the step symbol:
\begin{itemize}
\item \texttt{induction\_\allowbreak{}hypothesis(\allowbreak{}D,\allowbreak{} n,\allowbreak{} k)} is $D[n:=k]$ with $n$ replaced by $k$ in the scalar as well (\py!_with_scalar_index!).
\item \py!induction_renamed(D, n, k)! is $D[n:=k+1]$, scalar included, unpeeled.
\item \py!induction_successor(D, n, k)! is $D[n:=k+1]$ (scalar included) with one copy peeled off its one !-box of \py!index_owners(., k)!; unpeeled when that !-box's multiplicity has no positive constant term (as $b^{k+1}$) and it is no tree !-box; with several owners, none nested in another and each with a positive constant term, the joint peel; when only some have a positive constant term, the partial joint peel; unpeeled when no !-box carries $k$; \py!None! when a tree !-box is among several owners of which not all peel, when none peels, or when the joint peel raises.
\item \py!peels_shifted(D, n)! holds when an owner of $n$ peels its first copy and reads a copy index.
\item \py!induction_offset(D, n, s)! makes the one !-box carrying $n$ read its copy index as $\text{index}+s$ (returns $D$ unchanged when that !-box reads none); with several owners, each that peels its first copy and reads an index is so shifted; it raises \py!InductionError! when no !-box carries $n$.
\item \py!offset_symbol(*diagrams)! is the first of \code!s!, \code!s1!, \code!s2!, \ldots{} named in none of the diagrams (free symbols and copy indices), as a SymPy symbol with \texttt{integer=True,\allowbreak{} nonnegative=True}.
\end{itemize}

\paragraph{Errors caught as one family.} \py!_INDUCTION_ERRORS! of \module{archytaszx.semantics.prove} is the tuple of the checking errors of the prover (\py!BangBoxError!, \py!CheckError!, \py!DenoteError!, \py!ContractError!, \py!PhaseError!, \py!RewriteError!, \py!ScalarError!, \py!ValidateError!, and the other errors \cref{ch:proofs} lists) extended by \py!InductionError!. Every kernel rule that rebuilds diagrams converts an exception of this family into a \py!KernelError!.

\section{Induction over Multiplicities}\label{sec:kernel-induction}

\module{archytaszx.semantics.induction} proves an equation $L=R$ for every natural value of one multiplicity symbol $n$, the \term{induction index}. It builds a base case at $n=b$ ($b\in\{0,1\}$) and a hypothesis--successor pair at $n=k$ and $n=k+1$ for a fresh \term{step symbol} $k$, then runs a ladder of \term{discharge tiers} on the step case until one settles it or all decline. The module is deterministic: every iteration over nodes, wires, and !-boxes runs in sorted id order or in the order a documented function returns.

\subsection{Errors, verdicts, and constants}

\begin{signature}
\texttt{class InductionError(\allowbreak{}Exception)} --- the base class of every error the module raises.\\
\texttt{class InductionDomainError(\allowbreak{}InductionError)} --- a value outside the mathematical domain the module accepts (a multiplicity of unsupported shape, a negative base, a replacement whose interface does not fit).\\
\texttt{class InductionGrammarError(\allowbreak{}InductionError)} --- a malformed request, independent of any diagram's concreteness (no !-box, an index that names no live !-box, a symbol left uninstantiated).
\end{signature}

\begin{signature}
\py!class Verdict(enum.Enum)!: \texttt{PROVED\_\allowbreak{}UNIFORM = "proved\_\allowbreak{}uniform"}, \texttt{PROVED\_\allowbreak{}INDUCTION = "proved\_\allowbreak{}induction"}, \texttt{SCHEMA\_\allowbreak{}CHECKED = "schema\_\allowbreak{}checked"}, \py!REFUTED = "refuted"!, \texttt{INCONCLUSIVE = "inconclusive"}.

\texttt{class StepDischarge(\allowbreak{}enum.\allowbreak{}Enum)}: \texttt{UNIFORM\_\allowbreak{}REWRITE = "uniform\_\allowbreak{}rewrite"}, \texttt{INDUCTION\_\allowbreak{}REWRITE = "induction\_\allowbreak{}rewrite"}, \texttt{ORACLE\_\allowbreak{}WINDOW = "oracle\_\allowbreak{}window"}, \texttt{SYMBOLIC\_\allowbreak{}CONTRACTION = "symbolic\_\allowbreak{}contraction"}, \texttt{CHAIN\_\allowbreak{}WINDOW = "chain\_\allowbreak{}window"}, \py!TREE_WINDOW = "tree_window"!.
\end{signature}

\begin{table}[H]
\centering\small
\begin{tabularx}{\linewidth}{@{}l>{\raggedright\arraybackslash}p{0.2\linewidth}X@{}}
\toprule
Constant & Value & Meaning \\
\midrule
\texttt{DEFAULT\_WINDOW\_WIDTH} & $4$ & multiplicities the oracle window compares \\
\texttt{DEFAULT\_MAX\_STEPS} & $64$ & rewrite firings and symbolic-contraction step budget \\
\texttt{DEFAULT\_STEP\_SYMBOL} & \texttt{"k"} & stem of the step symbol \\
\texttt{DEFAULT\_RULES} & \texttt{(SPIDER\_FUSION,)} & rules of the rewriting tiers when none are given \\
\texttt{\_MAX\_CUT\_PLACEMENTS} & $2000$ & placements of cut ends per boundary list \\
\texttt{\_MAX\_GREEDY\_STEPS} & $16$ & exposing steps of one greedy run \\
\texttt{\_MAX\_LEMMA\_DEPTH} & $2$ & nesting of lemma decisions \\
\texttt{max\_regions} default & $12$ & connected regions combined by \texttt{find\_hypothesis\_region} \\
\bottomrule
\end{tabularx}
\caption{Constants of \texttt{archytaszx.semantics.induction}.}\label{tab:kernel-ind-constants}
\end{table}

The \term{default ladder} \py!DEFAULT_LADDER! is, in order, \texttt{UNIFORM\_REWRITE}, \texttt{INDUCTION\_REWRITE}, \texttt{SYMBOLIC\_CONTRACTION}, \texttt{CHAIN\_WINDOW}, \texttt{TREE\_WINDOW}, \texttt{HYPOTHESIS\_REWRITE}, and \texttt{ORACLE\_WINDOW}.
The induction rung of \py!decide_equal! passes its own \term{step ladder} \py!_STEP_LADDER!, the default ladder without \texttt{ORACLE\_WINDOW}, so a decision never rests on a finite window. The verdict a settling tier gives is fixed by the table \py!_VERDICT!: \texttt{UNIFORM\_REWRITE} gives \texttt{PROVED\_UNIFORM}; \texttt{INDUCTION\_REWRITE}, \texttt{SYMBOLIC\_CONTRACTION}, \texttt{CHAIN\_WINDOW}, \texttt{TREE\_WINDOW}, and \texttt{HYPOTHESIS\_REWRITE} give \texttt{PROVED\_INDUCTION}; \texttt{ORACLE\_WINDOW} gives \texttt{SCHEMA\_CHECKED}, a finite check that is not a proof.

\subsection{Records}

\begin{definition}[Induction obligation]\label{def:kernel-obligation}
\py!InductionObligation! (frozen, slotted, compared by identity) holds the six diagrams an induction over one index is discharged against, and the bindings they were built under:
\begin{center}\small
\begin{tabularx}{\linewidth}{@{}>{\raggedright\arraybackslash}p{0.33\linewidth}>{\raggedright\arraybackslash}p{0.2\linewidth}X@{}}
\toprule
Field & Type & Content \\
\midrule
\texttt{index} & \texttt{str} & the induction index $n$ \\
\texttt{base\_value} & \texttt{int} & the base $b\in\{0,1\}$ \\
\texttt{step\_symbol} & \texttt{str} & the fresh step symbol $k$ \\
\texttt{left\_at\_base}, \texttt{right\_at\_base} & \texttt{Diagram} & $L$ and $R$ at $n=b$ \\
\texttt{left\_at\_k}, \texttt{right\_at\_k} & \texttt{Diagram} & the hypothesis diagrams $L[n:=k]$, $R[n:=k]$ \\
\texttt{left\_at\_successor}, \texttt{right\_at\_successor} & \texttt{Diagram} & the successors $L[n:=k+1]$, $R[n:=k+1]$ \\
\texttt{witness} & \texttt{Mapping[str, int|Rational]} & values of held non-multiplicity symbols (default empty) \\
\texttt{held\_symbolic} & \texttt{frozenset[str]} & symbols the symbolic tiers keep formal (default empty) \\
\texttt{base\_proved} & \texttt{bool} & the base case is taken as proved (default \texttt{False}) \\
\bottomrule
\end{tabularx}
\end{center}
\end{definition}

\begin{definition}[Tier outcome]\label{def:kernel-tier}
\py!TierOutcome! (frozen, slotted, compared by value) is the result of one discharge tier: \py!discharge: StepDischarge!, \py!settled: bool!, \py!reason: str!, \texttt{comparisons: tuple[ComparisonResult,\allowbreak{} .\allowbreak{}.\allowbreak{}.\allowbreak{}] = (\allowbreak{})}, \texttt{steps: tuple[RewriteStep,\allowbreak{} .\allowbreak{}.\allowbreak{}.\allowbreak{}] = (\allowbreak{})} (the rewrite steps of a rewriting tier), \texttt{counterexample: int | None = None} (a multiplicity at which the oracle window found a mismatch), and, for a settled peeling tier, \py!left_peeled!, \py!left_copy!, \py!right_peeled!, \py!right_copy! (each peeled successor and the node ids of its peeled copy; \py!None! otherwise).
\end{definition}

\begin{definition}[Induction result]\label{def:kernel-result}
\py!InductionResult! (frozen, slotted, compared by identity) is a verdict on one equation over one index. Fields: \py!verdict!, \py!index!, \py!base_value!, \py!reason!, \py!obligation!, \py!witness! (default empty), \texttt{base: ComparisonResult | None} (the base comparison), \texttt{tiers: tuple[TierOutcome,\allowbreak{} .\allowbreak{}.\allowbreak{}.\allowbreak{}]} (every tier run, in order), \py!comparison! (the deciding comparison of a refutation), \py!counterexample: int | None!, \texttt{derivation: Derivation | None} (the induction certificate's derivation), and, for a restricted run, \py!restriction!, \texttt{family: InductionResult | None} (the unrestricted run), \py!family_holds: bool! (that run proved the claim), \py!restricted_pair! (the substituted pair the run inducted on otherwise), and \py!counterexample_original! (a counterexample in the original symbols). Its properties:
\begin{itemize}
\item \py!proved! is true when \py!verdict! is \texttt{PROVED\_UNIFORM} or \texttt{PROVED\_INDUCTION};
\item \py!discharge! is the discharge of the last settled tier, \py!None! when none settled;
\item \py!held_symbolic! is \py!obligation.held_symbolic! when \py!discharge! is \texttt{UNIFORM\_REWRITE}, \texttt{INDUCTION\_REWRITE}, \texttt{SYMBOLIC\_CONTRACTION}, \texttt{CHAIN\_WINDOW}, \texttt{TREE\_WINDOW}, or \texttt{HYPOTHESIS\_REWRITE}, and the empty set otherwise.
\end{itemize}
\end{definition}

\subsection{Free symbols and the index}\label{sec:kernel-index}

\py!_all_free_symbols(D)! is the set of names free in the scalar of $D$, in every port dimension and phase of every node, in every !-box multiplicity, and in every copy scalar, less the copy-index symbols of $D$ (\py!index_symbols!). \texttt{\_\allowbreak{}fresh\_\allowbreak{}symbol\_\allowbreak{}name(\allowbreak{}stem,\allowbreak{} avoid)} returns the first of \code!stem!, \code!stem1!, \code!stem2!, \ldots{} not in \code!avoid!. Write $\mathrm{fs}(D)$ for this set. \texttt{\_\allowbreak{}fresh\_\allowbreak{}step\_\allowbreak{}symbol(\allowbreak{}L,\allowbreak{} R,\allowbreak{} n,\allowbreak{} k)} freshens $k$ against $\mathrm{fs}(L)\cup\mathrm{fs}(R)\cup\{n\}$.

\begin{signature}
\texttt{choose\_\allowbreak{}index(\allowbreak{}left,\allowbreak{} right,\allowbreak{} *,\allowbreak{} index=None) -> str}

The lexicographically least multiplicity symbol free in either diagram (\py!min! of \texttt{free\_\allowbreak{}mult\_\allowbreak{}symbols(\allowbreak{}left) | free\_\allowbreak{}mult\_\allowbreak{}symbols(\allowbreak{}right)}), or \code!index! when given. Raises \py!InductionGrammarError! with \texttt{diagram carries no bang box; there is no index to induct on} when no multiplicity symbol is free, and with \texttt{no live bang box has multiplicity symbol 'n'} when \code!index! is not one of them.
\end{signature}

\py!_offset(expr, sym)! is the integer $c$ when $\mathit{expr}-\mathit{sym}$ expands to an integer, else \py!None!. \py!_least_offset(L, R, n)! is the least $c$ over all multiplicities $n+c$ of both diagrams' !-boxes, with $0$ included in the minimum.

\subsection{Renaming a multiplicity}\label{sec:kernel-rename}

\begin{algorithm}[\texttt{rename\_\allowbreak{}multiplicity}]\label{alg:kernel-rename}
Input: a diagram, a multiplicity name $n$, and a \py!Mult! $e$. Output: a new diagram.
\begin{enumerate}
\item Copy the diagram.
\item For each !-box in id order whose multiplicity reads $n$: when the multiplicity is the bare symbol $n$, set it to $e$. Otherwise let $x$ be the multiplicity as a SymPy expression and $\sigma$ its symbol named $n$. When $x$ is neither $\sigma+c$ for an integer $c$ nor a power $\beta^{\sigma}$ with integer base, raise \py!InductionDomainError! \texttt{induction requires a multiplicity n + c (\allowbreak{}c an integer) or a power b\^{}n,\allowbreak{} got M}. Else set the multiplicity to \texttt{Mult.\allowbreak{}\_\allowbreak{}from\_\allowbreak{}expr(\allowbreak{}powsimp(\allowbreak{}expand(\allowbreak{}x[$\sigma$ := e]),\allowbreak{} force=True))}; a \py!BangBoxError! there becomes \py!InductionDomainError! \texttt{multiplicity M does not take n = e: E}.
\item When the scalar reads $n$, substitute $e$ for $n$ in it; likewise in every copy scalar that reads $n$.
\item Substitute $e$ for $n$ in the phases that read it (\py!substitute_index_phases!).
\item For every node with a port dimension reading $n$, reset all its port dimensions to \py!_dim_at(dim, n, e)!.
\item Drop a parameter binding of $n$.
\end{enumerate}
\end{algorithm}

\texttt{hypothesis\_\allowbreak{}diagram(\allowbreak{}D,\allowbreak{} n,\allowbreak{} *,\allowbreak{} step\_\allowbreak{}symbol="k")} is $D[n:=k]$. \texttt{successor\_\allowbreak{}diagram(\allowbreak{}D,\allowbreak{} n,\allowbreak{} *,\allowbreak{} step\_\allowbreak{}symbol="k")} freshens $k$ against $D$ and $n$ and returns $D[n:=k+1]$.

\subsection{The obligation}\label{sec:kernel-obligation}

\begin{signature}
\texttt{choose\_\allowbreak{}base\_\allowbreak{}value(\allowbreak{}left,\allowbreak{} right,\allowbreak{} index,\allowbreak{} witness,\allowbreak{} *,\allowbreak{} mode=EXACT,\allowbreak{} tolerance=DEFAULT\_\allowbreak{}TOLERANCE,\allowbreak{} max\_\allowbreak{}elements=DEFAULT\_\allowbreak{}MAX\_\allowbreak{}ELEMENTS) -> int}

For the candidates $0$ then $1$: call the oracle's \texttt{compare(\allowbreak{}left,\allowbreak{} right,\allowbreak{} \{**witness,\allowbreak{} index: c\},\allowbreak{} .\allowbreak{}.\allowbreak{}.\allowbreak{})} (\cref{ch:proofs}); the first candidate for which it returns (whether or not the tensors match) is the base. A \py!BangBoxError!, \py!CheckError!, \py!DenoteError!, \py!ContractError!, or \py!ValidateError! rejects the candidate. When both are rejected, raise \py!InductionDomainError! \texttt{neither multiplicity 0 nor 1 yields an evaluable instance; at 0: E0; at 1: E1}.
\end{signature}

\begin{algorithm}[\texttt{build\_\allowbreak{}obligation}]\label{alg:kernel-obligation}
\begin{enumerate}
\item $n \gets$ \texttt{choose\_\allowbreak{}index(\allowbreak{}left,\allowbreak{} right,\allowbreak{} index=index)}.
\item $H \gets (\text{free multiplicities of both})\setminus\{n\}$, the \emph{held} multiplicities.
\item $Q \gets (\mathrm{fs}(L)\cup\mathrm{fs}(R))\setminus\{n\}\setminus H$, the symbols that need a value.
\item When $Q$ has a symbol not named in \code!witness! and \code!base_proved! is false, raise \py!InductionGrammarError! \texttt{induction over 'n' leaves symbol(\allowbreak{}s) uninstantiated: [.\allowbreak{}.\allowbreak{}.\allowbreak{}]; name them in the witness to hold them at a value}.
\item When \code!base! is given and is a \py!bool! or not in $\{0,1\}$, raise \py!InductionDomainError! \texttt{base index must be 0 or 1,\allowbreak{} got B}.
\item When \code!base_proved! is true and \code!base! is \py!None!, raise \py!InductionGrammarError! \texttt{a proved base needs its base value}.
\item When \code!base! is given and \code!base_proved! is false, run the oracle comparison at $\{\ldots\text{witness}, n: b\}$; an error of the list above becomes \py!InductionDomainError! \texttt{base b has no evaluable instance: E}.
\item When \code!base! is \py!None!, $b\gets$ \py!choose_base_value!.
\item When $b+c_{\min}<0$, for $c_{\min}$ the value of \py!_least_offset(L, R, n)!, raise \py!InductionDomainError! \texttt{a box of multiplicity n + c is negative at the base n = b}.
\item $k\gets$ \texttt{\_\allowbreak{}fresh\_\allowbreak{}step\_\allowbreak{}symbol(\allowbreak{}L,\allowbreak{} R,\allowbreak{} n,\allowbreak{} step\_\allowbreak{}symbol)}.
\item Build the obligation: the base diagrams are \texttt{\_\allowbreak{}at\_\allowbreak{}base(\allowbreak{}D,\allowbreak{} n,\allowbreak{} b,\allowbreak{} witness,\allowbreak{} base\_\allowbreak{}proved)} (with \code!base_proved!, only $n$ instantiated through \py!_instantiate_index! of \cref{ch:proofs}; otherwise the oracle's \texttt{instantiate(\allowbreak{}D,\allowbreak{} \{**witness,\allowbreak{} n: b\})}), the hypothesis diagrams \texttt{hypothesis\_\allowbreak{}diagram(\allowbreak{}D,\allowbreak{} n,\allowbreak{} step\_\allowbreak{}symbol=k)}, the successors \texttt{successor\_\allowbreak{}diagram(\allowbreak{}D,\allowbreak{} n,\allowbreak{} step\_\allowbreak{}symbol=k)}; \code!held_symbolic! is $Q\cup H$.
\item Validate each of the six diagrams with \py!validate_or_raise!; a \py!ValidateError! propagates.
\end{enumerate}
\end{algorithm}

\begin{signature}
\texttt{discharge\_\allowbreak{}base(\allowbreak{}obligation,\allowbreak{} *,\allowbreak{} mode,\allowbreak{} tolerance,\allowbreak{} max\_\allowbreak{}elements) -> ComparisonResult}

Score both base diagrams under the empty assignment (\texttt{score(\allowbreak{}D,\allowbreak{} \{\},\allowbreak{} max\_\allowbreak{}elements=.\allowbreak{}.\allowbreak{}.\allowbreak{})}) and compare their tensors with \texttt{compare\_\allowbreak{}tensors(\allowbreak{}.\allowbreak{}.\allowbreak{}.\allowbreak{},\allowbreak{} mode,\allowbreak{} tolerance)} (\cref{ch:proofs}).
\end{signature}

\subsection{The discharge ladder}\label{sec:kernel-ladder}

\begin{signature}
\texttt{prove\_\allowbreak{}by\_\allowbreak{}induction(\allowbreak{}left,\allowbreak{} right,\allowbreak{} *,\allowbreak{} index=None,\allowbreak{} base=None,\allowbreak{} witness=\{\},\allowbreak{} ladder=DEFAULT\_\allowbreak{}LADDER,\allowbreak{} rules=None,\allowbreak{} step\_\allowbreak{}symbol="k",\allowbreak{} width=4,\allowbreak{} max\_\allowbreak{}steps=64,\allowbreak{} certify\_\allowbreak{}proof=True,\allowbreak{} mode=EXACT,\allowbreak{} tolerance=DEFAULT\_\allowbreak{}TOLERANCE,\allowbreak{} max\_\allowbreak{}elements=DEFAULT\_\allowbreak{}MAX\_\allowbreak{}ELEMENTS,\allowbreak{} restrict=None,\allowbreak{} family\_\allowbreak{}first=True,\allowbreak{} base\_\allowbreak{}proved=False) -> InductionResult}
\end{signature}

\begin{algorithm}[\texttt{prove\_\allowbreak{}by\_\allowbreak{}induction}]\label{alg:kernel-prove-by-induction}
\begin{enumerate}
\item When \code!family_first! is not a \py!bool!, raise \py!InductionGrammarError! \texttt{family\_\allowbreak{}first must be a bool,\allowbreak{} got V}. When \code!base_proved! is true and \code!restrict! is given, raise \py!InductionGrammarError! \texttt{a proved base is not taken under a restriction}.
\item When \code!restrict! is given, return the restricted induction (\cref{alg:kernel-restricted}) with the options \code!base!, \code!ladder!, \code!rules!, \code!step_symbol!, \code!width!, \code!max_steps!, \code!certify_proof!, \code!mode!, \code!tolerance!, \code!max_elements!.
\item Build the obligation (\cref{alg:kernel-obligation}).
\item Unless \code!base_proved!, compare the base diagrams (\py!discharge_base!). A mismatch returns \texttt{REFUTED} with reason \texttt{the base case at n=b does not hold: R}, \py!counterexample! $=b$, and the comparison in \py!base! and \py!comparison!.
\item For each discharge of \code!ladder!, in order: run the tier through \py!_run_tier!, append its outcome to the tier list, and
  \begin{enumerate}
  \item when the outcome carries a counterexample, return \texttt{REFUTED} with the outcome's reason, its last comparison, and that counterexample;
  \item when the outcome is settled, take the verdict \py!_VERDICT[discharge]!; when \code!certify_proof! is true and the verdict is a proved one, build the derivation (\cref{alg:kernel-derivation}); when that fails, the reason becomes \texttt{R; the verdict stands and no certificate was built: W}. Return the result.
  \end{enumerate}
\item Return \texttt{INCONCLUSIVE} with reason \texttt{no discharge tier settled the step case}.
\end{enumerate}
\end{algorithm}

\paragraph{Running one tier.} \texttt{\_\allowbreak{}run\_\allowbreak{}tier(\allowbreak{}discharge,\allowbreak{} obligation,\allowbreak{} *,\allowbreak{} rules,\allowbreak{} max\_\allowbreak{}steps,\allowbreak{} width,\allowbreak{} mode,\allowbreak{} tolerance,\allowbreak{} max\_\allowbreak{}elements)} dispatches through \py!_DISPATCH! (\cref{tab:kernel-tiers}). The two rewriting tiers receive \code!rules! and \code!max_steps!; the oracle window receives \code!width!, \code!mode!, \code!tolerance!, and \code!max_elements!; the other tiers receive \code!max_steps!. An \py!InductionGrammarError! becomes the unsettled outcome \texttt{tier refused the obligation: E}; an \py!InductionDomainError!, \py!RewriteDomainError!, \texttt{SymbolicContractionDomainError}, or \py!ContractDomainError! becomes \code!tier hit a domain error: E!. Every other exception propagates. A direct call of a tier function still raises.

\begin{table}[H]
\centering\small
\begin{tabularx}{\linewidth}{@{}>{\raggedright\arraybackslash}p{0.27\linewidth}>{\raggedright\arraybackslash}p{0.33\linewidth}X@{}}
\toprule
Discharge & Function & Method \\
\midrule
\texttt{UNIFORM\_REWRITE} & \texttt{discharge\_uniform\_rewrite} & rewrite the hypothesis diagrams into each other, the multiplicity never instantiated (\cref{sec:kernel-tier-uniform}) \\
\texttt{INDUCTION\_REWRITE} & \texttt{discharge\_induction\_rewrite} & peel one copy off each successor; residuals are the hypothesis diagrams; compare the copies (\cref{sec:kernel-tier-peel}) \\
\texttt{SYMBOLIC\_CONTRACTION} & \texttt{discharge\_symbolic\_\allowbreak contraction} & contract both successors symbolically (\cref{sec:kernel-tier-symbolic}) \\
\texttt{CHAIN\_WINDOW} & \texttt{discharge\_chain\_window} & peel the last chain copy; replace the residual through cut wires (\cref{sec:kernel-chain}) \\
\texttt{TREE\_WINDOW} & \texttt{discharge\_tree\_window} & peel the leaf level of a tree !-box; graft it on the other side (\cref{sec:kernel-tree}) \\
\texttt{HYPOTHESIS\_REWRITE} & \texttt{discharge\_hypothesis\_\allowbreak rewrite} & rewrite the left successor with the hypothesis as a rule, tried first (\cref{sec:kernel-tier-hypothesis}) \\
\texttt{ORACLE\_WINDOW} & \texttt{discharge\_oracle\_window} & compare numerically at $b,\ldots,b+w-1$ (\cref{sec:kernel-tier-oracle}) \\
\bottomrule
\end{tabularx}
\caption{The discharge tiers.}\label{tab:kernel-tiers}
\end{table}

\paragraph{The induction certificate.} A proved verdict with \code!certify_proof! carries an induction \py!Derivation! (\cref{ch:proofs}, certificates). It records the tier's steps but no re-checkable step case, so the proof kernel refuses it (\cref{sec:kernel-check}); the kernel theorem of an induction comes instead from \cref{sec:kernel-indthm}.

\begin{algorithm}[\texttt{\_\allowbreak{}build\_\allowbreak{}derivation}]\label{alg:kernel-derivation}
\begin{enumerate}
\item The \emph{concrete child} at a value $v$ (\py!_concrete_child!): instantiate both hypothesis diagrams at $k=v$ by \py!_at_base! (with the obligation's witness and \py!base_proved!), and rewrite the left into the right by \py!_rewrite_to! (\cref{sec:kernel-tier-uniform}) with \code!rules! (default \py!DEFAULT_RULES!) sorted by rule name, \code!max_steps!, and no guard. When the target is not reached, fail with \code!at k=v: R!; otherwise the child is a \texttt{STEP\_SEQUENCE} derivation from the left instance to the right one with those steps.
\item Build the base child at $v=b$ and the step child at $v=b+1$; the first failure is returned.
\item The claim is \texttt{InductionClaim(\allowbreak{}symbol=n,\allowbreak{} base\_\allowbreak{}value=b,\allowbreak{} held\_\allowbreak{}symbolic=sorted(\allowbreak{}held\_\allowbreak{}symbolic),\allowbreak{} step\_\allowbreak{}discharge=outcome.\allowbreak{}discharge.\allowbreak{}value)}.
\item Return the \texttt{INDUCTION} derivation from \texttt{left\_at\_k} to \texttt{right\_at\_k} with the outcome's steps, the two children, and the claim.
\end{enumerate}
\end{algorithm}

\subsection{Restricted induction}\label{sec:kernel-restricted}

A \term{restriction} is a conjunction of clauses on symbols (\cref{ch:proofs}, restricted decisions): a multiplicity or dimension symbol equal to a constant, at least a constant, of the form $a m+c$ for a fresh $m$, a shift, a side condition, or a finite set.

\begin{algorithm}[\texttt{\_\allowbreak{}restricted\_\allowbreak{}induction}]\label{alg:kernel-restricted}
\begin{enumerate}
\item Parse and validate the restriction against both diagrams (\texttt{validate(\allowbreak{}parse\_\allowbreak{}restriction(\allowbreak{}restrict),\allowbreak{} left,\allowbreak{} right)}); a \py!DecideGrammarError! becomes \py!InductionGrammarError!. An empty restriction returns the unrestricted \py!prove_by_induction!.
\item $n\gets$ \texttt{choose\_\allowbreak{}index(\allowbreak{}left,\allowbreak{} right,\allowbreak{} index=index)}. Start the held values from \code!witness!. For each clause: on $n$, an equality clause raises \py!InductionGrammarError! \texttt{the induction index n is fixed by T}, and any other clause is skipped. On another symbol, an equality clause sets the held value to its constant (a conflicting held value raises \code!x = v lies outside T!); another clause raises the same message when a held value lies outside it (\py!clause.contains!).
\item When \code!family_first!, run the family: \texttt{prove\_\allowbreak{}by\_\allowbreak{}induction(\allowbreak{}left,\allowbreak{} right,\allowbreak{} index=n,\allowbreak{} witness=held,\allowbreak{} **options)}. A proved family returns that result with \py!restriction!, \py!family! (itself), and \py!family_holds=True!. A family counterexample $\{\ldots\text{witness}, n: c\}$ inside the restriction (\py!in_domain!) returns the family result with \py!restriction!, \py!family!, and that point as \py!counterexample_original!.
\item Substitute the restriction into the pair (\py!subst_pair!; a \py!DecideGrammarError! becomes \py!InductionGrammarError!). The inner index is the fresh symbol of the clause on $n$, or $n$ when no clause names it; the inner witness is \texttt{to\_\allowbreak{}restricted(\allowbreak{}restriction,\allowbreak{} held)} (empty when \py!None!).
\item Run \py!prove_by_induction! on the restricted pair over the inner index. A counterexample point is mapped back by \py!to_original!. Return the inner result with \py!restriction!, \py!family!, \py!family_holds=False!, \py!restricted_pair!, and \py!counterexample_original!.
\end{enumerate}
\end{algorithm}

\section{The Rewriting, Contraction, and Oracle Tiers}\label{sec:kernel-tiers}

\subsection{Uniform rewrite}\label{sec:kernel-tier-uniform}

\begin{algorithm}[\texttt{\_\allowbreak{}rewrite\_\allowbreak{}to}]\label{alg:kernel-rewrite-to}
Returns \py!(reached, steps, reason)!.
\begin{enumerate}
\item $W\gets$ a copy of \code!source!; the step list is empty.
\item Repeat \code!max_steps! $+\,1$ times (the last pass only tests the target):
  \begin{enumerate}
  \item When $W\equiv\text{target}$ and \py!_same_bang_boxes(W, target)!, return \texttt{(\allowbreak{}True,\allowbreak{} steps,\allowbreak{} "structurally reached the target")}.
  \item Take the first rule of \code!rules!, in the given order, whose \py!pattern.find_matches(W)! is nonempty, and its first match. When there is none, return \texttt{(\allowbreak{}False,\allowbreak{} steps,\allowbreak{} "no applicable rule")}.
  \item Apply it (\py!apply!; a \py!RewriteError! propagates). When \code!guard_symbols! is given and the free multiplicities of the result differ from it, return \texttt{(\allowbreak{}False,\allowbreak{} steps,\allowbreak{} "a step consumed the multiplicity symbol")} without that step.
  \item $W\gets$ the result; append its step.
  \end{enumerate}
\item Return \texttt{(\allowbreak{}False,\allowbreak{} steps,\allowbreak{} "step budget N exhausted")}.
\end{enumerate}
The test of step~2(a) also runs after the last permitted firing, so a target reached by it is reported as reached.
\end{algorithm}

\py!_same_bang_boxes(a, b)! holds when both diagrams carry the same !-box ids and, per id, equal multiplicities, parents, node scopes (sorted), port scopes (sorted by \py!_port_scope_key!, the triple of node id, direction value, and index), and extensions. The extension of a !-box (\py!_extension!) is the tuple of its links, cyclic flag, index, copy scalar, pair parent, pair index, left ports and right ports (sorted), tree ports, tree parent port, level, position, feed, feed port, alignment, and groups (each sorted).

\begin{signature}
\texttt{discharge\_\allowbreak{}uniform\_\allowbreak{}rewrite(\allowbreak{}obligation,\allowbreak{} *,\allowbreak{} rules=None,\allowbreak{} max\_\allowbreak{}steps=64) -> TierOutcome}

Run \py!_rewrite_to! from \texttt{left\_at\_k} to \texttt{right\_at\_k} with \code!rules! (default \py!DEFAULT_RULES!) sorted by rule name, \code!max_steps!, and the guard \py!free_mult_symbols(left_at_k)!. The outcome is settled exactly when the target is reached; it carries the steps and the reason.
\end{signature}

A uniform rewrite never instantiates $k$, so its steps rewrite $L$ into $R$ for every value of $n$; \cref{sec:kernel-indthm} turns them into a chain of rule inferences on $L$ itself.

\subsection{Induction rewrite and peeling}\label{sec:kernel-tier-peel}

\begin{signature}
\texttt{hypothesis\_\allowbreak{}rule(\allowbreak{}obligation) -> Rule | None} --- a rule rewriting the hypothesis diagram's left side to its right side; it returns \py!None! for every obligation.

\texttt{discharge\_\allowbreak{}induction\_\allowbreak{}rewrite(\allowbreak{}obligation,\allowbreak{} *,\allowbreak{} rules=None,\allowbreak{} max\_\allowbreak{}steps=64) -> TierOutcome}

Run \texttt{discharge\_\allowbreak{}peeled\_\allowbreak{}hypothesis(\allowbreak{}obligation,\allowbreak{} max\_\allowbreak{}steps=max\_\allowbreak{}steps)}. A settled outcome is returned; otherwise the unsettled outcome \texttt{hypothesis not expressible as a rule: H; peeling did not settle it either: P}, with $H$ from \texttt{\_\allowbreak{}hypothesis\_\allowbreak{}inexpressible\_\allowbreak{}reason} and $P$ the peeling reason. The hypothesis as a rule is the separate tier \texttt{HYPOTHESIS\_REWRITE} (\cref{sec:kernel-tier-hypothesis}).

\paragraph{The hypothesis rule.}\label{sec:kernel-tier-hypothesis} \py!hypothesis_rule(obligation)! is \py!None! when $L(k)$ has no node or no box whose multiplicity carries the step symbol $k$, when the interfaces or \py!boundary_layout!s of $L(k)$ and $R(k)$ differ, when $L(k)$'s scalar is zero, or when the ratio $\rho = R(k).\text{scalar}/L(k).\text{scalar}$ has a free symbol. Otherwise it is the rule \code!induction_hypothesis! (not automatic) with pattern \py!HypothesisPattern(L(k), k)! of \texttt{rewrite/match\_hypothesis.py} and scalar $\rho$. Its \py!find_matches(H)!:
\begin{enumerate}
\item For each $c\ge 0$ with $M' - m_b(k) = c$ for some box $b$ of $L(k)$ carrying $k$ and some host box multiplicity $M'$, ascending: peel $c$ copies off every owner of $k$ (\py!peel_owners! $c$ times); $k$ stays bound.
\item Embed $L(k)$ in the peeled host: an injective node map keeping generator, leg dimensions and phase, every wire of $L(k)$ port for port; a boundary port of $L(k)$ lands on a host port wired off the image or on the host boundary, a boxed slot (\py!slot_owners!) on the host boundary.
\item Each box of $L(k)$ maps onto one host box over the mapped footprint, every field (multiplicity, parent, links, pair, ports, groups) mapped; every host box meeting the image is such an image.
\end{enumerate}
\py!apply_hypothesis(match, L, R)! removes the image boxes and nodes, grafts $R(k)$ with its boxes, joins $R$'s $i$-th boundary port to the far end of the $i$-th cut wire or puts it in the host boundary slot $L$'s port held, and multiplies the scalar by $\rho$.

\texttt{discharge\_\allowbreak{}hypothesis\_\allowbreak{}rewrite(\allowbreak{}obligation,\allowbreak{} *,\allowbreak{} rules,\allowbreak{} max\_\allowbreak{}steps)} rewrites the left successor, at each step the first hypothesis match if there is one, else the first match of \code!rules! (default \py!SPIDER_FUSION!) by name, until it is the right successor or that successor with its owners of $k$ peeled, field for field or isomorphic with equal scalars. A derivation without a hypothesis step is refused with \texttt{derivation never used the hypothesis}.
\end{signature}

\texttt{\_\allowbreak{}hypothesis\_\allowbreak{}inexpressible\_\allowbreak{}reason(\allowbreak{}obligation)} returns \texttt{the hypothesis and successor left sides carry different bang-box ids} when the !-box ids of \texttt{left\_at\_k} and \texttt{left\_at\_successor} differ; else, for the first id (in sorted order) whose multiplicities differ, \texttt{bang box b carries multiplicity M in the hypothesis and M' in the successor; every Pattern here matches structure only and none unifies Mult arithmetic,\allowbreak{} and a match ignoring the multiplicity would rewrite the successor straight onto its own target}; else \texttt{the hypothesis left side already equals the successor left side field for field}.

\begin{figure}[H]
\centering
\begin{tikzpicture}[zxd]
  \node[zs] (a) at (0,0) {$\alpha$};
  \draw[wire] (a) -- (0,0.9); \draw[wire] (a) -- (0,-0.9);
  \node[bang,fit=(a),minimum width=1cm,minimum height=1cm] (B) {};
  \node[banglabel] at (B.north west) {$k{+}1$};
  \node at (1.4,0) {$=$};
  \node[zs] (b) at (2.6,0) {$\alpha$};
  \draw[wire] (b) -- (2.6,0.9); \draw[wire] (b) -- (2.6,-0.9);
  \node[bang,fit=(b),minimum width=1cm,minimum height=1cm] (C) {};
  \node[banglabel] at (C.north west) {$k$};
  \node[zs] (c) at (3.8,0) {$\alpha$};
  \draw[wire] (c) -- (3.8,0.9); \draw[wire] (c) -- (3.8,-0.9);
\end{tikzpicture}
\caption{Peeling one copy off a node-scoped !-box: the residual is the !-box at $k$, the copy is laid out after its block.}\label{fig:kernel-peel}
\end{figure}

\begin{algorithm}[\footnotesize\texttt{discharge\_peeled\_hypothesis}]\label{alg:kernel-peeled}
All outcomes have discharge \texttt{INDUCTION\_REWRITE}. Let $k$ be the step symbol.
\begin{enumerate}
\item Peel each successor at $k$ by \py!_peel_step_index!: \py!peel_owners(D, k)!; \py!None! on \py!BangBoxError!, on a non-separable peel, or on an empty copy. When either side gives \py!None!, return unsettled with \texttt{the step index is not carried by a single separable bang box on each side,\allowbreak{} so no copy can be peeled off to expose the hypothesis}.
\item Check each residual by \texttt{\_\allowbreak{}residual\_\allowbreak{}is(\allowbreak{}peeled,\allowbreak{} copy,\allowbreak{} hypothesis,\allowbreak{} k)} (below); on failure return \texttt{the left successor's peeled residual is not the left hypothesis diagram} (or the right).
\item When the copies occupy different boundary positions (\py!_boundary_positions!: the indices of the boundary inputs and of the boundary outputs on the copy's nodes), return \texttt{the peeled copies occupy different boundary slots on the two sides,\allowbreak{} so the two successors differ by a permutation the hypothesis does not settle}.
\item When the peeled successors' scalars differ, return \texttt{the two successors carry different scalars,\allowbreak{} S and S'}.
\item Extract each copy with its !-boxes (\texttt{\_\allowbreak{}extract(\allowbreak{}peeled,\allowbreak{} copy,\allowbreak{} with\_\allowbreak{}boxes=True)}, \cref{sec:kernel-extract}). When their interfaces differ, return \texttt{the peeled copies do not share one boundary interface: I against I'}.
\item When the copies are identical, return settled (\texttt{.\allowbreak{}.\allowbreak{}.\allowbreak{}the peeled copies are structurally identical}) with both peeled successors and copies.
\item Contract both copies symbolically (\texttt{contract\_\allowbreak{}symbolic\_\allowbreak{}pair(\allowbreak{}.\allowbreak{}.\allowbreak{}.\allowbreak{},\allowbreak{} max\_\allowbreak{}steps)}). On \texttt{SymbolicContractionUnsupportedError}, \texttt{SymbolicContractionDomainError}, or \py!ScalarBudgetError!: when \texttt{lemma\_\allowbreak{}decision(\allowbreak{}copy\_\allowbreak{}L,\allowbreak{} copy\_\allowbreak{}R)} (\cref{sec:kernel-lemma-decision}) is not \py!None!, return settled with reason \texttt{.\allowbreak{}.\allowbreak{}.\allowbreak{}the peeled copies are equal by a nested decision (\allowbreak{}lemma)}; else unsettled with the error text.
\item Compare the tensors (\py!compare_symbolic!). On a mismatch with a lemma decision, return the lemma outcome. On a match, return settled with \texttt{.\allowbreak{}.\allowbreak{}.\allowbreak{}the peeled copies agree under symbolic contraction with d formal} and the comparison; otherwise unsettled \texttt{the peeled copies are not equal: R}.
\end{enumerate}
Every settled reason begins \texttt{peeled one copy off each successor; the residuals are the hypothesis diagrams and}.
\end{algorithm}

\paragraph{Residual test.} \texttt{\_\allowbreak{}residual\_\allowbreak{}is(\allowbreak{}peeled,\allowbreak{} copy,\allowbreak{} expected,\allowbreak{} index="")} holds when (a) the nodes of \code!peeled! off \code!copy! are exactly the nodes of \code!expected! (by id); (b) \py!_extract! of those nodes is identical to \py!_extract! of all of \code!expected!; (c) no !-box of \code!peeled! has a footprint that meets \code!copy! without lying inside it; (d) the !-boxes whose footprint avoids \code!copy! have exactly the ids of \code!expected!'s, and per id equal multiplicities, parents, extensions, node scopes, and port scopes; and (e) none of them is a chain !-box whose multiplicity reads \code!index! (every chain !-box when \code!index! is empty).

\paragraph{Soundness of the tier.} With $L(k+1)=L(k)\otimes c_L$ and $R(k+1)=R(k)\otimes c_R$ in the same boundary positions and with equal scalars, the hypothesis $L(k)=R(k)$ and $c_L=c_R$ give $L(k+1)=R(k+1)$.

\subsection{Symbolic contraction}\label{sec:kernel-tier-symbolic}

\begin{signature}
\texttt{discharge\_\allowbreak{}symbolic\_\allowbreak{}contraction(\allowbreak{}obligation,\allowbreak{} *,\allowbreak{} max\_\allowbreak{}steps=64) -> TierOutcome}

Contract both successors with every symbol formal (\py!contract_symbolic_pair!, both doubled when either has a discard; \cref{ch:semantics}). A \texttt{SymbolicContractionUnsupportedError}, \texttt{SymbolicContractionDomainError}, or \py!ScalarBudgetError! gives an unsettled outcome with its text; every other error propagates. A match of \py!compare_symbolic! settles the step with \texttt{successor case closed by symbolic contraction with every symbol formal}; a mismatch gives the comparison's reason.
\end{signature}

\subsection{The oracle window}\label{sec:kernel-tier-oracle}

\begin{signature}
\texttt{discharge\_\allowbreak{}oracle\_\allowbreak{}window(\allowbreak{}obligation,\allowbreak{} *,\allowbreak{} width=4,\allowbreak{} mode=EXACT,\allowbreak{} tolerance,\allowbreak{} max\_\allowbreak{}elements) -> TierOutcome}

Raises \py!InductionDomainError! \texttt{window width must be a positive int,\allowbreak{} got W} when \code!width! is not an \py!int! (a \py!bool! is rejected) or is below $1$, and \py!InductionGrammarError! \texttt{a numeric tier leaves symbol(\allowbreak{}s) uninstantiated: [.\allowbreak{}.\allowbreak{}.\allowbreak{}]; name them in the witness to hold them at a value} when a held symbol has no witness value. For $v=b,\ldots,b+w-1$, compare the hypothesis diagrams through the oracle at $\{\ldots\text{witness}, k: v\}$. The first mismatch returns unsettled \code!mismatch at n=v! with \py!counterexample=v!. Otherwise the outcome is settled with \texttt{equality confirmed at n = b.\allowbreak{}.\allowbreak{}b+w-1; a finite schema check,\allowbreak{} not a proof}.
\end{signature}

\section{Regions, Cuts, and Replacement}\label{sec:kernel-regions}

The windowed tiers and the kernel rules \textsc{Congr}, \textsc{Induction}, and \textsc{LocalContraction} read a set of nodes of a diagram as a diagram of its own and replace it by another. This section specifies those operations.

\subsection{Extraction}\label{sec:kernel-extract}

\begin{algorithm}[\texttt{\_\allowbreak{}extract}]\label{alg:kernel-extract}
Output: the sub-diagram on \code!node_ids!, scalar one.
\begin{enumerate}
\item Add each node of \code!node_ids!, in id order, to an empty diagram with its generator, port dimensions, and phase; record the id map.
\item Add every wire with both ends in \code!node_ids!, in wire sort order.
\item The boundary inputs (outputs) are those of \code!diagram! on \code!node_ids!, in the enclosing order, on the mapped nodes.
\item With \code!with_boxes!: for each !-box in id order with $P = \text{node scope}\cap\text{node\_ids}$, keep it when its footprint lies inside \code!node_ids!, or when $P\neq\emptyset$ and it is a plain straddle or a grouped straddle (below). The image of a kept !-box is the !-box itself (footprint inside), the !-box cut to $P$ (a plain straddle, \py!with_node_scope(P)!), or \py!_group_part(box, P)! (a grouped straddle: its links whose first port lies in $P$, its groups inside $P$, and no groups when only one is left). Any other !-box whose footprint meets \code!node_ids! raises \py!InductionDomainError! \texttt{bang box b lies partly inside the region}; a kept pair !-box whose parent is not kept raises \texttt{a pair box lies inside a region its parent leaves}. Add the images in dependency order (a !-box after its kept parent, its kept pair parent, and its kept alignment target) with \py!add_box_image!; an alignment to a !-box not kept is dropped.
\item Copy the parameter environment.
\end{enumerate}
\end{algorithm}

\begin{definition}[Straddling !-boxes]\label{def:kernel-straddle}
A !-box is a \term{plain straddle} (\py!_plain_straddle!) when it is a top-level node-scoped !-box with no links, copy scalar, tree data, feed, pairing, alignment, or groups, and no !-box is aligned to it. A !-box partly inside a region $X$ is a \term{grouped straddle} (\py!_grouped_straddle!) when it is a top-level node-scoped !-box with groups and with no index, copy scalar, alignment, tree data, feed, pairing, or cyclic chain; no other !-box names it as parent, alignment target, pair parent, or feed; each of its blocks (\py!_blocks!: its groups, then its ungrouped scope nodes when there are any) lies wholly in or wholly off $X$; and, for each link, its two ports and the ports they are wired to lie on one side of $X$.
\end{definition}

\subsection{Connected regions and the hypothesis region}

\py!connected_regions(D)! partitions the nodes of $D$ into the sets closed under wires and !-box footprints: a union--find structure~\cite{tarjan1975unionfind} with path halving joins the two ends of each wire and all nodes of each footprint, always keeping the smaller root id; the classes are returned ordered by root, that is, by smallest node id.

\begin{signature}
\texttt{find\_\allowbreak{}hypothesis\_\allowbreak{}region(\allowbreak{}diagram,\allowbreak{} hypothesis,\allowbreak{} *,\allowbreak{} max\_\allowbreak{}regions=12) -> frozenset[NodeId] | None}

The smallest union of connected regions whose extraction with !-boxes is isomorphic to \py!_bare(hypothesis)! with symmetric legs. Unions are tried by increasing size and, within a size, in the order of \py!itertools.combinations! over the region list; a union is skipped when its node count or its sorted list of generator names differs from the hypothesis's, or when extraction raises \py!InductionDomainError!. Past \code!max_regions! regions only single regions are tried.
\end{signature}

\subsection{Boundary layout}\label{sec:kernel-layout}

\begin{definition}[Slot owners and layout]\label{def:kernel-layout}
The \term{owners} of a boundary entry $r$ (\py!_slot_owners!) are the !-boxes that replicate it: the node-scoped !-boxes over its node (a chain !-box not, when $r$ is one of its link ports) and the port-scoped !-boxes holding $r$, outermost first (sorted by nesting depth, then id). The \term{end dimension} of $r$ (\py!_end_dim!) is its port dimension, except on an open link end of an indexed chain: the copy index read at the copy holding it, the last copy ($\iota := \mu-1$) for a link's $p$ port and the first ($\iota:=0$) for its $q$ port. The \term{boundary layout} (\py!boundary_layout!) is, per boundary list, the tuple of rows $(\text{owner multiplicities}, \text{opens}, \text{end dimension})$, where \emph{opens} tells per owner whether the entry starts a new run of that !-box (the owner was not an owner of the previous entry, or the entry lies in another group of the !-box than the previous entry; an entry in no group counts as group $-1$). The layout is \py!None! when a !-box's entries in one list are not contiguous, that is, when an owner reappears after a run of it closed.
\end{definition}

\subsection{Replacing a closed region}

\begin{algorithm}[\texttt{replace\_\allowbreak{}region}]\label{alg:kernel-replace-region}
\begin{enumerate}
\item When \code!replacement! has a different number of boundary inputs (outputs) than \code!region! holds, raise \py!InductionDomainError! \texttt{the replacement's interface does not fit the region's slots}.
\item When the layout of \code!hypothesis! is \py!None! or differs from that of \code!replacement!, raise \texttt{the replacement's boundary layout differs from the region's}; when the layout of \code!diagram! is \py!None!, raise \texttt{a box's boundary slots are not contiguous}.
\item Every !-box whose footprint meets \code!region! must lie inside it, else \texttt{bang box b straddles the replaced region}; remove those !-boxes (in decreasing id order) and the region's nodes.
\item Graft \code!replacement! (\py!_graft!: its nodes in id order, its wires, its !-boxes in dependency order).
\item Refill each boundary list in order: each entry that was on the region takes the next entry of the replacement's list.
\item Multiply the scalar by $\text{replacement.scalar}/\text{hypothesis.scalar}$.
\end{enumerate}
\end{algorithm}

\subsection{Cut regions}\label{sec:kernel-cut}

A region that meets the rest of the diagram through wires is read as an open diagram by cutting those wires.

\begin{definition}[Cut region]\label{def:kernel-cut}
\py!CutRegion! (frozen, slotted, compared by value) has fields \py!region: frozenset[NodeId]!, \py!inputs: tuple[PortRef, ...]!, \py!outputs: tuple[PortRef, ...]!, and \texttt{cuts: Mapping[PortRef,\allowbreak{} PortRef]}. The \term{cut wires} of a region $X$ (\py!_cut_wires!) are the wires with exactly one end in $X$, as a map from the end in $X$ to the far end. Each entry of \code!inputs! and \code!outputs! is either a boundary entry of the diagram on $X$ or the region end of a cut wire. \py!_cut_extract(D, cut)! is \texttt{\_\allowbreak{}extract(\allowbreak{}D,\allowbreak{} region,\allowbreak{} with\_\allowbreak{}boxes=True)} with its boundary lists replaced by the images of \code!cut.inputs! and \code!cut.outputs! (the region's nodes mapped in sorted order onto the extracted nodes in sorted order).
\end{definition}

\begin{algorithm}[\texttt{find\_\allowbreak{}cut\_\allowbreak{}region}]\label{alg:kernel-find-cut}
Output: a \py!CutRegion! under which the region is the hypothesis up to scalar, or \py!None!.
\begin{enumerate}
\item Compute the cut wires; per direction (inputs, then outputs), let \emph{own} be the region ends of that port direction in port sort order, \emph{kept} the indices of the boundary entries on the region, and $w$ the hypothesis's number of entries in that list. When $|\text{kept}|+|\text{own}|\neq w$, return \py!None!.
\item The candidate \emph{spots} are the indices of the boundary entries off the region (the holes) followed by the positions past the end of the list. For each permutation of $|\text{own}|$ spots (\py!itertools.permutations! order), place the cut ends at those spots and read off the list over the sorted union of kept indices and chosen spots. Above $2000$ (\py!_MAX_CUT_PLACEMENTS!) options in one list, return \py!None!.
\item For each input option and, inside it, each output option, build the cut region; when \py!_cut_extract! raises \py!InductionDomainError!, return \py!None!; when the bare extraction is isomorphic (legs not symmetric) to the bare hypothesis, return it.
\item Return \py!None!.
\end{enumerate}
\end{algorithm}

\begin{algorithm}[\texttt{replace\_\allowbreak{}cut\_\allowbreak{}region}]\label{alg:kernel-replace-cut}
\begin{enumerate}
\item When the replacement's list lengths differ from the cut's, raise \py!InductionDomainError! \texttt{the replacement's interface does not fit the cut region}; the layout checks are those of \cref{alg:kernel-replace-region}, the contiguity of \code!diagram!'s layout only when \code!contiguous!.
\item Every !-box whose footprint meets the region and does not lie inside it must be a mergeable plain straddle (\py!_straddle_merges!: a plain straddle each of whose children, pair !-boxes, and aligned !-boxes lies inside the region, and whose index no node of its scope off the region reads) or a grouped straddle; else raise \texttt{bang box b straddles the replaced region}. Record these \emph{outside parts}.
\item Record the \emph{alignment partners} (\py!_outside_partners!): per alignment group with members both meeting the region and off it, its members off the region and its root when the root is off the region. Clear the alignment of those members.
\item Remove every !-box meeting the region (decreasing id), graft the replacement, and map slot $i$ of each cut list to the image of slot $i$ of the replacement's list. Rewrite the diagram's boundary lists through that map, remove the region's nodes, and wire each cut wire's far end to the image of its region end, in port sort order.
\item Merge each outside part into the grafted !-boxes: a grouped straddle by \py!_merge_grouped!, a plain one by \py!_merge_straddle! (below).
\item Realign each alignment group (\py!_realign!, below).
\item Set the boundary lists and multiply the scalar by $\text{replacement.scalar}/\text{hypothesis.scalar}$.
\end{enumerate}
\end{algorithm}

\paragraph{Merging a plain straddle.} \py!_merge_straddle! splits the straddle's nodes off the region into wire-connected parts (\py!_parts!, smallest id first). Each part goes into the one grafted top-level !-box that every cut wire from the part reaches, which must hold the grafted end alone among the grafted !-boxes, carry the straddle's multiplicity, and be no tree, fed, or pair !-box; its node scope grows by the part. A part reaching no such single !-box raises \texttt{bang box b straddles the replaced region}.

\paragraph{Merging a grouped straddle.} \py!_merge_grouped! requires every cut wire from the straddle's part $P$ off the region to start at a node of the straddle and to end at a node held by exactly one !-box, a grafted one; all these must be one target. The target must carry the straddle's multiplicity and sheets, be top level, and carry no index, copy scalar, alignment, tree data, feed, pairing, or cyclic chain, and no other !-box may name it as parent, alignment target, pair parent, or feed. The target's node scope grows by $P$, its links by the straddle's links on $P$, and its groups become its blocks followed by the straddle's blocks inside $P$. Any failure raises \texttt{bang box b straddles the replaced region}.

\paragraph{Realignment.} \py!_realign(working, members, root, cut, slot_of, grafted, replaced=None)! collects the grafted !-boxes that hold the grafted ends of the cut wires from the members' scopes; each such end must lie in exactly one grafted !-box (\texttt{aligned boxes [.\allowbreak{}.\allowbreak{}.\allowbreak{}] meet no single box}), and at most one target is allowed (\texttt{aligned boxes [.\allowbreak{}.\allowbreak{}.\allowbreak{}] meet several boxes}). With no target and a shape $(\mu,\mathit{index},\mathit{links})$ of the replaced alignment group (\py!_replaced_shape!: the one member inside the region of the replaced group's first member, read as its multiplicity, index, and whether it has links), the target is the one grafted !-box that is top level, node scoped, no pair, tree, fed, grouped or aligned !-box, without copy scalar, and of that shape, when exactly one exists. The new group is the members, plus the alignment group of the target when there is one; its root is the old root when it lay off the region, else the root of the target's group, else the least member. Every member must carry the root's multiplicity, sit inside every !-box that holds the root, be node scoped, and be no pair, tree, or fed !-box (\texttt{bang box b cannot align with r}). Each member but the root is then aligned to the root.

\subsection{The part off a region}\label{sec:kernel-copy-side}

\begin{signature}
\texttt{\_\allowbreak{}copy\_\allowbreak{}side(\allowbreak{}exposed,\allowbreak{} cut,\allowbreak{} scalar) -> (\allowbreak{}Diagram,\allowbreak{} held)}

The part of \code!exposed! off \code!cut.region! as an open diagram: \py!_extract! with !-boxes of the remaining nodes; its boundary inputs (outputs) are its own boundary entries in order, followed by the far ends of the cut wires whose far port is an input (output), ordered by the place of their region end in the cut's lists (inputs before outputs, then by position); its scalar is \code!scalar!. \code!held! is the tuple (positions of the region's boundary inputs, positions of its boundary outputs, places of the region ends in cut order, directions of the far ends).
\end{signature}

\py!rest_cut(exposed, cut)! is the part off the region as a \py!CutRegion! of \code!exposed!, its lists ordered as \py!_copy_side! orders them and its cuts the cut wires of the remaining nodes. \py!_open_region(D, X)! is $X$ as an open diagram: its boundary entries, then its ends of the cut wires in port order, with $D$'s scalar (computed as \py!_copy_side! of the cut region of the complement).

\section{Exposing the Hypothesis}\label{sec:kernel-expose}

After a peel, the residual of a chain or a spider-heavy successor is often the hypothesis diagram only after a local rearrangement: the peeled copy shares a spider with the residual, or legs sit on the wrong spider. The functions of this section produce those rearrangements; each is an equation of the ZXW calculus (spider fusion and unfusion, and identity removal), and each has a kernel rule that re-derives it (\cref{sec:kernel-rules}).

\subsection{Spider splits}

\begin{figure}[H]
\centering
\begin{tikzpicture}[zxd]
  \node[zs] (v) at (0,0) {$\alpha$};
  \draw[wire] (v) -- (0,1.1) node[above] {\tiny slot};
  \draw[wire] (v) -- (-0.8,-1) node[below] {\tiny other};
  \draw[wire] (v) -- (0.8,-1) node[below] {\tiny into copy};
  \node at (1.8,0) {$=$};
  \node[zs] (h) at (3.2,-0.4) {$\alpha$};
  \node[zdot] (l) at (3.9,0.5) {};
  \draw[wire] (l) -- (3.9,1.3) node[above] {\tiny slot};
  \draw[wire] (h) -- (l);
  \draw[wire] (h) -- (2.6,-1.2) node[below] {\tiny other};
  \draw[wire] (l) .. controls (4.5,0) .. (4.5,-1.2) node[below] {\tiny into copy};
  \node[font=\tiny] at (2.7,-0.1) {hub};
  \node[font=\tiny] at (4.25,0.75) {leaf};
\end{tikzpicture}
\caption{A spider split: a phaseless leaf takes the boundary entry and the legs into the copy; the hub keeps the phase and the other legs and is wired to the leaf in the entry's place.}\label{fig:kernel-split}
\end{figure}

\texttt{\_\allowbreak{}split\_\allowbreak{}spider(\allowbreak{}D,\allowbreak{} v,\allowbreak{} slot,\allowbreak{} inward,\allowbreak{} copy)} splits the Z or X spider $v$ at its boundary leg \code!slot!: a hub node of $v$'s generator with $v$'s phase and its legs other than \code!inward!, and a phaseless leaf of the same generator whose legs are, in the slot's direction, the slot leg first, in the other direction a connector first, and then the \code!inward! legs appended in their directions. Every wire of $v$ is moved onto the hub or the leaf (a self-loop onto the images of both ends), the leaf's connector is wired to the hub's port at the slot's index, the boundary entry moves to the leaf, $v$ is removed, and the returned copy is $(\text{copy}\setminus\{v\})\cup\{\text{leaf}\}$.

\begin{signature}
\py!_fan_split(D, copy)! yields, in order:
\begin{enumerate}
\item for each Z or X spider $v$ of \code!copy! (by id) that has a leg into \code!copy! and a leg to a boxed node off \code!copy!, and for each boundary leg of $v$: \texttt{\_\allowbreak{}split\_\allowbreak{}spider(\allowbreak{}D,\allowbreak{} v,\allowbreak{} slot,\allowbreak{} inward,\allowbreak{} copy)};
\item the splits of \py!_gather_split!: for each unboxed Z or X spider off \code!copy! with a leg into \code!copy! and a leg to a boxed node off it, and for each of its boundary legs, the same split (the leaf joins the copy);
\item the splits of \py!_boxed_split! (below).
\end{enumerate}
\py!_gather_splits(D, copy)! yields only the second kind; \py!boxed_splits(D, copy)! only the third.
\end{signature}

\paragraph{Boxed splits.} \texttt{\_\allowbreak{}boxed\_\allowbreak{}split(\allowbreak{}D,\allowbreak{} copy,\allowbreak{} partner,\allowbreak{} slots)} takes each Z or X spider $v$ off \code!copy! held by at least one node-scoped !-box (its \emph{owners}). For each leg of $v$ wired into \code!copy!, the leaf's \emph{home} is \py!_leaf_home(owners, far)!: the owners themselves when the far node lies in all of them; the one top-level !-box holding the far node when the owners are one top-level !-box aligned with it and the far node lies in no other !-box; else none. A leg whose home is not the owners and whose far node has $v$'s generator is skipped. Homes are tried with the owners first; per home and per boundary leg of $v$, \py!_split_boxed_spider! builds the split, which fails (\py!None!) when an inward leg precedes a kept leg of its direction. In a boxed split the hub takes over every reference to $v$ (wires, boundary entries, and !-box fields, by \py!_move_node!) at the same direction and index, the leaf carries the slot leg, a connector, and the inward legs, the leaf is wired to the hub's port in the slot's place, the boundary entry moves to the leaf, and the leaf joins the node scope of each home !-box; in each such !-box a group holding the hub gets the leaf, and loses the hub when the hub holds no boundary entry. The yielded copy is $\text{copy}\cup\{\text{leaf}\}$.

\subsection{Leg moves}

\begin{figure}[H]
\centering
\begin{tikzpicture}[zxd]
  \node[zs] (w) at (0,0.5) {};
  \node[zs] (v) at (0,-0.5) {};
  \draw[wire] (w) -- (0,1.4) node[above] {\tiny slot};
  \draw[wire] (w) -- (v);
  \draw[wire] (v) -- (-0.6,-1.4);
  \draw[wire] (v) -- (0.7,-1.4) node[below] {\tiny into copy};
  \node at (1.6,0) {$=$};
  \node[zs] (w2) at (3,0.5) {};
  \node[zs] (v2) at (3,-0.5) {};
  \draw[wire] (w2) -- (3,1.4) node[above] {\tiny slot};
  \draw[wire] (w2) -- (v2);
  \draw[wire] (v2) -- (2.4,-1.4);
  \draw[wire] (w2) .. controls (3.8,0.2) .. (3.8,-1.4) node[below] {\tiny into copy};
\end{tikzpicture}
\caption{A leg move: legs into the copy move from $v$ (holding no boundary entry) to the same-type spider $w$ it is wired to, which holds one.}\label{fig:kernel-legmove}
\end{figure}

\begin{signature}
\texttt{leg\_\allowbreak{}moves(\allowbreak{}D,\allowbreak{} copy) -> Iterator[Diagram]}

Let \emph{fixed} be the nodes a !-box holds or whose ports a !-box names (\py!_box_referenced!: links, pair ports, tree ports, tree parent, feed port, port scope). For each Z or X spider $v$ (by id) off \code!copy!, not fixed, holding no boundary entry, with legs into \code!copy!, and for each leg of $v$ wired to a node $w\neq v$ off \code!copy!, not fixed, holding a boundary entry, of $v$'s generator: yield \py!_move_legs(D, v, inward, w)!, in which $v$ and $w$ are rebuilt as new nodes, $v$'s inward legs removed and appended after $w$'s legs of their direction, each wired to its old far end, and every other wire and boundary entry re-mapped.
\end{signature}

\subsection{Identity slides}

\begin{signature}
\texttt{identity\_\allowbreak{}slides(\allowbreak{}D,\allowbreak{} copy) -> Iterator[Diagram]}

For each phaseless Z spider $v\notin\text{copy}$ with one input and one output of equal dimension, exactly one leg a boundary entry, no leg named by a !-box (\py!_box_ports!), and the other leg wired to a port $p$ that no !-box names, that is not a boundary entry, and whose node lies in exactly the node-scoped !-boxes holding $v$ (\py!_holders!); and for each wire (in sort order) from a port $a$ off \code!copy! to a port $b$ of a node of \code!copy!, $a$'s node $\neq v$, $b$ not named by a !-box, $a$ named by a !-box only as a link port, $b$'s node held by the same !-boxes as $v$, and $b$'s port dimension equal to $v$'s: yield \py!_slide!: $v$'s wires and the wire $a$--$b$ removed, $p$ put in $v$'s boundary entry, $v$ wired between $a$ and $b$ (from the output end to $v$'s input, from $v$'s output to the input end), and, in each grouped !-box holding $v$ where $p$'s group is not $v$'s, $p$ moved into a new group of its own.
\end{signature}

\subsection{Copy candidates and the exposing search}

\begin{signature}
\texttt{copy\_\allowbreak{}candidates(\allowbreak{}peeled,\allowbreak{} old,\allowbreak{} copy) -> list[frozenset[NodeId]]}

Let $N$ be the nodes of \code!peeled! not in \code!old!, \emph{bare} those of $N$ in no node-scoped !-box, and \emph{extra} the boxed nodes of $N$, in no pair !-box, wired to a bare one. The list is: \code!copy!; then, when \emph{extra} is nonempty and $\text{bare}\cup\text{extra}\neq\text{copy}$, $\text{bare}\cup\text{extra}$; then, when \emph{extra} is nonempty and some bare spider is wired to a node of \code!old!, $(\text{bare}\setminus\text{those spiders})\cup\text{extra}$; then, when some spider of \code!copy! is wired to a node of $\text{old}\setminus\text{copy}$ (a \emph{hub}) and $\text{copy}\setminus\text{hubs}$ is not yet listed, $\text{copy}\setminus\text{hubs}$.
\end{signature}

\begin{signature}
\py!ExposeStep! (frozen, slotted): \py!kind! (\code!"split"! or \code!"move"!), \py!before!, \py!after!, \py!copy_before!, \py!copy_after!.
\end{signature}

\begin{algorithm}[\texttt{expose\_\allowbreak{}cut}]\label{alg:kernel-expose-cut}
Output: \py!(diagram, cut, steps)! or \py!None!.
\begin{enumerate}
\item For each copy of \code!copies!: when the complement of the copy is the hypothesis through cut wires (\py!find_cut_region!), return \py!(peeled, cut, ())!; else, for each split of \py!_fan_split(peeled, copy)!, when the complement of the split copy is the hypothesis through cut wires, return \texttt{(\allowbreak{}split,\allowbreak{} cut,\allowbreak{} (\allowbreak{}ExposeStep(\allowbreak{}"split",\allowbreak{} .\allowbreak{}.\allowbreak{}.\allowbreak{}),))}.
\item For each copy of \code!copies[1:]!: run \py!_greedy_steps(peeled, copy)!; try the prefixes of the steps in the order all steps, then lengths $1, 2, \ldots$, one fewer than all; for each, when the complement of the last step's copy in its result is the hypothesis through cut wires, return that result, the cut, and the prefix.
\item Return \py!None!.
\end{enumerate}
\end{algorithm}

\py!_greedy_steps(peeled, copy)! runs three stages in order: boxed splits (\py!boxed_splits!), leg moves (\py!leg_moves!), and unboxed splits at an entry (\py!_gather_splits!). Each stage repeatedly takes the first rearrangement its generator yields for the current diagram and copy and records it as an \py!ExposeStep! (kind \code!"move"! for a leg move, \code!"split"! otherwise), until the generator is empty or $16$ steps (\py!_MAX_GREEDY_STEPS!) are recorded in all.

\section{The Chain Window}\label{sec:kernel-chain}

A \term{chain !-box} is a node-scoped !-box with links (\cref{ch:diagrams}): expanding it wires the $p$ port of each copy to the $q$ port of the next. Peeling a chain leaves a residual that meets the peeled copy through link wires, so the residual is the hypothesis diagram only through cut wires. The chain window reads it so, replaces it by the other side's hypothesis diagram, and compares the result with the other successor.

\paragraph{Chain helpers.} \py!_chain_owner(D, k)! is the single !-box whose multiplicity reads $k$ when that !-box has links, else \py!None!; \py!_chain_owned(D, k)! holds when an owner of $k$ (\py!index_owners!) has links. \py!first_peelers(D, k)! is the set of owners of $k$ that peel their first copy. \py!peel_peelable(D, k)! is the joint peel, else the partial joint peel, else \py!None!.

\begin{definition}[Chain core]\label{def:kernel-chain-core}
For a chain !-box $c$ of $D$ (\py!box_core(D, c)!), the \emph{head} is the set of nodes off $c$'s scope reached from the far ends of the $q$ ports of its links by paths that avoid the scope, and the \emph{tail} likewise from the $p$ ports. The \term{chain core} is the scope together with the head, and is \py!None! when head and tail meet. \py!_chain_core(D, k)! is the core of \py!_chain_owner(D, k)!. \py!chain_cores(H, a, b, k)! returns \texttt{(\allowbreak{}\_\allowbreak{}open\_\allowbreak{}region(\allowbreak{}H,\allowbreak{} core\_\allowbreak{}H),\allowbreak{} core\_\allowbreak{}a,\allowbreak{} core\_\allowbreak{}b)}, or the text \texttt{no chain box with separate head and tail on both sides} when a core is missing, or \texttt{the chain boxes have multiplicities M1,\allowbreak{} M2,\allowbreak{} M3} when the three chain !-boxes differ in multiplicity.
\end{definition}

\paragraph{Closing a window.} \texttt{\_\allowbreak{}close\_\allowbreak{}window(\allowbreak{}replaced,\allowbreak{} target,\allowbreak{} max\_\allowbreak{}steps)} returns \texttt{(\allowbreak{}True,\allowbreak{} "the result is the other successor")} when the two are identical with the same !-boxes; \texttt{(\allowbreak{}True,\allowbreak{} "the result is isomorphic to the other successor")} when they are isomorphic, with symmetric legs unless either carries a chain; else contracts both symbolically and returns \texttt{the result and the other successor agree under symbolic contraction} on a match, or \code!symbolic contraction: R! on a mismatch or on an unsupported shape, a domain error, or an exhausted scalar budget.

\begin{algorithm}[\texttt{\_\allowbreak{}compare\_\allowbreak{}chain\_\allowbreak{}copies}]\label{alg:kernel-chain-copies}
\begin{enumerate}
\item For each side, \texttt{\_\allowbreak{}exposed\_\allowbreak{}chain(\allowbreak{}successor,\allowbreak{} hypothesis,\allowbreak{} k)}: the successor must be chain owned; peel it (\py!_peel_step_index!); the complement of the copy must be the hypothesis through cut wires and the peeled diagram's layout must exist. Failing either side: \texttt{the two successors are not both peelable chains over their hypotheses}.
\item The two hypothesis diagrams must have one layout (\texttt{the two hypothesis diagrams lay out their boundaries differently}).
\item Take each copy side, \texttt{\_\allowbreak{}copy\_\allowbreak{}side(\allowbreak{}exposed,\allowbreak{} cut,\allowbreak{} exposed.\allowbreak{}scalar / hypothesis.\allowbreak{}scalar)}. The \code!held! values must agree (\texttt{the residuals sit in different boundary places on the two sides}) and so must the interfaces (\texttt{the peeled chain copies have different interfaces}).
\item Close the window between the two copy sides; when that fails and \py!lemma_decision! of the two is not \py!None!, the copies are \texttt{equal by a nested decision (\allowbreak{}lemma)}.
\item Settled: \texttt{peeled the last copy off both chains; the residuals are the hypothesis diagrams in the same places and the copies agree: W}; else \texttt{the peeled chain copies differ: W}.
\end{enumerate}
\end{algorithm}

\begin{signature}
\texttt{\_\allowbreak{}side\_\allowbreak{}replaced(\allowbreak{}own\_\allowbreak{}successor,\allowbreak{} own\_\allowbreak{}k,\allowbreak{} other\_\allowbreak{}k,\allowbreak{} k,\allowbreak{} *,\allowbreak{} last=False) -> Diagram | str}

Peel \code!own_successor! (with \code!last!, the joint peel with every first-peeling chain at its last copy; else \py!_peel_step_index!). Read the residual as \code!own_k! through cut wires by \py!expose_cut! over \texttt{copy\_\allowbreak{}candidates(\allowbreak{}exposed,\allowbreak{} own\_\allowbreak{}successor's nodes,\allowbreak{} copy)}, and replace the cut region by \code!other_k! (\py!replace_cut_region!). Failures return the texts \texttt{no chain peels its first copy}, the \py!BangBoxError! text, \code!the chain box does not peel!, \texttt{the residual is not the hypothesis diagram through its cut wires}, or the \py!InductionDomainError! text.
\end{signature}

\paragraph{Closing by the chain core.} \texttt{\_\allowbreak{}close\_\allowbreak{}by\_\allowbreak{}chain\_\allowbreak{}core(\allowbreak{}replaced,\allowbreak{} target,\allowbreak{} H,\allowbreak{} k,\allowbreak{} max\_\allowbreak{}steps)} computes \texttt{chain\_\allowbreak{}cores(\allowbreak{}H,\allowbreak{} replaced,\allowbreak{} target,\allowbreak{} k)}; when that gives a text, it falls back to \texttt{\_\allowbreak{}scope\_\allowbreak{}cores(\allowbreak{}H,\allowbreak{} replaced,\allowbreak{} target,\allowbreak{} k)}: the node scope of $H$'s chain !-box over $k$ as an open diagram (\py!_open_region!) and the scopes of the other two diagrams' chain !-boxes, all three of one multiplicity (else \py!None!, and the text stands). It then calls the kernel rule \texttt{local\_\allowbreak{}contraction(\allowbreak{}replaced,\allowbreak{} core\_\allowbreak{}a,\allowbreak{} target,\allowbreak{} core\_\allowbreak{}b,\allowbreak{} opened,\allowbreak{} max\_\allowbreak{}steps)} (\cref{sec:kernel-rules}). A \py!KernelError!, \py!InductionDomainError!, or \py!ValidateError! fails with \code!chain core: E!; success reads \texttt{both hold the other side's chain core (\allowbreak{}box and head) in the same places and the parts off it agree under symbolic contraction}.

\paragraph{Closing by lemmas.} \texttt{\_\allowbreak{}close\_\allowbreak{}by\_\allowbreak{}lemmas(\allowbreak{}replaced,\allowbreak{} target,\allowbreak{} source,\allowbreak{} k)} calls \texttt{close\_\allowbreak{}with\_\allowbreak{}lemmas(\allowbreak{}replaced,\allowbreak{} target,\allowbreak{} source,\allowbreak{} k)} of \module{archytaszx.semantics.lemmas} (\cref{ch:proofs}: library lemmas rewritten on both sides, the other side re-peeled at its chains' last copies, and the parts off the chain cores compared by symbolic contraction). A \py!KernelError!, \py!LemmaError!, \py!InductionError!, \py!BangBoxError!, or \py!ValidateError! fails with \code!lemmas: E!.

\begin{algorithm}[\texttt{\_\allowbreak{}chain\_\allowbreak{}window}]\label{alg:kernel-chain-window}
Returns \py!(settled, reason)!; the reasons of failed attempts are joined by \code!; !.
\begin{enumerate}
\item Try \cref{alg:kernel-chain-copies}; on success return it.
\item The two \emph{orientations} are (left successor, left hypothesis, right hypothesis, right successor) and its mirror. For (\code!last!, orientation) in (false, first), (false, second), (true, first), (true, second): when the own successor is not chain owned, record \texttt{no chain box carries the step index} (only when \code!last! is false) and skip; with \code!last!, skip when the own successor has no first peeler. Compute \texttt{\_\allowbreak{}side\_\allowbreak{}replaced(\allowbreak{}.\allowbreak{}.\allowbreak{}.\allowbreak{},\allowbreak{} last=last)} (a failure is recorded, prefixed \code!last copy: ! with \code!last!); the targets are the joint peel of the other successor (with \code!last!, its first peelers peeling their last copy), when it peels, and then the other successor itself. For each target close the window; on failure close by the chain core of the other hypothesis. A success returns \texttt{peeled the last chain copy; the residual is the hypothesis diagram through its cut wires,\allowbreak{} replaced by the other side's; W}.
\item For (\code!last!, orientation) in (false, first), (false, second), (true, first), (true, second): skip when the other successor is not chain owned, when \texttt{\_\allowbreak{}side\_\allowbreak{}replaced(\allowbreak{}.\allowbreak{}.\allowbreak{}.\allowbreak{},\allowbreak{} last=last)} fails, or when \py!peel_peelable! of the other successor is \py!None!; otherwise close by lemmas against that peel with the other successor as source. A success returns \texttt{peeled one copy (\allowbreak{}chains at their last copy); the residual .\allowbreak{}.\allowbreak{}.\allowbreak{}} (the parenthesis only with \code!last!).
\item \emph{Plain !-box.} For each orientation whose own successor is not chain owned but has an owner of $k$: compute \py!_side_replaced! (failure: \code!plain box: W!), take the targets as in step~2 without \code!last!, and close the window against each; a success returns \texttt{peeled one copy of a plain box; the residual is the hypothesis diagram through its cut wires,\allowbreak{} replaced by the other side's; W}.
\item \emph{Radix window.} For each orientation, \texttt{\_\allowbreak{}close\_\allowbreak{}by\_\allowbreak{}radix(\allowbreak{}own\_\allowbreak{}successor,\allowbreak{} own\_\allowbreak{}k,\allowbreak{} other\_\allowbreak{}successor,\allowbreak{} other\_\allowbreak{}k,\allowbreak{} k)}: the own successor must be chain owned with a first peeler (\texttt{radix: no first-peeling chain}); then \texttt{radix\_\allowbreak{}window.\allowbreak{}step(\allowbreak{}own\_\allowbreak{}successor,\allowbreak{} None,\allowbreak{} own\_\allowbreak{}k,\allowbreak{} other\_\allowbreak{}successor,\allowbreak{} None,\allowbreak{} other\_\allowbreak{}k,\allowbreak{} k)} of \module{archytaszx.semantics.radix_window} runs (first-copy peels exposed by library lemmas at the chain !-boxes, aligned rows merged, boxed spiders fused, and the other side's first-copy peel met by isomorphism; \cref{ch:proofs}). A \py!RadixWindowError!, \py!BangBoxError!, \py!ValidateError!, or \py!InductionError! records \code!radix: E!; success returns \texttt{every box peeled at its first copy; library lemmas expose the hypothesis at the chain boxes,\allowbreak{} replaced by the other side's; aligned boxes merged,\allowbreak{} boxed spiders fused,\allowbreak{} and the other side's first-copy peel met by isomorphism}.
\item Return \py!(False, reasons)!.
\end{enumerate}
\end{algorithm}

\begin{signature}
\texttt{offset\_\allowbreak{}obligation(\allowbreak{}obligation) -> tuple[InductionObligation,\allowbreak{} str] | None}

\py!None! unless a successor \py!peels_shifted! at $k$. Otherwise, with $s$ the \py!offset_symbol! of the four step diagrams, the obligation whose successors read their copy index as $\text{index}+s$ and whose hypothesis diagrams read it as $\text{index}+s+1$ (\py!induction_offset!), and the name of $s$; \py!None! when a \py!BangBoxError! or \py!InductionError! arises.

\texttt{discharge\_\allowbreak{}chain\_\allowbreak{}window(\allowbreak{}obligation,\allowbreak{} *,\allowbreak{} max\_\allowbreak{}steps=64) -> TierOutcome}

Run \py!_chain_window!; on success return a settled \texttt{CHAIN\_WINDOW} outcome with its reason. Otherwise, when \py!offset_obligation! gives a shifted obligation, run \py!_chain_window! on it; success reads \texttt{successors reading index + s,\allowbreak{} hypotheses index + s + 1 for every s: W}; failure appends \code!; offset s: W! to the reason. Then, for each pair $(a, v)$ of \texttt{parameter\_\allowbreak{}candidates(\allowbreak{}obligation)} (\cref{sec:kernel-parameters}) whose \py!parameter_obligation! exists, run \py!_chain_window! on it; success reads \texttt{claim for every a,\allowbreak{} hypotheses at a = v: W}; failure appends \code!; a = v: W!. The outcome records no peeled successors.
\end{signature}

\paragraph{Why the offset.} A chain that peels its first copy and reads a copy index leaves, after the peel, a residual whose copies read indices $1,\ldots,k$ instead of $0,\ldots,k-1$. Inducting over the claim ``both sides with the copy index read as $\text{index}+s$, for every $s$'' makes the residual an instance of the hypothesis at $s+1$. The kernel's \textsc{InductionOffset} rule (\cref{sec:kernel-rules}) checks this form.

\subsection{Phase parameters}\label{sec:kernel-parameters}

A successor's residual can also be the hypothesis diagram only at another value of a real symbol read in phases. The chain window and the induction-theorem builder then induct on the claim for every real value of that symbol, the hypothesis diagrams taken at the value the successor gives it.

\begin{definition}[Generalized constants]\label{def:kernel-generalized-constant}
\py!generalized_constant(left, right, index) -> (Diagram, Diagram, str, sympy.Expr) | None! collects the magnitudes $|c|$ of the rational coefficients of every term that reads \code!index! in the phase expressions of both sides (pair phases, phase-vector entries, and rules). It requires exactly one magnitude $c$, rational, and neither $0$ nor $1$; then, with a fresh parameter symbol \code!a! (\py!_fresh_symbol_name!), every such term $\pm c\cdot r$ becomes $\pm a\cdot r$ on both sides, and it returns the two lifted sides, the parameter name, and $c$. The decision procedure's generalized induction route (\cref{ch:proofs}) inducts on the lifted claim and takes it at $a=c$ by \py!kernel.phase_instance!.
\end{definition}

\begin{definition}[Phase parameter]\label{def:kernel-phase-parameter}
A \term{phase parameter} of a diagram (\py!_phase_symbols(D)!) is a free symbol read in some node phase and read nowhere else: not in the scalar, a multiplicity, a port dimension, or a copy scalar, and not a copy index. \py!parameter_symbol(name)! is the SymPy symbol of that name with \code!real=True!.
\end{definition}

\begin{signature}
\texttt{substitute\_\allowbreak{}parameter(\allowbreak{}diagram,\allowbreak{} name,\allowbreak{} value,\allowbreak{} *,\allowbreak{} normalize=True) -> Diagram} --- the diagram with \code!name! replaced by \code!value! in every node phase that reads it; with \code!normalize!, each turns expression of the phase's entries, rule, or pair is then expanded and its powers combined (\py!_mapped_phase!, through \py!_phase_map!, which keeps the weights of a weighted phase). Raises \py!InductionDomainError! \texttt{'a' is read outside node phases} when the name is free in the diagram but is no phase parameter of it.

\texttt{generalized\_\allowbreak{}constant(\allowbreak{}left,\allowbreak{} right,\allowbreak{} index) -> (\allowbreak{}Diagram,\allowbreak{} Diagram,\allowbreak{} str,\allowbreak{} Expr) | None} --- collects the coefficient magnitudes $|c|$ of the additive terms, reading \code!index!, of every phase expression (entries, rule, and pair phases) of both diagrams. When there is exactly one, rational and neither $0$ nor $1$, it takes a fresh name $a$ (\py!_fresh_symbol_name("a", ...)!), replaces every term $\pm c\,t$ by $\pm a\,t$ in every phase of both diagrams, and returns the two lifted diagrams, $a$, and $c$; else \py!None!.

\texttt{parameter\_\allowbreak{}value(\allowbreak{}successor,\allowbreak{} hypothesis,\allowbreak{} name,\allowbreak{} avoid) -> Expr | None} --- the value $v$ with every phase of \code!hypothesis! at $a=v$ equal to the phase of the same node id in \code!successor!. For each node whose phase reads $a$, the node of the same id must exist in the successor with a phase of the same class; for a plain phase vector, neither weighted, with the same entry indices and both or neither with a rule. Each pair of turns expressions (own, target) without $a$ must agree after expansion; one with $a$ must be linear in $a$ with a slope free of $a$, and $v=(\text{target}-\text{own}|_{a=0})/\text{slope}$ (simplified, expanded, powers combined) must be free of \code!avoid! and of the phase indices, and equal for every pair. It is \py!None! when any check fails, when no phase reads $a$, when $v$ simplifies to $a$, or when $v$ is not real.

\texttt{parameter\_\allowbreak{}candidates(\allowbreak{}obligation) -> list[(\allowbreak{}str,\allowbreak{} Expr)]} --- for each phase parameter of either hypothesis diagram, by name, not in the witness and not the index or the step symbol, the values \py!parameter_value! gives on the left pair and then on the right pair (avoiding the successors' copy indices), repeats dropped.

\texttt{parameter\_\allowbreak{}obligation(\allowbreak{}obligation,\allowbreak{} name,\allowbreak{} value) -> InductionObligation | None} --- the obligation with both hypothesis diagrams at $a=v$ (\py!substitute_parameter!); \py!None! on a \py!BangBoxError! or \py!InductionError!.
\end{signature}

The kernel's parametric induction (\textsc{InductionParam}, \cref{sec:kernel-rules}) checks the form.

\section{The Tree Window and W Fusion}\label{sec:kernel-tree}

A \term{tree !-box} of depth $\mu$ and branching $\beta$ expands to a $\beta$-ary tree of copies; a \term{fed !-box} of multiplicity $\beta^{\mu_T}$ is fed by the leaves of a tree !-box $T$ (\cref{ch:diagrams}).

\paragraph{Tree helpers.} \py!_tree_owner(D, k)! is the single !-box whose multiplicity reads $k$ when it is a tree !-box. \py!_slot_counts(D, refs)! is the list of expanded replication counts (\py!_port_multiplicity!) of the entries, \py!None! when one does not resolve.

\begin{signature}
\texttt{unfed(\allowbreak{}exposed,\allowbreak{} tree\_\allowbreak{}id) -> Diagram}

\code!exposed! without the !-box that the tree !-box \code!tree_id! feeds: in the output list and then the input list, the block of entries on the fed !-box's nodes must be one contiguous run of length $\beta$ (else \py!InductionDomainError! \texttt{the fed box's boundary slots are not its leaf block}) and is replaced by the tree's child ports; the fed !-box and its nodes are removed. With no fed !-box: \code!tree box t feeds no box!.
\end{signature}

\begin{algorithm}[\texttt{fed\_\allowbreak{}replacement}]\label{alg:kernel-fed-replacement}
Output: \code!other_k! with the fed !-box of \code!exposed! grafted on its leaf run.
\begin{enumerate}
\item Let $T$ be the tree !-box, $F$ the !-box it feeds (none, or none with a feed port: \code!tree box t feeds no box!), and $\beta$ the number of $T$'s child ports.
\item The residual \py!unfed(exposed, T)! must be isomorphic to \code!own_k! (\texttt{the residual is not the hypothesis diagram}).
\item The list is the output list when $T$'s first child port is an output, else the input list. Let $s$ be the position of $T$'s first child port in the residual's list and $w=|\text{other}|-|\text{own}|+\beta$. With the slot counts of both lists, $w\ge1$ is required (\texttt{the boundary slot counts do not resolve}), and the counts before $s$ and after the leaf run must agree (\texttt{the two hypothesis diagrams place the leaf run differently}). The \emph{middle} is \code!other_k!'s entries $s,\ldots,s+w-1$.
\item Copy $F$'s nodes and inner wires into a copy of \code!other_k!; the \emph{block} is $F$'s entries of \code!exposed!'s list, mapped. Let $\lambda=\beta^{\mu_T}$.
\item When the middle is one entry held by a port-scoped !-box of multiplicity $\lambda$ (expanded difference zero): that !-box must be neither nested nor a parent (\texttt{the leaf leg's port box is nested}); remove it and wire the middle entry to the image of $F$'s feed port. Else, when a tree !-box of \code!other_k! with $\beta$ child ports equal to the middle, feeding nothing, carries $T$'s multiplicity: it becomes the feeder. Else raise \texttt{the other hypothesis diagram's leaf run is neither one leg port-boxed b\^{}k times nor a tree box's child run}.
\item The list becomes the entries before $s$, the block, and the entries after the middle. Add the image of $F$ without feed; with a feeder, set its feed and feed port. Multiply the scalar by $\text{exposed.scalar}/\text{residual.scalar}$ and validate.
\end{enumerate}
\end{algorithm}

\subsection{W fusion}

\begin{figure}[H]
\centering
\begin{tikzpicture}[zxd]
  \node[wnode] (w) at (0,-0.6) {};
  \draw[wire] (w) -- (0,-1.4);
  \draw[wire] (w) -- (-0.7,0.3);
  \node[zs] (g) at (0.7,0.5) {$w$};
  \node[wnode] (w2) at (0.7,1.3) {};
  \draw[wire] (w) -- (g); \draw[wire] (g) -- (w2);
  \draw[wire] (w2) -- (0.4,2.1); \draw[wire] (w2) -- (1.0,2.1);
  \node[bang,fit=(g)(w2),minimum width=1.2cm] (B) {};
  \node[banglabel] at (B.north west) {$M$};
  \node at (2.2,0.4) {$=$};
  \node[wnode] (u) at (3.8,-0.1) {};
  \node[zs] (z) at (3.8,-0.9) {$w$};
  \draw[wire] (z) -- (3.8,-1.6); \draw[wire] (z) -- (u);
  \draw[wire] (u) -- (3.8,0.9);
  \node[bang,minimum width=0.6cm,minimum height=0.5cm] (P) at (3.8,0.6) {};
  \node[banglabel] at (P.north west) {$Mm$};
  \node[font=\tiny] at (5.0,0) {$\times\,(c\,s)^M$};
\end{tikzpicture}
\caption{W fusion: copies of $c\cdot W(1\to m)$ after a Z with weights $w$, fed by one W leg, fuse into that leg port-boxed $Mm$ times, the weighted Z moved before the W node (shown for a W node with one output).}\label{fig:kernel-wfusion}
\end{figure}

\py!_w_with_weights(d, m, w, c)! is $c$ times a $W(1\to m)$ after a $Z(1\to1)$ whose phase vector has the weights $w$ (no phase when $w$ is empty), at dimension $d$, boundary input the Z's input, outputs the W's outputs in order.

\begin{algorithm}[\texttt{\_\allowbreak{}w\_\allowbreak{}body}]\label{alg:kernel-w-body}
Output: $(c, w)$ with the copy of the !-box exactly equal to $c\cdot W(1\to m)$ after $Z(w)$, or \py!None!.
\begin{enumerate}
\item The input's dimension $d$ must be concrete and every output's must equal it.
\item Extract the !-box's scope (no !-boxes) with input \code!inside! and outputs \code!outs!; it must have no free symbol.
\item Contract it symbolically; an unsupported shape, a domain error, or a variable-rank tensor gives \py!None!.
\item For $j=0,\ldots,d-1$, let $r_j$ be the simplified entry at the axis values $(j,0,\ldots,0,j)$ in the tensor's axis order. When $r_0=0$, return \py!None!.
\item $w_j=r_j/r_0$ for $j\ge1$, with the weights equal to one dropped; the candidate is \texttt{\_\allowbreak{}w\_\allowbreak{}with\_\allowbreak{}weights(\allowbreak{}d,\allowbreak{} m,\allowbreak{} w,\allowbreak{} r\_\allowbreak{}0)}.
\item When the symbolic contractions of the extracted copy and the candidate agree, return $(r_0, w)$; else \py!None!.
\end{enumerate}
\end{algorithm}

\begin{algorithm}[\texttt{w\_\allowbreak{}fused}]\label{alg:kernel-w-fused}
For each !-box in id order:
\begin{enumerate}
\item Skip it unless it is node scoped and top level, has no links, tree data, feed, or pair parent, is the parent or pair parent of no !-box, and holds no boundary input.
\item Its cut wires must be exactly one, from an input port \code!inside! of the scope to a port $q$ of a W node in no !-box footprint.
\item Its boundary outputs must form one nonempty contiguous run \code!outs!, $m=|\text{outs}|$.
\item Its copy scalar $\kappa$ (one when absent) must not read its copy index.
\item \py!_w_body! must give $(c,w)$; with $w$ nonempty, the W node must have exactly one output.
\item Build the result: remove the !-box and its nodes; with $w$ nonempty, insert a Z of weights $w$ before the W node's input (on its boundary entry, or between its source and the W input); add a port-scoped !-box over $q$ of multiplicity $\mu m$ (\py!powsimp(expand(μ m))!); replace the run of outputs by $q$; set the scalar to $\text{scalar}\cdot(c\kappa)^{\mu}$; validate and return.
\end{enumerate}
With no such !-box, raise \py!InductionDomainError! \texttt{no box is a W fan on one W leg}. The rewrite is W associativity together with moving the weighted Z through the W node's only output.
\end{algorithm}

\begin{algorithm}[\texttt{discharge\_\allowbreak{}tree\_\allowbreak{}window}]\label{alg:kernel-tree-window}
For each orientation (as in \cref{alg:kernel-chain-window}):
\begin{enumerate}
\item The own successor's tree owner must exist (\texttt{no tree box carries the step index}).
\item Peel the leaf level (\texttt{peel\_\allowbreak{}one(\allowbreak{}own\_\allowbreak{}successor,\allowbreak{} owner)}) and build \texttt{fed\_\allowbreak{}replacement(\allowbreak{}exposed,\allowbreak{} owner,\allowbreak{} own\_\allowbreak{}k,\allowbreak{} other\_\allowbreak{}k)}; a \py!BangBoxError!, \py!InductionDomainError!, or \py!ValidateError! records its text.
\item The targets are the other successor with its tree owner peeled (when it has one and the peel succeeds), then the other successor. Close the window against each; success returns settled with \texttt{peeled the tree's leaf level; the rest is the hypothesis diagram,\allowbreak{} replaced by the other side's with the leaf level on its leaf run; W} and the exposed diagram as \py!left_peeled! when the peeled side is the left one, else as \py!right_peeled!.
\item Otherwise W-fuse the replaced diagram (\py!w_fused!; failure records \code!W fusion: E!) and close the window against each target; success adds \texttt{the leaf copies are a weighted W fused into the W leg feeding them} to the reason.
\end{enumerate}
Return unsettled with the joined reasons.
\end{algorithm}

\section{The Radix Window}\label{sec:kernel-radix-window}

\module{archytaszx.semantics.radix_window} closes a step case whose chain boxes peel at their first copy: library lemmas rewrite the re-peeled successor until the hypothesis is a cut region, the region is replaced by the other side's hypothesis, and the result is tidied and met with the other side's first-copy peel. Every theorem it returns is built by kernel rules (\cref{sec:kernel-rules}). A failure raises \py!RadixWindowError! with its reason; \py!MAX_ROUNDS! $=12$.

\begin{algorithm}[Exposing the hypothesis]\label{alg:kernel-radix-expose}
\py!_first_peels(source, k)! is \py!kernel.repeel(source, k, a_first=B)!, $B$ every box that is no pair box: each such box peeled at its first copy. \py!expose(peeled, first, own_k, k)! starts from $t=$ \py!kernel.symm(first)!, whose left side must be \py!_same! as \py!peeled! when given (\code!the re-peel does not end at the peeled successor!). In each of at most \py!MAX_ROUNDS! rounds:
\begin{enumerate}
\item For \py!joint! false, then true: the \emph{core} is the nodes of every box with links whose multiplicity reads $k$, grown by every unboxed part (one node, or with \py!joint! a connected set of unboxed nodes outside the core) all of whose wire neighbors outside the part lie in the core, until no part joins. When the core is nonempty and \py!find_cut_region(t.rhs, core, own_k)! (\cref{ch:proofs}) finds a region, return $(t,\text{region})$.
\item Otherwise rewrite with the first library lemma (\cref{tab:proofs-lemmas}) that applies to $t$'s right side and is not being proved (\py!_lemma_round!; \py!LemmaError!, \py!KernelError!, and \py!InductionDomainError! skip a lemma). When none applies, try the pair splits (\py!_pair_splits!): for each Z spider with at least three outputs, four ports, or a phase, each output wired to a \code!Zp! node is split off through the first other output (\py!_split_off!, \py!_split_theorem!), and the first split after which a lemma applies is taken. When nothing applies, raise \code!no library lemma applies and the core is not the hypothesis!.
\end{enumerate}
After the last round raise \code!the core is not the hypothesis after the lemma rounds!.
\end{algorithm}

\begin{algorithm}[Closing]\label{alg:kernel-radix-close}
\py!close(replaced, other_source, other_peeled, k)!: from \py!kernel.refl(replaced)!, apply in order
\begin{enumerate}
\item \py!_align_merges!: while some alignment root merges (\py!align_merged!, \py!kernel.align_merge!), merge it;
\item \py!_slot_groups!: keep in each box's groups only the boundary nodes, the groups sorted by their first boundary position, and \py!kernel.regroup! when this changes anything;
\item \py!_box_fusions!: while some unboxed-phase Z spider $\ell$ in a box, with a boundary port, is wired to a Z spider $h$ in exactly the same boxes, fuse $\ell$ into $h$ (\py!_fuse_into!: $h$'s port toward $\ell$ becomes $\ell$'s boundary port, $\ell$'s other ports are appended to $h$'s sides, box scopes, groups, and pair ports follow) and justify the step by the first \py!boxed_splits! candidate that \py!kernel.iso! and \py!kernel.box_split! accept;
\item \py!_tidy!: while possible, fuse two unboxed Z spiders joined by a wire (\py!_fusion!), else remove an unboxed identity (\code!identity_removal!);
\item \py!_slot_groups! again.
\end{enumerate}
Then $m=$ \py!_first_peels(other_source, k)! must end at \py!other_peeled! when given (\code!the other re-peel does not end at the other peeled successor!); its left side is tidied and slot-grouped to $u$; the result is $t$, \py!kernel.iso! of the two tidied sides, $u$ reversed, and $m$, joined by \py!kernel.trans!.
\end{algorithm}

\begin{signature}
\py!step(own_source, own_peeled, own_k, other_source, other_peeled, other_k, k) -> tuple[Theorem, CutRegion, Theorem]!
\end{signature}
\noindent \py!expose! on the own side, \py!replace_cut_region! of the exposed region by \py!other_k!, and \py!close!; a \py!KernelError! or \py!InductionDomainError! becomes \py!RadixWindowError!. The induction machinery calls it from its cut window (\cref{sec:kernel-induction}).

\begin{algorithm}[\texttt{regrouped\_induction(left, right, index, k)}]\label{alg:kernel-regrouped-induction}
\begin{enumerate}
\item Slot-group both sides (\py!_slot_grouped!): each outermost box without groups and pair partner that holds at least two boundary nodes gets one group per boundary node, in boundary order.
\item For the own side $0$, then $1$: split the own side's chain boxes (\py!_chain_splits!: while some outermost chain box's nodes fall into at least two parts connected by wires or links (\py!_chain_parts!, ordered by group rank), \py!chain_splitted! and \py!kernel.chain_split!); no split raises \code!no chain box splits!. Run \py!step! on the renamed successors, peeled successors, and hypotheses of the split own side and the grouped other side over \code!index! and \code!k!; prove the base case by \py!kernel.contraction!; build \py!kernel.induction! from the start side with the exposing theorem, the cut, and the step; prefix the splits, reverse for own side $1$, and return \py!kernel.regroup! of it to the original \code!left! and \code!right!.
\item A \py!RadixWindowError!, \py!KernelError!, or \py!InductionDomainError! records \code!side I: E!; when both sides fail raise \py!RadixWindowError! of the joined reasons.
\end{enumerate}
\end{algorithm}

\section{Lemma Decisions and Nesting}\label{sec:kernel-lemma-decision}

\begin{signature}
\texttt{lemma\_\allowbreak{}decision(\allowbreak{}left,\allowbreak{} right) -> Decision | None}

When both diagrams carry a free multiplicity and the context variable \py!_LEMMA_DEPTH! (default $0$) is below $2$ (\py!_MAX_LEMMA_DEPTH!): increment it, run \texttt{decide\_\allowbreak{}equal(\allowbreak{}left,\allowbreak{} right,\allowbreak{} samples=[])} (\cref{ch:proofs}; no oracle sample), restore it, and return the decision when its verdict is \texttt{EQUAL}. A \py!DecideError!, \py!InductionError!, \py!BangBoxError!, or \py!ValidateError! gives \py!None!, as does every other verdict.
\end{signature}

\begin{table}[H]
\centering\small
\begin{tabularx}{\linewidth}{@{}lrX@{}}
\toprule
Bound & Value & What it limits \\
\midrule
\texttt{\_MAX\_LEMMA\_DEPTH} & $2$ & nested \texttt{decide\_equal} calls of lemma decisions \\
\texttt{INDUCTED\_DEPTH} & $2$ & contraction steps re-proved by induction inside one opaque Lean export \\
\texttt{max\_depth} of the induction rung & $2$ & passed by \texttt{\_induction\_rung\_run}; the base case runs at one less \\
\texttt{\_MAX\_GREEDY\_STEPS} & $16$ & exposing steps of one greedy run \\
\texttt{\_MAX\_CUT\_PLACEMENTS} & $2000$ & cut placements per boundary list \\
\bottomrule
\end{tabularx}
\caption{Nesting and search bounds of the induction machinery.}\label{tab:kernel-bounds}
\end{table}

\section{The Kernel Theorem of an Induction}\label{sec:kernel-indthm}

\module{archytaszx.semantics.induction_kernel} builds the proof-kernel theorem of an \texttt{EQUAL} that the induction rung found, from the rung's own data: the proved \py!InductionResult!, the side order, the base pair and its decision, and the rule set. It never trusts the tier's verdict: every piece goes through a kernel rule, which re-checks it.

\begin{signature}
\texttt{class NoInductionTheorem(\allowbreak{}Exception)} --- why an induction verdict has no kernel theorem.

\texttt{induction\_\allowbreak{}theorem(\allowbreak{}left,\allowbreak{} right,\allowbreak{} result,\allowbreak{} reversed\_\allowbreak{},\allowbreak{} base\_\allowbreak{}pair,\allowbreak{} base\_\allowbreak{}decision,\allowbreak{} rules=None) -> tuple[Theorem | None,\allowbreak{} str]}

$L=R$ from a proved \code!result! (run on $R, L$ when \code!reversed_!) and the decision of \code!base_pair! at index $0$; or \py!None! and why.
\end{signature}

\begin{algorithm}[\texttt{induction\_\allowbreak{}theorem}]\label{alg:kernel-induction-theorem}
\begin{enumerate}
\item $(\textit{first}, \textit{second}) \gets (R, L)$ when \code!reversed_!, else $(L, R)$. The settling tier is the last settled tier of \code!result! (\py!_settling_tier!; none raises \texttt{no tier settled the step case}).
\item By its discharge:
  \begin{itemize}
  \item \texttt{UNIFORM\_REWRITE}: the theorem is \texttt{\_\allowbreak{}uniform(\allowbreak{}first,\allowbreak{} second,\allowbreak{} result,\allowbreak{} tier)}.
  \item \texttt{INDUCTION\_REWRITE}: $B\gets$ \py!_base(...)!; the theorem is \texttt{\_\allowbreak{}peeled(\allowbreak{}first,\allowbreak{} second,\allowbreak{} result,\allowbreak{} tier,\allowbreak{} B,\allowbreak{} rules)}.
  \item \texttt{CHAIN\_WINDOW}, \texttt{TREE\_WINDOW}: $B\gets$ \py!_base(...)!; the theorem is \texttt{\_\allowbreak{}windowed(\allowbreak{}first,\allowbreak{} second,\allowbreak{} result,\allowbreak{} w,\allowbreak{} B,\allowbreak{} rules)} with $w$ \code!"cut"! or \code!"tree"!.
  \item \texttt{HYPOTHESIS\_REWRITE}: $B\gets$ \py!_base(...)!; the theorem is the first of \texttt{\_\allowbreak{}windowed(\allowbreak{}first,\allowbreak{} second,\allowbreak{} result,\allowbreak{} w,\allowbreak{} B,\allowbreak{} rules,\allowbreak{} contract=x)} over $w\in\{$\code!"cut"!, \code!"region"!$\}$, $x\in\{$\code!False!, \code!True!$\}$ (\py!_hypothesis!).
  \item \texttt{SYMBOLIC\_CONTRACTION}: $B\gets$ \py!_base(...)!; the theorem is \texttt{\_\allowbreak{}contracted(\allowbreak{}first,\allowbreak{} second,\allowbreak{} result,\allowbreak{} B,\allowbreak{} rules)}.
  \item Any other: \texttt{the step case was settled by D,\allowbreak{} which records no rewrite steps}.
  \end{itemize}
\item When \code!reversed_!, apply \textsc{Symm}. Return \py!(theorem, "")!.
\item A \py!NoInductionTheorem! returns \py!(None, message)!; a \py!KernelError! returns \texttt{the kernel refused a piece: E}; any other exception returns \code!building it raised T: E!.
\end{enumerate}
\end{algorithm}

\paragraph{Uniform rewrite.} \texttt{\_\allowbreak{}uniform(\allowbreak{}first,\allowbreak{} second,\allowbreak{} result,\allowbreak{} tier)} requires the hypothesis diagram \texttt{left\_at\_k} to have exactly the node ids of \textit{first} (\texttt{the hypothesis diagram does not keep the left side's node ids}), transports the tier's steps onto \textit{first} with the identity id map (\py!transport!, below), chains them (\py!kernel.chain!), and closes with \textsc{Iso} onto \textit{second}. No induction inference appears: the steps hold for every value of the index.

\paragraph{Base case.} \texttt{\_\allowbreak{}base(\allowbreak{}start,\allowbreak{} goal,\allowbreak{} n,\allowbreak{} pair,\allowbreak{} decision,\allowbreak{} reversed\_\allowbreak{})} reads the base decision as a theorem (\texttt{kernel.\allowbreak{}from\_\allowbreak{}decision(\allowbreak{}decision,\allowbreak{} *pair)}). When the kernel refuses it, it runs the goal-directed search \py!prove(*pair)! (\cref{ch:proofs}) and checks its proof certificate, or else its induction proof, with \py!kernel.check!; with neither, it raises \code!base case: E; search: R! (\code!no certificate! when the search gives no reason). Under \code!reversed_! it applies \textsc{Symm}. It returns $\textsc{Iso}(\textit{start}|_{n=0}, \Theta.\mathrm{lhs})\mathbin{;}\Theta\mathbin{;}\textsc{Iso}(\Theta.\mathrm{rhs}, \textit{goal}|_{n=0})$, where $;$ is \textsc{Trans}.

\begin{algorithm}[\texttt{\_\allowbreak{}peeled}]\label{alg:kernel-peeled-theorem}
\begin{enumerate}
\item The tier must record both peeled successors and both copies (\texttt{the peeling tier recorded no peeled successors}).
\item $P\gets$ \texttt{induction\_\allowbreak{}successor(\allowbreak{}start,\allowbreak{} n,\allowbreak{} k)}, $P'\gets$ \texttt{induction\_\allowbreak{}successor(\allowbreak{}goal,\allowbreak{} n,\allowbreak{} k)}; either \py!None!: \texttt{a successor has no single box to peel}.
\item $E\gets$ the tier's left peeled successor; the exposing theorem is \textsc{Refl}$(P)$ when $P\equiv E$, else \textsc{Iso}$(P,E)$. The region is $E$'s nodes less the left copy; the replaced diagram $X\gets$ \texttt{replace\_\allowbreak{}region(\allowbreak{}E,\allowbreak{} region,\allowbreak{} induction\_\allowbreak{}hypothesis(\allowbreak{}start),\allowbreak{} induction\_\allowbreak{}hypothesis(\allowbreak{}goal))}.
\item Extract both copies with !-boxes ($c_L$ from $E$, $c_R$ from the right peeled successor $E'$). Find a rewrite derivation between them (\py!copy_derivation!).
\item With a derivation (forward steps $f$, backward steps $g$): $\textit{there}\gets$ \texttt{chain(\allowbreak{}X,\allowbreak{} transport(\allowbreak{}X,\allowbreak{} f,\allowbreak{} ids))}, $\textit{back}\gets$ \texttt{chain(\allowbreak{}E',\allowbreak{} transport(\allowbreak{}E',\allowbreak{} g,\allowbreak{} ids'))}, where \code!ids! maps the extracted copy's nodes, in sorted order, onto the copy's nodes in sorted order (\py!_ids!). The step is $\textit{there}\mathbin{;}\textsc{Iso}\mathbin{;}\textsc{Symm}(\textit{back})\mathbin{;}\textsc{Iso}(E',P')$.
\item Without one: try \textsc{Contraction}$(X, E')$; failing that, form the closed cut region of the left copy in $X$ (its boundary inputs and outputs on the copy, no cuts) and try \texttt{\_\allowbreak{}lemma\_\allowbreak{}congr(\allowbreak{}X,\allowbreak{} closed,\allowbreak{} c\_\allowbreak{}L,\allowbreak{} c\_\allowbreak{}R,\allowbreak{} E')}; failing that, raise \code!D; C; lemma: M!. The step is that theorem $;\ \textsc{Iso}(E',P')$.
\item Return \textsc{Induction}$(\textit{start},\textit{goal},n,k,\textit{base},\texttt{START},\textit{expose},\text{region},\textit{step})$.
\end{enumerate}
\end{algorithm}

\paragraph{Symbolic-contraction discharge.} \texttt{\_\allowbreak{}contracted(\allowbreak{}start,\allowbreak{} goal,\allowbreak{} result,\allowbreak{} base,\allowbreak{} rules)} tries \py!_windowed! with the windows \code!"cut"!, \code!"tree"!, and \code!"region"!, in that order, each with \code!contract=True!, and returns the first theorem; it raises the joined reasons \code!window: E! otherwise.

\begin{algorithm}[\texttt{\_\allowbreak{}windowed}]\label{alg:kernel-windowed}
\begin{enumerate}
\item The offsets are \py!None!, and also $s=$ \texttt{offset\_\allowbreak{}symbol(\allowbreak{}start,\allowbreak{} goal,\allowbreak{} left\_\allowbreak{}at\_\allowbreak{}successor,\allowbreak{} right\_\allowbreak{}at\_\allowbreak{}successor)} when \py!peels_shifted! holds of either side at $n$.
\item The combinations are each offset with no parameter, then no offset with each pair $(a,v)$ of \texttt{parameter\_\allowbreak{}candidates(\allowbreak{}result.\allowbreak{}obligation)} (\cref{sec:kernel-parameters}). For each combination and, inside it, each side (\texttt{START}, then \texttt{GOAL}), with \textit{own} the side's diagram and \textit{other} the other one; the label is the side's value, or \code!side, offset s!, with \code!, a = v! appended for a parameter:
  \begin{enumerate}
  \item With an offset, shift both by \py!induction_offset(., n, s)!. $P\gets$ \texttt{induction\_\allowbreak{}successor(\allowbreak{}own\_\allowbreak{}s,\allowbreak{} n,\allowbreak{} k)}, $P'\gets$ that of \textit{other}; an exception records \code!label: peeling raised T: E!, a \py!None! records \texttt{label: a successor has no single box to peel}.
  \item $H\gets$ \texttt{induction\_\allowbreak{}hypothesis(\allowbreak{}own,\allowbreak{} n,\allowbreak{} k)}, $H'\gets$ that of \textit{other}; with an offset, shifted by \texttt{induction\_\allowbreak{}offset(\allowbreak{}.\allowbreak{},\allowbreak{} k,\allowbreak{} s + 1)}; with a parameter, taken at $a=v$ by \py!substitute_parameter!.
  \item The exposing theorem starts as \textsc{Refl}$(P)$.
    \begin{itemize}
    \item Window \code!"cut"!: the copy is $P$'s nodes not in $H$, the region its complement. Find the cut region of the complement read as $H$; failing that, for each spider split of \py!_fan_split(P, copy)! find it in the split, the exposing theorem then being \py!_split_theorem!; failing that, when \py!copy_candidates! gives more than the copy, run \py!expose_cut! over the copy and the later candidates, the exposing theorem being \py!_steps_theorem! of the steps found. With none: \texttt{the residual is not the hypothesis through cut wires}. $X\gets$ \texttt{replace\_\allowbreak{}cut\_\allowbreak{}region(\allowbreak{}expose.\allowbreak{}rhs,\allowbreak{} cut,\allowbreak{} H,\allowbreak{} H')}.
    \item Window \code!"region"!: the region is $P$'s nodes that are $H$'s; $X\gets$ \texttt{replace\_\allowbreak{}region(\allowbreak{}P,\allowbreak{} region,\allowbreak{} H,\allowbreak{} H')}.
    \item Window \code!"tree"!: $P$ must hold exactly one fed !-box (\texttt{the peeled successor holds no single fed box}); its feeder $T$ fixes the region, the nodes of \py!unfed(P, T)!, and $X\gets$ \py!fed_replacement(P, T, H, H')!.
    \end{itemize}
  \item Outside the tree window, the \emph{local} data are the triple of $X$'s nodes less the rest, the nodes of $P'$ that are nodes of $H'$, and $H'$, the rest being the exposed diagram's nodes off the region; in the cut window the \emph{chain} data are $(H', k)$.
  \item The step is \texttt{\_\allowbreak{}step\_\allowbreak{}theorem(\allowbreak{}X,\allowbreak{} P',\allowbreak{} rules,\allowbreak{} contract=contract,\allowbreak{} local=local,\allowbreak{} chain=chain)}. When it raises \py!NoInductionTheorem!: in the cut window, the step is \texttt{\_\allowbreak{}lemma\_\allowbreak{}step(\allowbreak{}X,\allowbreak{} P',\allowbreak{} induction\_\allowbreak{}renamed(\allowbreak{}other\_\allowbreak{}s,\allowbreak{} n,\allowbreak{} k),\allowbreak{} k,\allowbreak{} exc)} and the induction is built at once; in the region window it is re-raised and recorded by step~(g); in the tree window, apply \textsc{WFusion} to $X$ (a refusal raises \code!E; W fusion: E'!), close from its right side to $P'$ by \py!_step_theorem!, and compose.
  \item Return \textsc{Induction}$(\textit{start},\textit{goal},n,k,\textit{base},\text{side},\textit{expose},\text{region},\textit{step})$ with the cut (cut window), the tree id (tree window), the offset name (with an offset), and the parameter $(a,v)$ (with a parameter); the kernel's inference kind is then \code!induction_offset! or \code!induction_param! (\cref{sec:kernel-rules}).
  \item A \py!NoInductionTheorem!, \py!KernelError!, or \py!InductionDomainError! records \code!label: E! and continues.
  \end{enumerate}
\item In the cut window, try \py!_last_window! for \texttt{START} then \texttt{GOAL}; failures record \texttt{side,\allowbreak{} chains at their last copy: E}. Then try \py!_radix! for \texttt{START} then \texttt{GOAL}; failures record \code!side, radix window: E!.
\item Raise \py!NoInductionTheorem! with the joined reasons.
\end{enumerate}
\end{algorithm}

\begin{algorithm}[\texttt{\_\allowbreak{}last\_\allowbreak{}window}]\label{alg:kernel-last-window}
\begin{enumerate}
\item $S\gets$ \py!induction_renamed(own, n, k)!; the first peelers of $S$ at $k$ must be nonempty and both successors must exist (\texttt{no chain peels its first copy}).
\item $\textit{moved}\gets\textsc{Repeel}(S, k, \emptyset, \text{first peelers})$: from the joint peel at the first copies to the joint peel with the first peelers at their last copy. Its left side must be the successor (\texttt{the re-peel does not start at the peeled successor}).
\item Read the exposed diagram's complement of the copy as the own hypothesis through cut wires by \py!expose_cut! over \py!copy_candidates! (\texttt{the residual is not the hypothesis through cut wires}). The exposing theorem is $\textit{moved}\mathbin{;}$\py!_steps_theorem!.
\item Replace the cut region by the other hypothesis diagram; the step is \texttt{\_\allowbreak{}lemma\_\allowbreak{}step(\allowbreak{}X,\allowbreak{} P',\allowbreak{} induction\_\allowbreak{}renamed(\allowbreak{}other,\allowbreak{} n,\allowbreak{} k),\allowbreak{} k,\allowbreak{} .\allowbreak{}.\allowbreak{}.\allowbreak{})}; when that raises, the step is \texttt{\_\allowbreak{}scope\_\allowbreak{}step(\allowbreak{}X,\allowbreak{} P',\allowbreak{} induction\_\allowbreak{}renamed(\allowbreak{}other,\allowbreak{} n,\allowbreak{} k),\allowbreak{} H',\allowbreak{} k,\allowbreak{} exc)}.
\item Return \textsc{Induction} with that cut.
\end{enumerate}
\end{algorithm}

\paragraph{Scope step.} \texttt{\_\allowbreak{}scope\_\allowbreak{}step(\allowbreak{}replaced,\allowbreak{} target,\allowbreak{} source,\allowbreak{} hypothesis,\allowbreak{} k,\allowbreak{} failed)} re-raises \code!failed! when \code!source! has no first peeler at $k$. Otherwise \textit{moved} is \textsc{Repeel} of \textit{source} at $k$ from the joint peel to the one with the first peelers at their last copy, whose left side must be \code!target! (\texttt{F; scope: the re-peel does not start at the target}); the scopes of the chain !-boxes over $k$ (\py!_scope_cores! of \code!hypothesis!, \code!replaced!, and the re-peel's right side) must exist (\texttt{F; scope: no chain box over the index on all sides}); and the step is \textsc{LocalContraction} of \textit{replaced} and the re-peel's right side over the two scopes and the opened scope of \code!hypothesis!, followed by \textsc{Symm}(\textit{moved}) (a refusal or \py!InductionDomainError! raises \code!F; scope: E!).

\paragraph{Radix window.} \texttt{\_\allowbreak{}radix(\allowbreak{}start,\allowbreak{} goal,\allowbreak{} n,\allowbreak{} k,\allowbreak{} base,\allowbreak{} side)} calls \py!radix_window.step! of \module{archytaszx.semantics.radix_window} with, for the own side and then the other side, the renamed successor (\py!induction_renamed!), the peeled successor (\py!induction_successor!), and the hypothesis diagram, and $k$. That function returns an exposing theorem, a cut region, and a step theorem, all built by kernel rules; a \py!RadixWindowError! becomes \py!NoInductionTheorem!. The result is \textsc{Induction} over that cut.

\paragraph{Exposing theorems.} \texttt{\_\allowbreak{}steps\_\allowbreak{}theorem(\allowbreak{}peeled,\allowbreak{} steps)} starts at \textsc{Refl} and appends, by \textsc{Trans}, \textsc{LegMove}$(\text{before},\text{after},\text{copy\_before})$ for a move and \py!_split_theorem! for a split. \texttt{\_\allowbreak{}split\_\allowbreak{}theorem(\allowbreak{}peeled,\allowbreak{} split,\allowbreak{} copy,\allowbreak{} split\_\allowbreak{}copy)} returns \textsc{BoxSplit}$(\text{peeled},\text{split},\text{copy})$ when a !-box's node scope contains the nodes new to the split copy; otherwise the new leaf and the new hub must be unique, a spider-fusion match on exactly those two nodes is found, and the theorem is $\textsc{Symm}(\textsc{Rule}(\text{split}, \text{fusion step})\mathbin{;}\textsc{Iso}(\cdot,\text{peeled}))$; without such a match: \texttt{the split spider does not fuse back}.

\begin{algorithm}[\texttt{\_\allowbreak{}step\_\allowbreak{}theorem}]\label{alg:kernel-step-theorem}
The first success is returned; the failures are joined by \code!; !.
\begin{enumerate}
\item \textsc{Iso}$(X, P')$.
\item Unless \code!contract!: with \code!local!, a rewrite derivation between the parts off the two regions (\py!_local_derivation!), giving $\textit{there}\mathbin{;}\textsc{Iso}\mathbin{;}\textsc{Symm}(\textit{back})$ (failure: \code!local derivation: E!); without \code!local!, \texttt{copy\_\allowbreak{}derivation(\allowbreak{}X,\allowbreak{} P',\allowbreak{} rules)} chained on both sides likewise.
\item With \code!local!: \textsc{LocalContraction}$(X, r_X, P', r_{P'}, H')$ (\code!local: E!); then \py!_local_lemma! (\code!lemma: E!).
\item With \code!chain!: the chain cores of $(H', X, P')$; \textsc{LocalContraction} over the cores and the opened core of $H'$ (\code!chain core: E!).
\item \textsc{Contraction}$(X, P')$.
\end{enumerate}
\end{algorithm}

\paragraph{Scope steps.} \texttt{\_\allowbreak{}scope\_\allowbreak{}step(\allowbreak{}replaced,\allowbreak{} target,\allowbreak{} source,\allowbreak{} hypothesis,\allowbreak{} k,\allowbreak{} failed)} re-raises \texttt{failed} when \texttt{source} has no first peelers at $k$; otherwise re-peels \texttt{source} at its first peelers (\py!kernel.repeel(source, k, frozenset(), firsts)!), whose left side must be \py!_same! as \texttt{target} (\code!F; scope: the re-peel does not start at the target!). \py!_scope_cores! takes the chain owner over $k$ of the hypothesis, of \texttt{replaced}, and of the re-peel's right side; all three must exist with one multiplicity (\code!F; scope: no chain box over the index on all sides!); it returns the hypothesis's box scope as an open region and the two other scopes. The step is \py!kernel.local_contraction! of \texttt{replaced} at its scope against the re-peel's right side at its scope over the open region (a \py!KernelError! or \py!InductionDomainError! raises \code!F; scope: E!), followed by the reversed re-peel. \py!_close_by_chain_core! also falls back to \py!_scope_cores! when \py!chain_cores! gives a reason.

\paragraph{Lemma steps.} \texttt{\_\allowbreak{}lemma\_\allowbreak{}congr(\allowbreak{}context,\allowbreak{} cut,\allowbreak{} part,\allowbreak{} other\_\allowbreak{}part,\allowbreak{} target)} takes \texttt{lemma\_\allowbreak{}decision(\allowbreak{}part,\allowbreak{} other\_\allowbreak{}part)} (\texttt{no nested decision proves the parts equal} when \py!None!), reads it as the theorem $\Lambda$ by \py!from_decision!, and returns $\textsc{Congr}(\Lambda,\text{context},\text{cut})\mathbin{;}\textsc{Iso}(\cdot,\text{target})$. \texttt{\_\allowbreak{}local\_\allowbreak{}lemma(\allowbreak{}replaced,\allowbreak{} target,\allowbreak{} local)} finds the cut regions of the two regions read as the local hypothesis (\texttt{a region is not the hypothesis through cut wires}), takes the parts off them (\py!_copy_side! with the scalar ratio), and calls \py!_lemma_congr! at \py!rest_cut! of the replaced diagram. \texttt{\_\allowbreak{}lemma\_\allowbreak{}step(\allowbreak{}replaced,\allowbreak{} target,\allowbreak{} source,\allowbreak{} k,\allowbreak{} failed)} returns \texttt{close\_\allowbreak{}with\_\allowbreak{}lemmas(\allowbreak{}replaced,\allowbreak{} target,\allowbreak{} source,\allowbreak{} k)} of \cref{ch:proofs}, which builds a kernel theorem; a \py!KernelError!, \py!LemmaError!, or \py!InductionDomainError! raises \code!F; lemmas: E!.

\paragraph{Derivations between copies.} \texttt{\_\allowbreak{}local\_\allowbreak{}derivation(\allowbreak{}replaced,\allowbreak{} target,\allowbreak{} local,\allowbreak{} rules)} takes the parts off the two regions with their cut ends open (\py!kernel._plugged!), finds \py!copy_derivation! between them, and transports the forward steps onto \code!replaced! and the backward steps onto \code!target!, returning both chains.

\begin{signature}
\texttt{copy\_\allowbreak{}derivation(\allowbreak{}left,\allowbreak{} right,\allowbreak{} rules) -> tuple[tuple[RewriteStep,\allowbreak{} .\allowbreak{}.\allowbreak{}.\allowbreak{}],\allowbreak{} tuple[RewriteStep,\allowbreak{} .\allowbreak{}.\allowbreak{}.\allowbreak{}]]}

The steps of the normal forms of both (\py!normal_form(., rules=rules)!, \cref{ch:rewriting}) when the two normal forms pass \textsc{Iso}; otherwise (on a \py!RewriteError! or a refusal, recorded as \code!normal forms: E!) the two halves of a proof certificate of \py!prove(left, right)!; otherwise \py!NoInductionTheorem! \texttt{no rewrite derivation between the peeled copies (\allowbreak{}R)} with \code!search: R! appended.
\end{signature}

\begin{algorithm}[\texttt{transport}]\label{alg:kernel-transport}
Re-applies steps recorded on another diagram, whose node ids \code!ids! sends into \code!start!.
\begin{enumerate}
\item $W\gets\textit{start}$; the map $\phi\gets\textit{ids}$.
\item For each step at position $i$: the wanted nodes are $\phi$ of its consumed nodes (a missing node raises \texttt{step i consumes a node outside the copy}). Look up its rule; for each match of the rule's pattern in $W$, in \py!find_matches! order, apply it (skipping a \py!RewriteError!) and accept the first result whose consumed nodes are the wanted ones and whose new-node count equals the step's (none: \texttt{step i (\allowbreak{}R) has no match on the corresponding nodes}). Extend $\phi$ by the step's new nodes onto the result's new nodes, in order; $W\gets$ the result.
\item Return the accepted steps.
\end{enumerate}
\end{algorithm}

\begin{figure}[H]
\centering
\begin{tikzpicture}[x=1cm,y=1cm,every node/.style={font=\small}]
  \node (a) at (0,3) {$\mathit{own}(k{+}1)$};
  \node (b) at (3.4,3) {$P$};
  \node (c) at (6.6,3) {$E = C[\mathit{own}(k)]$};
  \node (d) at (6.6,1.4) {$C[\mathit{other}(k)] = X$};
  \node (e) at (3.4,1.4) {$P'$};
  \node (f) at (0,1.4) {$\mathit{other}(k{+}1)$};
  \draw[-{Latex}] (a) -- node[above] {peel} (b);
  \draw[-{Latex}] (b) -- node[above] {\textit{expose}} (c);
  \draw[-{Latex}] (c) -- node[right] {hypothesis} (d);
  \draw[-{Latex}] (d) -- node[above] {\textit{step}} (e);
  \draw[-{Latex}] (f) -- node[above] {peel} (e);
\end{tikzpicture}
\caption{The induction rule: the side's successor $P$, exposed as a context $C$ around the side at $k$; the hypothesis replaces it by the other side at $k$; the step reaches the other side's successor $P'$.}\label{fig:kernel-induction}
\end{figure}

\section{The Proof Kernel}\label{sec:kernel-kernel}

\module{archytaszx.semantics.kernel} defines the type \py!Theorem! and the only functions that construct it. Every certified \texttt{EQUAL} of the engine is a \py!Theorem! (or, for the symbolic-sum normalizer, a \py!Symbolic! record; \cref{sec:kernel-refute}), and every exported Lean proof of an equation starts from one.

\subsection{Trust model}\label{sec:kernel-trust}

\begin{enumerate}
\item \emph{Construction.} \texttt{Theorem.\allowbreak{}\_\allowbreak{}\_\allowbreak{}init\_\allowbreak{}\_\allowbreak{}(\allowbreak{}token,\allowbreak{} lhs,\allowbreak{} rhs,\allowbreak{} inference,\allowbreak{} assumptions=(\allowbreak{}),\allowbreak{} domain=(\allowbreak{}),\allowbreak{} origin=None,\allowbreak{} *,\allowbreak{} floors=True)} raises \py!KernelError! \texttt{a Theorem is built only by the kernel's inference rules} unless \code!token! is the module-private sentinel \py!_TOKEN! (an \py!object()!). The inference-rule functions of the module are the only callers that pass it. The guard is a module-boundary convention of Python: code that imports the private name can forge a theorem, and the kernel does not defend against it.
\item \emph{Checking.} Each inference rule re-checks its premises' shapes and its side conditions with the engine's own algorithms (rule replay with match rediscovery, graph isomorphism, comparison views, number-field certificates, symbolic contraction, path-sum reduction, peeling, splitting, and region extraction) and raises \py!KernelError! at the first violation. A theorem is therefore trusted to the extent that these algorithms are correct; no step is accepted on the word of the search that found it.
\item \emph{Lean.} The export of \cref{sec:kernel-lean} hands the theorem to the Lean~4 kernel~\cite{demoura2021lean4}. In an \emph{opaque} export each kernel-checked step is a Lean axiom over an opaque type \lean!Diagram! and Lean checks only how the steps compose. In a \emph{semantic} export (\cref{ch:lean-export-net,ch:lean-export}) the diagrams are tensor networks over $\mathbb{C}$ in the bundled Mathlib-based~\cite{mathlib2020} library and each rule step is an instance of a lemma proved there, so Lean re-checks the steps themselves. An axiom whose name begins \lean!unproved_! marks a step with no Lean proof (a discard rule); a semantic export containing one is replaced by a symbolic-sum proof when one exists.
\end{enumerate}

\begin{signature}
\py!class KernelError(Exception)! --- a proof piece the kernel does not accept. Every refusal of every rule, of \py!check!, of \py!from_decision!, and of \py!lean! raises it, with the messages of \cref{sec:kernel-refusals}.
\end{signature}

\subsection{Theorems and inferences}

\begin{definition}[Inference]\label{def:kernel-inference}
\py!Inference! (frozen, compared by identity) records the rule that built a theorem: \py!kind: str!, \texttt{premises: tuple[Theorem,\allowbreak{} .\allowbreak{}.\allowbreak{}.\allowbreak{}] = (\allowbreak{})}, \py!rule: str = ""! (the rewrite rule's name for kind \code!rule!), \py!detail: str = ""! (the consumed nodes \code!n3, n5! for kind \code!rule!, the substitution \code!x = e! for \code!instance!, the clause text for \code!subst!), and \py!data: dict[str, object]! (the rule-specific record of \cref{tab:kernel-inferences}).
\end{definition}

\begin{definition}[Theorem]\label{def:kernel-theorem}
A \py!Theorem! (slots \code!lhs!, \code!rhs!, \code!inference!, \code!assumptions!, \code!domain!, \code!origin!; compared by identity) states
\[
\Gamma \mathrel{;} A \vdash L = R,
\]
read as: for every assignment of natural numbers to the free multiplicities of $L$ and $R$ that satisfies every clause of the \term{domain} $\Gamma$, and for every value of the other symbols that satisfies the dimension equations $A$ (the \term{assumptions}, strings such as \code!d = 2*e!), $\sem{L}=\sem{R}$; for diagrams with a discard, the equality of their completely positive maps~\cite{selinger2007cpm}. \code!origin! is the unrestricted pair the domain's clauses were substituted into, \py!None! for an unrestricted theorem.
\begin{itemize}
\item \py!variables! is the sorted tuple of free multiplicities of both sides.
\item \py!size()! is $1$ plus the sizes of the premises, the number of inferences of the proof tree.
\item \py!axioms()! is the list of rule names of every \code!rule! inference in proof order, repeats kept.
\item \py!statement()! is the quantifier text of \cref{def:kernel-statement}.
\end{itemize}
The constructor sets $\Gamma$ to \texttt{\_\allowbreak{}with\_\allowbreak{}box\_\allowbreak{}floors(\allowbreak{}lhs,\allowbreak{} rhs,\allowbreak{} domain)} (\cref{sec:kernel-floors}) unless \code!floors=False!.
\end{definition}

\begin{definition}[Statement text]\label{def:kernel-statement}
\texttt{statement(\allowbreak{}domain,\allowbreak{} variables,\allowbreak{} assumptions=(\allowbreak{}))} joins with \code!, ! the parts below. Let $\Phi$ be the fresh symbols of the clauses together with the symbols of the equality clauses.
\begin{enumerate}
\item For each clause in order: a shift, side, or finite clause gives its hypothesis text (below); a dimension clause gives nothing here; a multiplicity equality $x=c$ gives \code!at x = c!; a lower bound $x\ge c$ gives \code!for all x >= c (x = e)!; an affine clause gives \code!for all x = e (m >= 0)!, where $e$ is the clause's expression (\code!m + 2!, \code!2*m + 1!) and $m$ its fresh symbol.
\item When some variables are not in $\Phi$, \code!for all v1, v2! is put first.
\item For each dimension clause that is no shift, side, or finite clause: an equality gives \code!at d = c!; a lower bound gives \code!for d >= c! when \code!c <= d! is among the assumptions, else \code!for every d!.
\end{enumerate}
A shift clause gives \code!for all x = e (R)! with $R$ the range text of its fresh symbol's domain, followed by \code!; ! and the range text of $x$'s own domain when it has one; a finite clause gives \code!at x in {v1, v2}!; a side clause gives \code!for R! with $R$ the range text of $x$. The \term{range text} of a domain over $x$ (\py!_range_text!) lists, comma separated, \code!x >= lo! (or \code!x > lo!), \code!x <= hi! (or \code!x < hi!), \code!x in Z! for $q=1$ or \code!x in (1/q)Z! for $q>1$, \code!x in o + tZ! for a lattice of step $t$ and offset $o$, and \code!x != v! per excluded value.
\end{definition}

\begin{example}[Statements]
The four calls below print the texts shown.
\begin{center}\small
\begin{tabular}{@{}ll@{}}
\toprule
Domain; variables; assumptions & \texttt{statement} \\
\midrule
$n\ge2$ (fresh $m$); $(m)$ & \texttt{for all n >= 2 (n = m + 2)} \\
$n=2m+1$; $(m)$ & \texttt{for all n = 2*m + 1 (m >= 0)} \\
$n=2$, dimension $d\ge3$; $(k)$; \texttt{3 <= d} & \texttt{for all k, at n = 2, for d >= 3} \\
dimension $d=3$; $()$ & \texttt{at d = 3} \\
\bottomrule
\end{tabular}
\end{center}
\end{example}

\subsection{Domains, restricting clauses, and box floors}\label{sec:kernel-floors}

A domain clause is a \py!Clause! of \module{archytaszx.semantics.restrict} (\cref{ch:proofs}): a symbol, a kind (\code!multiplicity! or \code!dimension!), a form (\code!eq!, \code!ge!, \code!affine!, \code!shift!, \code!side!, \code!finite!), integers $a$ and $c$, a fresh symbol, and, after validation, the symbol's \code!domain! and the fresh symbol's \code!fresh_domain!. A clause \emph{substitutes} unless it is a side or finite clause or a dimension lower bound.

\begin{definition}[Restricting clauses]\label{def:kernel-restricting}
\py!restricting(t)! is the tuple of clauses of \code!t.domain! that constrain a symbol free in $t$'s sides (every symbol of the sides, \py!_all_symbols!, and every free multiplicity): the clause's symbol is still free, or its fresh symbol is free and its range is narrower than the naturals (a \code!fresh_domain!, a negative constant $c<0$, or a scale $a\le0$). \py!_domain(*ts)! concatenates the restricting clauses of the theorems, repeats dropped. \texttt{\_\allowbreak{}refuse\_\allowbreak{}on(\allowbreak{}where,\allowbreak{} ts,\allowbreak{} S,\allowbreak{} by=None)} raises \texttt{where: a premise holds only under T} (the clause text, or its symbol) when a restricting clause of one of \code!ts! names a symbol of $S$ as its symbol or its fresh symbol and is not discharged by the restriction \code!by!.
\end{definition}

\begin{definition}[Discharged clause]\label{def:kernel-discharged}
A clause $\gamma$ with a fresh symbol $m$ and a fresh domain $\Delta$ is \term{discharged} by a restriction (\py!_discharged!) when the first clause of the restriction on $m$ is an equality $m=c$ with $c\in\Delta$, or is a lower bound $m\ge c$ with no domain of its own while $\Delta$ is an unbounded integer domain ($q=1$, no upper end, no lattice, no excluded value) containing $c$. Every value the restriction allows for $m$ then lies in $\Delta$. A clause without fresh symbol or fresh domain is never discharged.
\end{definition}

\paragraph{Lower bounds of a domain.} \py!_lows(domain)! maps each symbol to the least integer it takes: for each clause, the fresh symbol by its \code!fresh_domain! and, for a side clause, the symbol by its \code!domain!; the bound is $\lceil \mathit{lo}\rceil$, plus one when the bound is strict and integral; the result is the maximum of these bounds and $0$.

\begin{algorithm}[\texttt{\_\allowbreak{}mult\_\allowbreak{}floor}]\label{alg:kernel-mult-floor}
Input: a !-box multiplicity $\mu$ as a SymPy expression and lower bounds. Output: \py!None!, or $(x, v)$.
\begin{enumerate}
\item Sort the symbols of $\mu$ by name. When $\mu$ is a polynomial in them of total degree $t>0$ and a monomial of degree $t$ has a negative coefficient, raise \texttt{the box multiplicity M is negative for large values}.
\item Let each symbol start at its lower bound (default $0$). When $\mu\ge0$ at every point of the grid $\prod_x\{\mathit{lo}_x,\ldots,\mathit{lo}_x+3\}$ ($4$ points per symbol, \py!_PROBE!), return \py!None!.
\item With one symbol $x$: for $v=\mathit{lo}_x+1,\ldots,\mathit{lo}_x+63$ ($<\mathit{lo}_x+$\py!_MAX_FLOOR!, $64$), return $(x,v)$ at the first $v$ with $\mu\ge0$ at $v,\ldots,v+3$.
\item Raise \texttt{the box multiplicity M is negative and no floor of one symbol helps}.
\end{enumerate}
\end{algorithm}

\begin{algorithm}[\texttt{\_\allowbreak{}with\_\allowbreak{}box\_\allowbreak{}floors}]\label{alg:kernel-box-floors}
\begin{enumerate}
\item An \emph{implied} clause is a multiplicity side clause whose source text begins \code!x >= ! (a floor added earlier). Compute the lower bounds of the other clauses, and keep an implied clause only when it raises its symbol's bound above them.
\item For each side and each !-box whose multiplicity has a free symbol and reads no copy index of the side: run \py!_mult_floor! with the bounds of the kept domain and the floors found so far; a result $(x, v)$ raises the floor of $x$ to at least $v$.
\item Append, in name order, one multiplicity side clause per floor: \texttt{Clause(\allowbreak{}x,\allowbreak{} MULTIPLICITY,\allowbreak{} SIDE,\allowbreak{} source="x >= v",\allowbreak{} domain=Domain(\allowbreak{}lo=v,\allowbreak{} q=1))}.
\end{enumerate}
\end{algorithm}

A !-box whose multiplicity is negative below $x=v$ thus restricts its theorem to $x\ge v$; a multiplicity negative under every floor of one symbol makes the theorem unconstructible.

\section{The Inference Rules}\label{sec:kernel-rules}

Each rule is a Python function that checks its inputs and returns a \py!Theorem!. Below, $\Theta$, $\Theta_i$ are premises with $\Theta = (\Gamma;A\vdash L=R)$; $\lfloor\Gamma\rfloor$ is the domain after the box floors of the conclusion's sides (\cref{alg:kernel-box-floors}); $\rho(\Theta)$ is \py!restricting!$(\Theta)$; $A\sqcup A'$ is \py!_merge! (concatenation, repeats dropped, first occurrence kept). Every rule converts a \py!_INDUCTION_ERRORS! exception raised while it rebuilds diagrams into a \py!KernelError!. Every rule function except \py!chain!, \py!restricted!, and \py!_onto! runs inside a span \code!kernel.<name>! of \module{archytaszx.trace} (\py!local_contraction! inside \code!kernel.contraction!).

\subsection{Axiom rules}

\[
\frac{}{\lfloor\emptyset\rfloor;\emptyset\vdash d=d}\;\textsc{Refl}
\qquad
\frac{}{\lfloor\emptyset\rfloor; A_\sigma\vdash d=\sigma(d)}\;\textsc{Rule}(\sigma)
\]
\textsc{Refl} (\py!refl(d)!) has no condition. \textsc{Rule} (\py!rule(d, step)!) takes a recorded \py!RewriteStep! $\sigma$. Conditions: when $\sigma$'s rule is not one of the discard rules (\py!MIXED_RULES!), none of its consumed nodes present in $d$ is a discard (\texttt{the pure rule R consumes a discard}); the step replays on a copy of $d$ by \texttt{\_\allowbreak{}replay\_\allowbreak{}one(\allowbreak{}.\allowbreak{}.\allowbreak{}.\allowbreak{},\allowbreak{} rediscover=True)} of \cref{ch:proofs}: the rule is looked up by name, every consumed node and wire is present, the recorded match is among the matches the pattern rediscovers, \py!apply! succeeds, and the resulting step equals the record. $\sigma(d)$ is the replay's result. $A_\sigma$ is the tuple of equations of \texttt{assumed(\allowbreak{}step.\allowbreak{}dimension\_\allowbreak{}constraints)} (the \texttt{DEFERRED} and \texttt{BOUND} constraints that are not identities, one per equation). The detail is the consumed node ids, sorted, as \code!n3, n5!.

\[
\frac{}{\lfloor\emptyset\rfloor;\emptyset\vdash a=b}\;\textsc{Iso}
\quad\text{if } a\equiv b \text{ or } a\cong_v b.
\]
\textsc{Iso} (\py!iso(a, b)!) refuses with \texttt{the diagrams have different comparison views}.

\[
\frac{}{\lfloor\emptyset\rfloor;\emptyset\vdash a=b}\;\textsc{Algebraic}(\kappa)
\]
\py!algebraic(a, b, cert=None)!. Conditions: neither side has a !-box (\texttt{an algebraic meet of diagrams with bang boxes}); the \emph{values} of a side are its scalar and, per node with a phase, the coefficients $1,\ldots,d-1$ of the phase vector, whose dimension $d$ must be concrete (\texttt{a spider parameter at a symbolic dimension}). The number-field certificate $\kappa$ (\py!numfield.certificate! of all values of both sides when not given, \cref{ch:algebra}) must pass \py!numfield.check!. Each side's \emph{keyed copy} replaces every phase coefficient by the integer token of its field key (\py!numfield.key!; tokens numbered from $2$ in order of first appearance, shared by both sides), with scalar one and no parameters. A \py!NumberFieldError! or \py!NotAlgebraicError! refuses with \texttt{no number-field certificate: E}. The keyed copies must be isomorphic with symmetric legs and the two scalars must have one key (\texttt{the diagrams differ beyond equal algebraic values}). The result is $\textsc{Trans}(\Theta_1, \textsc{Iso}(a', b))$, where $a'$ is $a$ with each node's phase replaced by the phase of its image under the isomorphism and with $b$'s scalar, and $\Theta_1$ is the axiom $a=a'$ of kind \code!algebraic! carrying $\kappa$ as \code!data["field"]!.

\[
\frac{}{\lfloor\emptyset\rfloor;\emptyset\vdash a=b}\;\textsc{Contraction}
\]
\texttt{contraction(\allowbreak{}a,\allowbreak{} b,\allowbreak{} max\_\allowbreak{}steps=64,\allowbreak{} contracted=None)}. The symbolic contractions of both sides with every symbol formal (\texttt{contract\_\allowbreak{}symbolic\_\allowbreak{}pair(\allowbreak{}a,\allowbreak{} b,\allowbreak{} max\_\allowbreak{}steps)}, both doubled when either has a discard; \cref{ch:semantics}) must agree under \py!compare_symbolic! (\cref{ch:proofs}). When \code!contracted! is a decision's tuple $(a_0, b_0, T_a, T_b)$ whose first two entries are these very objects (\py!is!) and both sides are pure, the tensors $T_a$, $T_b$ are used without contracting again. An unsupported shape, a domain error, or an exhausted scalar budget refuses with \code!no symbolic contraction: E!; a mismatch with \texttt{the symbolic contractions differ: R}.

\[
\frac{}{\lfloor\emptyset\rfloor;\emptyset\vdash a=b}\;\textsc{PathSum}
\]
\py!path_sum(a, b)!: \py!compare_path_sums(a, b)! of \cref{ch:semantics} reduces the path sums~\cite{amy2018pathsums} of both to one form (status \code!"equal"!). A \py!PathSumError! refuses with \code!no path sum: E!; another status with \texttt{the path sums are not equal: R}.

\begin{signature}
\texttt{local\_\allowbreak{}contraction(\allowbreak{}a,\allowbreak{} region\_\allowbreak{}a,\allowbreak{} b,\allowbreak{} region\_\allowbreak{}b,\allowbreak{} hypothesis,\allowbreak{} max\_\allowbreak{}steps=64)} --- kind \code!contraction!, data \code!region_a!, \code!region_b!, \code!hypothesis!, \code!parts!.
\end{signature}
\[
\frac{}{\lfloor\emptyset\rfloor;\emptyset\vdash a=b}\;\textsc{LocalContraction}
\]
Conditions: each region, read through cut wires, is the hypothesis up to scalar (\py!_plugged!: \texttt{find\_\allowbreak{}cut\_\allowbreak{}region(\allowbreak{}d,\allowbreak{} region,\allowbreak{} hypothesis)}, else \texttt{the region is not the hypothesis diagram through cut wires}); the part off each region, cut ends open, with scalar $\text{scalar}(d)/\text{scalar}(\text{hypothesis})$ (\py!_copy_side!; an error refuses with \code!the part off the region: E!), and the region's places in $d$'s boundary lists (the held positions and the places, in the cut's lists, of each boundary entry of $d$ on the region) agree between $a$ and $b$ (\texttt{the regions sit in different places of the two diagrams}); the two parts have one interface (\texttt{the parts off the regions have different interfaces}); and \textsc{Contraction} of the two parts succeeds, retried once with $4\cdot$\code!max_steps! when its refusal mentions \code!step budget!. Since both sides are one context around the same diagram, equal contexts give equal sides.

\subsection{Peeling and splitting axioms}

These axioms state equalities that hold by the expansion semantics of !-boxes (\cref{ch:diagrams}); each re-runs the transformation and compares results id for id.

\[
\frac{}{\lfloor\emptyset\rfloor;\emptyset\vdash d=\pi_{n,\Lambda}(d)}\;\textsc{Peel}
\qquad
\frac{}{\lfloor\emptyset\rfloor;\emptyset\vdash \pi_{n,\Lambda_1}(s)=\pi_{n,\Lambda_2}(s)}\;\textsc{Repeel}
\]
\texttt{peel(\allowbreak{}d,\allowbreak{} index,\allowbreak{} last=frozenset(\allowbreak{}))}: $\pi_{n,\Lambda}(d)$ is \texttt{peel\_\allowbreak{}owners(\allowbreak{}d,\allowbreak{} n,\allowbreak{} last=last).\allowbreak{}diagram} with $\Lambda$ the set \code!last!; a \py!BangBoxError! refuses with \code!peel: E!; data \code!index!, \code!last!. \texttt{repeel(\allowbreak{}source,\allowbreak{} index,\allowbreak{} a\_\allowbreak{}last,\allowbreak{} b\_\allowbreak{}last,\allowbreak{} a\_\allowbreak{}first,\allowbreak{} b\_\allowbreak{}first)} (each set empty by default): the two joint peels of \code!source! with the !-boxes of \code!a_last! and of \code!b_last! peeling their last copy and those of \code!a_first! and of \code!b_first! their first copy (\texttt{peel\_\allowbreak{}owners(\allowbreak{}.\allowbreak{}.\allowbreak{}.\allowbreak{},\allowbreak{} last=.\allowbreak{}.\allowbreak{}.\allowbreak{},\allowbreak{} first=.\allowbreak{}.\allowbreak{}.\allowbreak{})}, \cref{ch:diagrams}); when the full joint peel raises \py!BangBoxError!, both are taken as partial joint peels; a second failure refuses with \code!repeel: E!; data \code!source!, \code!index!, and, when a first set is nonempty, \code!first! $=(\text{a\_first},\text{b\_first})$.

\[
\frac{}{\lfloor\emptyset\rfloor;\emptyset\vdash d=d'}\;\textsc{Box}
\]
where \textsc{Box} is one of \textsc{IdentitySlide}, \textsc{BoxSplit}, \textsc{LegMove}, \textsc{PowerSplit}, \textsc{WFusion}, \textsc{AlignMerge}, and \textsc{ChainSplit}:
\begin{itemize}
\item \texttt{identity\_\allowbreak{}slide(\allowbreak{}d,\allowbreak{} slid,\allowbreak{} copy)}: some element of \py!identity_slides(d, copy)! is identical to $d'$; else \texttt{identity\_\allowbreak{}slide: the diagram is no identity slide of the other}.
\item \py!box_split(d, split, copy)!: some diagram of \py!boxed_splits(d, copy)! is identical to $d'$; else \texttt{box\_\allowbreak{}split: the diagram is no boxed spider split of the other}.
\item \py!leg_move(d, moved, copy)!: some element of \py!leg_moves(d, copy)! is identical to $d'$; else \texttt{leg\_\allowbreak{}move: the diagram is no leg move of the other}.
\item \py!power_split(d, split)!: for some !-box of $d$ (in id order) with a power multiplicity $\beta^{e+1}$, \py!split_power(d, box)! of \cref{ch:diagrams} (the !-box as $\beta$ !-boxes of $\beta^e$ copies, the $c$-th reading the copy index as $\text{index}+c\beta^e$) is identical to $d'$; \py!BangBoxError!s skip a !-box; else \texttt{power\_\allowbreak{}split: the diagram is no power box split of the other}; data \code!box!.
\item \py!w_fusion(d)!: $d'$ is \py!w_fused(d)! (\cref{alg:kernel-w-fused}); an error refuses with \code!w_fusion raised T: E!.
\item \py!align_merge(d, merged)! (\textsc{AlignMerge}): for some alignment root of $d$ (sorted; a root no longer present skipped) whose group passes \py!_align_group!, \py!align_merged(d, root)! of \cref{ch:diagrams} (copy $j$ of each member, wired to copy $j$ of the others, is copy $j$ of one !-box over the union of their scopes) is identical to $d'$; a failing group or a \py!BangBoxError! skips a root; else \texttt{align\_\allowbreak{}merge: the diagram is no aligned box merge of the other}; data \code!box! (the root).
\item \py!chain_split(d, split)! (\textsc{ChainSplit}): for some alignment root of $d'$ (sorted) present in both diagrams whose group in $d'$ passes \py!_chain_group!, \py!chain_merged(split, root)! of \cref{ch:diagrams} with the root's links set to $d$'s root's links is identical to $d$, where $d$'s root has no repeated link and its links, sorted by their text, are those of all members; else \texttt{chain\_\allowbreak{}split: the diagram is no chain box split of the other}; data \code!box!. A chain !-box over independent chains equals the aligned chain !-boxes, one per chain.
\end{itemize}

\paragraph{Alignment groups that may merge.} \py!_align_group(d, root)! returns the alignment group of \code!root! (\py!aligned_members!) when it has at least two members, starts at \code!root!, and \code!root! is aligned to nothing (\texttt{align\_\allowbreak{}merge: no alignment group at the box}); every member is top level and node scoped, with no port scope, links, cyclic chain, tree, feed, pair parent, or groups, and with the root's sheets flag, index, and multiplicity (expanded difference zero) (\texttt{align\_\allowbreak{}merge: the aligned boxes differ in kind,\allowbreak{} count or index}); at most one member carries a copy scalar (\texttt{align\_\allowbreak{}merge: two aligned boxes carry copy scalars}); the scopes are pairwise disjoint (\texttt{align\_\allowbreak{}merge: the aligned scopes overlap}); no other !-box lies in a member or aligns to a member other than the root (\texttt{align\_\allowbreak{}merge: a box lies in or aligns to a member}), and none meets the scopes by node or port scope (\texttt{align\_\allowbreak{}merge: another box meets the aligned scopes}). \py!_chain_group(d, root)! makes the same checks under the prefix \code!chain_split:!, except that links are allowed: every link's two ports lie in its own !-box's scope (\texttt{chain\_\allowbreak{}split: a link leaves its box}), at least one member has links (\texttt{chain\_\allowbreak{}split: no aligned box is a chain}), and no other !-box may pair over a member either (\texttt{chain\_\allowbreak{}split: a box lies in,\allowbreak{} aligns to or pairs over a member}).

\subsection{Structural rules}

\[
\frac{\Gamma;A\vdash L=R}{\lfloor\rho(\Theta)\rfloor;A\vdash R=L}\;\textsc{Symm}
\qquad
\frac{\Gamma_1;A_1\vdash L=M \qquad \Gamma_2;A_2\vdash M'=R}{\lfloor\rho(\Theta_1)\cdot\rho(\Theta_2)\rfloor;A_1\sqcup A_2\vdash L=R}\;\textsc{Trans}\quad\text{if } M\equiv M'.
\]
\textsc{Trans} (\py!trans(a, b)!) refuses with \texttt{trans: the middle diagrams differ}; the middle diagrams must be identical id for id, so every proof inserts \textsc{Iso} steps where a search met its goal up to isomorphism. \py!chain(start, steps)! is the derived rule $\textsc{Rule}(\sigma_1)\mathbin{;}\cdots\mathbin{;}\textsc{Rule}(\sigma_m)$, each step applied to the previous right side, and \textsc{Refl}$(\textit{start})$ when there is no step. \py!_onto(t, left, right)! re-states $t$ over \code!left! and \code!right! by $\textsc{Iso}(\textit{left},L)\mathbin{;}t\mathbin{;}\textsc{Iso}(R,\textit{right})$, each \textsc{Iso} only when the sides are not identical.

\[
\frac{\Gamma;A\vdash L=R}{\lfloor\rho(\Theta)\rfloor;A\vdash C=C[R/L]}\;\textsc{Congr}(\chi)
\]
\py!congr(t, context, cut)!, data \code!cut!. With $\chi$ a \py!CutRegion! of the context $C$ over the region $X$: $X$ lies in $C$ (\texttt{congr: the region is not part of the context}); $\chi$'s cuts are exactly the cut wires of $X$ (\texttt{congr: the cut wires are not the region's}); per port direction, $\chi$'s list holds exactly the boundary entries of $C$ on $X$ of that direction and the region ends of that direction (\texttt{congr: the region's boundary is not its slots and cut ends}); $L$'s scalar is not zero (\texttt{congr: the theorem's left side has scalar zero}); the bare cut extraction of $X$ is isomorphic to the bare $L$ (\texttt{congr: the region is not the theorem's left side}). $C[R/L]$ is \texttt{replace\_\allowbreak{}cut\_\allowbreak{}region(\allowbreak{}C,\allowbreak{} cut,\allowbreak{} L,\allowbreak{} R,\allowbreak{} contiguous=False)}, which multiplies the scalar by $\text{scalar}(R)/\text{scalar}(L)$.

\[
\frac{\Gamma;A\vdash L=R}{\lfloor\rho(\Theta)\rfloor;A\vdash L'=R'}\;\textsc{Regroup}
\]
\py!regroup(t, lhs, rhs)!. Conditions: each new side has the !-box ids of the old one (\code!regroup: the boxes differ!); the old side with each !-box's groups set to the new side's groups is identical to the new side (\texttt{regroup: the sides differ beyond box groups}); and the slot patterns (\py!_slot_pattern!) of $L$ and $R$ agree, those of $L'$ and $R'$ agree, and none is \py!None! (\texttt{regroup: the sides lay out their slots differently}). The \emph{slot pattern} gives, per boundary list and per entry, its replication count (the expanded multiplicity of the one node-scoped !-box over it, $1$ off every !-box) and the position of the first entry of its group block; it is \py!None! when a !-box is port scoped, or an entry lies in more than one !-box or in a nested one. One permutation of the expanded boundary then carries both sides.

\[
\frac{\Gamma;A\vdash L=R}{\lfloor\rho(\Theta)\rfloor;A\vdash L[x:=e]=R[x:=e]}\;\textsc{Instance}
\]
\py!instance(t, symbol, value)!, detail \code!x = e!, data \code!symbol!, \code!value!. Conditions: $e$ is a \py!Mult! (\texttt{instance: V is not a multiplicity}); $x$ is a free multiplicity of a side (\texttt{instance: 'x' is not a free multiplicity}); neither $x$ nor a symbol of $e$ is a copy index (\texttt{instance: a symbol is a copy index}); no restricting clause names $x$ (\py!_refuse_on("instance", ...)!). Both sides are rebuilt by \py!rename_multiplicity!; an error refuses with \code!instance raised T: E!.

\subsection{Induction}

\begin{signature}
\texttt{induction(\allowbreak{}start,\allowbreak{} goal,\allowbreak{} index,\allowbreak{} step\_\allowbreak{}symbol,\allowbreak{} base,\allowbreak{} side,\allowbreak{} expose,\allowbreak{} region,\allowbreak{} step,\allowbreak{} *,\allowbreak{} cut=None,\allowbreak{} tree=None,\allowbreak{} offset=None,\allowbreak{} parameter=None) -> Theorem}
\end{signature}

Let $n$ be \code!index!, $k$ \code!step_symbol!, \textit{own} the side named by \code!side! (\texttt{START}: \textit{start}; \texttt{GOAL}: \textit{goal}) and \textit{other} the other one. Without an offset, $P=$ \texttt{induction\_\allowbreak{}successor(\allowbreak{}own,\allowbreak{} n,\allowbreak{} k)}, $P'=$ \texttt{induction\_\allowbreak{}successor(\allowbreak{}other,\allowbreak{} n,\allowbreak{} k)}, $H=$ \texttt{induction\_\allowbreak{}hypothesis(\allowbreak{}own,\allowbreak{} n,\allowbreak{} k)}, and $H'=$ \texttt{induction\_\allowbreak{}hypothesis(\allowbreak{}other,\allowbreak{} n,\allowbreak{} k)}.
\[
\frac{\Gamma_0;A_0\vdash \mathit{start}|_{n=0}=\mathit{goal}|_{n=0}
\qquad \Gamma_1;A_1\vdash P=E
\qquad \Gamma_2;A_2\vdash X=P'}
{\lfloor\rho(\Theta_0)\cdot\rho(\Theta_1)\cdot\rho(\Theta_2)\rfloor;A_0\sqcup A_1\sqcup A_2\vdash \mathit{start}=\mathit{goal}}\;\textsc{Induction}
\]
Conditions, checked in order:
\begin{enumerate}
\item $n$ is a free multiplicity of a side and $k$ is not (\texttt{'n' is not a free multiplicity or 'k' is not fresh}).
\item No restricting clause of the three premises names $n$, $k$, or the offset (\py!_refuse_on("induction", ...)!).
\item The base premise runs from $\mathit{start}|_{n=0}$ to $\mathit{goal}|_{n=0}$, both id for id (\texttt{the base theorem does not start at the start side at 0}, \texttt{the base theorem does not end at the goal side at 0}).
\item With an offset (below), the successors are built from the shifted sides. Both successors exist (\texttt{a successor has no single box to peel}).
\item The exposing premise starts at $P$ (\texttt{the exposing theorem does not start at the peeled successor}); $E$ is its right side.
\item The window. \emph{Cut} ($\chi$ given): $\chi$'s region is \code!region! and lies in $E$ (\texttt{the cut region is not the hypothesis region}), and its bare cut extraction is isomorphic to the bare $H$ (\texttt{the cut region is not the hypothesis diagram}); $X=$ \texttt{replace\_\allowbreak{}cut\_\allowbreak{}region(\allowbreak{}E,\allowbreak{} cut,\allowbreak{} H,\allowbreak{} H')}. \emph{Tree} (a tree !-box id $t$): \py!unfed(E, t)! has exactly the nodes of \code!region! and is isomorphic to $H$ (\texttt{the unfed exposed diagram is not the hypothesis diagram}); $X=$ \py!fed_replacement(E, t, H, H')!. \emph{Region} (neither): \code!region! lies in $E$ and no wire leaves it (\texttt{the hypothesis region is not a closed part of the exposed diagram}), and its bare extraction with !-boxes is isomorphic to the bare $H$ (\texttt{the hypothesis region is not the hypothesis diagram}); $X=$ \texttt{replace\_\allowbreak{}region(\allowbreak{}E,\allowbreak{} region,\allowbreak{} H,\allowbreak{} H')}.
\item The step premise runs from $X$ to $P'$, id for id (\texttt{the step theorem does not run from the replaced to the other side}).
\end{enumerate}
An \py!_INDUCTION_ERRORS! exception in steps 3--7 refuses with \texttt{rebuilding the induction raised T: E}. The data record \code!index!, \code!step!, \code!side!, \code!peeled! ($P$), \code!other_peeled! ($P'$), \code!own_k! ($H$), \code!other_k! ($H'$), \code!replaced! ($X$), \code!region!, \code!exposed! ($E$), \code!window! (\code!"cut"!, \code!"tree"!, or \code!"region"!), and \code!offset!.

\paragraph{Soundness.} Induct on $n$. At $n=0$ the base premise gives the claim. At $n=k+1$, by the hypothesis $\mathit{own}(k)=\mathit{other}(k)$ in every context: $\mathit{own}(k{+}1)=P=E=C[H]=C[H']=X=P'=\mathit{other}(k{+}1)$, where the first and last equalities are peeling, $C$ is the context of the window, and the scalar correction of the replacement accounts for the scalars of $H$ and $H'$ (\cref{fig:kernel-induction}).

\paragraph{Offset induction.} With \code!offset=s!, the inference kind is \code!induction_offset! and the claim inducted on is $\forall s.\ \mathit{start}^{+s}=\mathit{goal}^{+s}$, where $D^{+s}$ reads the copy index of the !-box carrying $n$ as $\text{index}+s$ (\py!induction_offset!). The kernel first checks (\py!_offset_shift!) that $s$ is fresh, named neither among the free symbols of either side, nor as $n$, $k$, nor as any copy index (\code!the offset 's' is not fresh!), and that each side at $n=0$ is unchanged by the shift (\texttt{a side at 0 depends on the offset}). The successors are then built from $\mathit{own}^{+s}$ and $\mathit{other}^{+s}$, and the hypothesis diagrams $H$, $H'$ read their copy index as $\text{index}+s+1$; the base premise is unchanged. The hypothesis at $s+1$ covers the residual of a chain that peels its first copy.

\paragraph{Parametric induction.} With \code!parameter=(a, v)!, the inference kind gains the suffix \code!_param! (\code!induction_param!, or \code!induction_offset_param! with an offset), the data record \code!parameter!, and the claim inducted on is $\forall a\in\mathbb{R}.\ \mathit{start}=\mathit{goal}$; the hypothesis diagrams $H$, $H'$ (after any offset) are taken at $a=v$ by \py!substitute_parameter!. Before anything else, \py!_parameter_move! checks the pair through \texttt{\_\allowbreak{}parameter\_\allowbreak{}checked(\allowbreak{}start,\allowbreak{} goal,\allowbreak{} a,\allowbreak{} v)}: $a$ is a phase parameter (\cref{def:kernel-phase-parameter}) of at least one side (\texttt{'a' is not a phase parameter of the claim}) and is read outside node phases in neither (\texttt{'a' is read outside node phases}), and $v$ is real and free of copy indices and of the phase indices (\texttt{V is not a real value free of the copy and phase indices}); and neither $a$ nor a symbol of $v$ is $n$, $k$, or the offset (\texttt{the parameter 'a' or its value reads the induction symbols}). $a$ joins the symbols on which no restricting premise clause may lie. Soundness: the hypothesis holds for every real $a$, in particular at $a=v$.

\[
\frac{\Gamma;A\vdash L=R}{\lfloor\rho(\Theta)\rfloor;A\vdash L[a:=v]=R[a:=v]}\;\textsc{PhaseInstance}
\]
\texttt{phase\_\allowbreak{}instance(\allowbreak{}t,\allowbreak{} symbol,\allowbreak{} value)}, detail \code!a = v!, data \code!symbol!, \code!value!: the checks of \py!_parameter_checked! on $t$'s sides; no restricting clause names $a$; both sides by \texttt{substitute\_\allowbreak{}parameter(\allowbreak{}.\allowbreak{}.\allowbreak{}.\allowbreak{},\allowbreak{} normalize=False)} (an error refuses with \code!phase_instance raised T: E!).

\subsection{Periodic closure}

\[
\frac{\emptyset;A_0\vdash \mathit{start}|_{n=0}=\mathit{goal}|_{n=0}\ \ \cdots\ \ \emptyset;A_{p-1}\vdash \mathit{start}|_{n=p-1}=\mathit{goal}|_{n=p-1}}{\lfloor\emptyset\rfloor;A_0\sqcup\cdots\sqcup A_{p-1}\vdash \mathit{start}=\mathit{goal}}\;\textsc{Periodic}
\]
\texttt{periodic(\allowbreak{}start,\allowbreak{} goal,\allowbreak{} index,\allowbreak{} bases)}, data \code!index!, \code!step! ($p$), \code!ratio!. Conditions: $n$ is a free multiplicity and there is at least one base (\texttt{'n' is not a free multiplicity or there is no base theorem}); no base has a domain (\texttt{a base theorem holds under a restriction}); base $r$ runs between the sides at $n=r$ (\py!_instantiate_index!), id for id (\texttt{base r does not start at the start side at r}, \texttt{base r does not end at the goal side at r}); and both sides scale by one constant from $n$ to $n+p$: \py!scaling(side, n, p)! of \module{archytaszx.semantics.periodic} (\cref{ch:proofs}; $n$ read only in the scalar and in the multiplicity of at most one top-level chain whose copy transfer matrix has a scalar power with period dividing $p$) gives $a_{\mathit{start}}$ and $a_{\mathit{goal}}$ with $a_{\mathit{start}}-a_{\mathit{goal}}$ simplifying to $0$ (\texttt{periodic: the sides scale by A and B}). A \py!PeriodicError! refuses with \code!periodic: E!. Soundness: strong induction with step $p$, $\mathit{start}(n+p)=a\,\mathit{start}(n)=a\,\mathit{goal}(n)=\mathit{goal}(n+p)$.

\subsection{Substitution and restriction}

\[
\frac{\Gamma;A\vdash L=R}{\lfloor\Gamma\cdot\gamma\rfloor;A\vdash L\gamma=R\gamma}\;\textsc{Subst}(\gamma)
\]
\texttt{subst(\allowbreak{}t,\allowbreak{} clause,\allowbreak{} lhs=None,\allowbreak{} rhs=None)}, detail the clause text, data \code!clause!, origin \code!t.origin! or $(L,R)$. Conditions: $\gamma$ substitutes (\code!subst: T substitutes nothing!); its fresh symbol is not a symbol of the sides (\texttt{subst: m is not fresh for the theorem}); no restricting clause names its symbol unless $\gamma$ alone discharges it (\cref{def:kernel-discharged}); for a multiplicity clause with a fresh symbol and no fresh domain, $a\ge1$, and $a=1$ when $c<0$ (\code!subst: T leaves the naturals!). A shift $n=m-c'$ ($c=-c'<0$) is kept with the fresh domain $m\ge c'$ and substituted into phases and port dimensions too (\code!everywhere=True!). Each side is rebuilt by \py!apply_clause! (\cref{ch:proofs}); a stated side must be identical to the rebuilt one (\texttt{subst: a stated side is not the theorem's side at T}), and the conclusion's sides are the stated ones when given. The conclusion keeps the whole premise domain, not only its restricting clauses, followed by $\gamma$.

\paragraph{Restricted reading.} \texttt{restricted(\allowbreak{}t,\allowbreak{} left,\allowbreak{} right,\allowbreak{} restriction,\allowbreak{} used=(\allowbreak{}))} is no new inference: it returns a theorem with $t$'s sides and inference, read as the restricted claim of $\mathit{left}=\mathit{right}$. Each side of $t$ must be identical or view-isomorphic to the corresponding side of \texttt{subst\_\allowbreak{}pair(\allowbreak{}left,\allowbreak{} right,\allowbreak{} restriction)} (\texttt{the theorem is not about the restricted pair}; a rebuilding error refuses with \code!restricting raised T: E!). The domain is the restriction's clauses followed by $t$'s other restricting clauses; the assumptions gain the floors \code!c <= d! of each dimension clause $d\ge c$ whose symbol is in \code!used! (\py!_floors_used!); the origin is $(\mathit{left},\mathit{right})$.

\[
\frac{\Gamma;A\vdash L\gamma=R\gamma}{\lfloor\gamma\cdot\rho\rfloor;A\vdash L=R}\;\textsc{Shift}(\gamma)
\qquad
\frac{\Gamma;A\vdash L\gamma_1\cdots\gamma_r=R\gamma_1\cdots\gamma_r}{\lfloor\gamma_1\cdots\gamma_r\cdot\rho\rfloor;A\vdash L=R}\;\textsc{Lift}(\gamma_1,\ldots,\gamma_r)
\]
Both state the original pair under the clauses from a theorem about the substituted pair. The premise is first moved onto the rebuilt sides by \py!_onto!; $\rho$ is the restricting clauses of that moved premise other than the clauses themselves; the origin is $(L,R)$; the inference kind of both is \code!shift!.
\begin{itemize}
\item \texttt{shift(\allowbreak{}t,\allowbreak{} left,\allowbreak{} right,\allowbreak{} clause)}: a multiplicity equality clause is passed to \py!lift! with that one clause. Otherwise the clause must be a multiplicity clause with a fresh symbol and not $n=m$ ($a=1$, $c=0$) (\texttt{shift: T is not n = a*m + c with a > 1 or c >= 1}); $n$ must be a free multiplicity of the sides and the fresh symbol not a symbol of them (\texttt{'n' is not a free multiplicity or 'm' is not fresh}). Each side is rebuilt by \py!apply_clause! (\texttt{rebuilding the successor raised T: E}), and no restricting clause of the moved premise names $n$ or $m$. Data: \code!index!, \code!step! ($m$), \code!count! ($c$), \code!scale! ($a$), \code!peeled!, \code!other_peeled!.
\item \texttt{lift(\allowbreak{}t,\allowbreak{} left,\allowbreak{} right,\allowbreak{} clauses)}: every clause is a multiplicity clause (\texttt{lift: T is not a multiplicity clause}) on a free multiplicity (\texttt{'n' is not a free multiplicity}); a non-equality clause has a fresh symbol not among the sides' symbols (\code!lift: T has no fresh symbol!); there is at least one clause and one per symbol (\texttt{lift: one clause per multiplicity}). The clauses are applied in order (\texttt{rebuilding the instance raised T: E}); no restricting clause names a symbol of a clause. Data as \textsc{Shift} with \code!scale!~$0$, \code!step! \py!None!, \code!fixed! (the pairs $(n,c)$ of the equality clauses), and \code!lift! (per clause $(n, 0, c, \text{None})$ for an equality and $(n, a, c, m)$ otherwise).
\end{itemize}

\begin{table}[H]
\centering\small
\begin{tabularx}{\linewidth}{@{}>{\raggedright\arraybackslash}p{0.3\linewidth}>{\raggedright\arraybackslash}p{0.24\linewidth}c>{\raggedright\arraybackslash}X@{}}
\toprule
Kind & Built by & Premises & Data keys \\
\midrule
\texttt{refl} & \texttt{refl} & 0 & --- \\
\texttt{rule} & \texttt{rule} & 0 & \texttt{step} (also \texttt{rule}, \texttt{detail}) \\
\texttt{iso} & \texttt{iso} & 0 & --- \\
\texttt{algebraic} & \texttt{algebraic} & 0 & \texttt{field} \\
\texttt{contraction} & \texttt{contraction}, \texttt{local\_contraction} & 0 & ---; \texttt{region\_a}, \texttt{region\_b}, \texttt{hypothesis}, \texttt{parts} \\
\texttt{path\_sum} & \texttt{path\_sum} & 0 & --- \\
\texttt{peel}, \texttt{repeel} & \texttt{peel}, \texttt{repeel} & 0 & \texttt{index}, \texttt{last}; \texttt{source}, \texttt{index} \\
\texttt{identity\_slide}, \texttt{box\_split}, \texttt{leg\_move}, \texttt{w\_fusion} & same names & 0 & --- \\
\texttt{power\_split}, \texttt{align\_merge}, \texttt{chain\_split} & same names & 0 & \texttt{box} \\
\texttt{phase\_instance} & \texttt{phase\_instance} & 1 & \texttt{symbol}, \texttt{value} \\
\texttt{symm}, \texttt{trans} & \texttt{symm}, \texttt{trans} & 1, 2 & --- \\
\texttt{congr} & \texttt{congr} & 1 & \texttt{cut} \\
\texttt{regroup} & \texttt{regroup} & 1 & --- \\
\texttt{instance} & \texttt{instance} & 1 & \texttt{symbol}, \texttt{value} \\
\texttt{induction}, \texttt{induction\_offset}, \texttt{induction\_param}, \texttt{induction\_offset\_param} & \texttt{induction} & 3 & see the induction rule; \texttt{offset}, \texttt{parameter} \\
\texttt{periodic} & \texttt{periodic} & $p$ & \texttt{index}, \texttt{step}, \texttt{ratio} \\
\texttt{subst} & \texttt{subst} & 1 & \texttt{clause} \\
\texttt{shift} & \texttt{shift}, \texttt{lift} & 1 & see the shift and lift rules \\
\bottomrule
\end{tabularx}
\caption{Inference kinds, their constructors, premise counts, and data.}\label{tab:kernel-inferences}
\end{table}

\section{From Certificates and Decisions to Theorems}\label{sec:kernel-check}

\begin{signature}
\texttt{@dataclass(\allowbreak{}frozen=True,\allowbreak{} eq=False) class RestrictedProof}: \py!proof: object!, \py!restriction: Restriction!, \py!left: Diagram!, \py!right: Diagram!, \py!family: bool!, \texttt{floors\_\allowbreak{}used: tuple[str,\allowbreak{} .\allowbreak{}.\allowbreak{}.\allowbreak{}] = (\allowbreak{})}.

A proof of $\mathit{left}=\mathit{right}$ under \code!restriction!: with \code!family!, \code!proof! proves the unrestricted claim and is specialized by \textsc{Subst}; otherwise it proves the restricted pair. \code!floors_used! names the dimension floors a normalizer proof assumed.
\end{signature}

\begin{algorithm}[\texttt{check}]\label{alg:kernel-check}
Runs inside the span \code!proof.kernel_check!.
\begin{enumerate}
\item A \py!Theorem! is returned as it is.
\item A \py!RestrictedProof!: check its proof, giving $\Theta$. Without \code!family!, return \texttt{restricted(\allowbreak{}inner,\allowbreak{} left,\allowbreak{} right,\allowbreak{} restriction,\allowbreak{} floors\_\allowbreak{}used)} for $\Theta$ the inner theorem. With \code!family!, every restricting clause of $\Theta$ must be discharged by the restriction (\cref{def:kernel-discharged}; else \texttt{a family proof holds only under a restriction}); apply \textsc{Subst} once per substituting clause of the restriction, in order, skipping a multiplicity equality $n=v$ on a free multiplicity of $\Theta$'s sides (\py!_pinned_box!), and return a theorem with the final sides, inference, and assumptions, the domain $\rho(\Theta)$ followed by the restriction's clauses, and the origin $(\Theta.\mathrm{lhs},\Theta.\mathrm{rhs})$.
\item A step-sequence \py!Certificate! (\cref{ch:proofs}): an induction certificate is refused (\texttt{an induction certificate carries no re-checkable step}). Otherwise $\Theta\gets$ \py!chain(initial, steps)!, whose right side must be identical to the recorded final diagram (\texttt{the steps do not reach the recorded final diagram}).
\item A \py!ProofCertificate!: the forward half must begin at the start and the backward half at the goal, id for id (\texttt{a proof half does not begin at its end of the goal}). Check both halves, $\Phi$ and $\Psi$. The \emph{meet} is \textsc{Iso}$(\Phi.\mathrm{rhs},\Psi.\mathrm{rhs})$, else \textsc{Algebraic} of the two (whose refusal is discarded in favor of the \textsc{Iso} refusal). Return $\Phi\mathbin{;}\text{meet}\mathbin{;}\textsc{Symm}(\Psi)$.
\item An \py!InductionProof! (\cref{ch:proofs}): return \textsc{Induction} over its start, goal, index, step symbol, side, and region, with the region window and the premises $\mathrm{check}(\mathit{base})$, $\mathrm{check}(\mathit{expose})$, and $\mathrm{check}(\mathit{step})$.
\item Anything else: \code!cannot check a T!.
\end{enumerate}
\end{algorithm}

\texttt{instantiated(\allowbreak{}theorem) -> (\allowbreak{}Diagram,\allowbreak{} Diagram)} returns the theorem's own sides when it has no origin; otherwise it applies every substituting clause of the domain, in order, to the origin pair by \py!apply_clause! (with \code!everywhere! when $c<0$), as \textsc{Subst} applies them.

\begin{algorithm}[\texttt{from\_\allowbreak{}decision}]\label{alg:kernel-from-decision}
Reads a \py!Decision! of \py!decide_equal! (\cref{ch:proofs}) as $\mathit{left}=\mathit{right}$; runs inside the span \code!proof.build!. Anything but a \py!Decision! refuses with \code!cannot read a T!.
\begin{enumerate}
\item \texttt{DOUBLED}: the decision must carry its doubled pair and an inner method other than \texttt{DOUBLED} (\texttt{a doubled verdict without its doubled pair and settling rung}). When \textit{left} and \textit{right} are not the pair itself, each must be identical to its pair member or double to it (\py!double!, \cref{ch:mixed}) (\texttt{the doubled pair is not the doubling of left and right}). Recurse on the decision with the inner method and the doubled pair.
\item \texttt{INDUCTION} with an \py!induction_proof!: $\Theta\gets$ \py!check(induction_proof)!; keep it in the induction memory (below); return $\textsc{Iso}(\mathit{left},\Theta.\mathrm{lhs})\mathbin{;}\Theta\mathbin{;}\textsc{Iso}(\Theta.\mathrm{rhs},\mathit{right})$.
\item \texttt{NORMAL\_FORM} with exactly two certificates: check both as $\Phi$ (from $\mathit{left}$'s side) and $\Psi$; return $\textsc{Iso}(\mathit{left},\Phi.\mathrm{lhs})\mathbin{;}\Phi\mathbin{;}\textsc{Iso}(\Phi.\mathrm{rhs},\Psi.\mathrm{rhs})\mathbin{;}\textsc{Symm}(\Psi)\mathbin{;}\textsc{Iso}(\Psi.\mathrm{lhs},\mathit{right})$.
\item \texttt{SATURATION}: start at \textsc{Refl}$(\mathit{left})$; for each certificate (an edge of the e-graph path, \cref{ch:rewriting}), check it as $\Theta$; when the current right side is view-isomorphic to $\Theta.\mathrm{lhs}$, append $\textsc{Iso}\mathbin{;}\Theta$; when it is view-isomorphic to $\Theta.\mathrm{rhs}$, append $\textsc{Iso}\mathbin{;}\textsc{Symm}(\Theta)$; else refuse with \texttt{a saturation edge does not continue the path}. Close with \textsc{Iso} onto \textit{right}.
\item \texttt{SYMBOLIC\_CONTRACTION}: \texttt{contraction(\allowbreak{}left,\allowbreak{} right,\allowbreak{} contracted=decision.\allowbreak{}contracted)}.
\item \texttt{PATH\_SUM}: \py!path_sum(left, right)!.
\item Any other method, and \texttt{INDUCTION} without a proof: \texttt{a M verdict carries no re-checkable proof} ($M$ the method's value).
\end{enumerate}
\end{algorithm}

\subsection{The induction memory and the induction rung of the export}\label{sec:kernel-memory}

The kernel keeps a module-level list \py!_INDUCTED! of triples $(x, y, \Theta)$, $\Theta$ a theorem or \py!None!: the induction theorems of decisions read by \py!from_decision! and the theorems (or failures) built by \py!_induction_rung!. \texttt{\_\allowbreak{}remember\_\allowbreak{}induction(\allowbreak{}proof,\allowbreak{} x=None,\allowbreak{} y=None)} appends one triple (its sides default to the theorem's) and keeps the latest $16$ (\py!INDUCTED_KEPT!). \py!_recalled_induction(x, y)! searches from the latest for a triple with both diagrams identical to $x$ and $y$ and returns \py!(True, theorem_or_None)!, else \py!(False, None)!. The memory is not thread safe and persists for the life of the process.

\begin{algorithm}[\texttt{\_\allowbreak{}induction\_\allowbreak{}rung}]\label{alg:kernel-induction-rung}
\begin{enumerate}
\item A remembered triple for $(x,y)$ returns its theorem (or \py!None!).
\item Otherwise run the induction rung of \py!decide_equal! directly (\py!_decide_by_induction!, \cref{ch:proofs}) with the first $4$ (\py!INDUCTED_SAMPLES!) points of \py!sample_grid(x, y)! as the evaluated samples, \code!samples=None!, \code!rules=None!, the default rewrite guard, \code!max_samples=24!, \code!use_symbolic=True!, \code!max_depth=2!, the default tolerance and element budget, \code!use_saturation=True!, and the decision saturation limits. Any exception gives \py!None!. The result is the induction theorem when the rung returns \texttt{EQUAL} with an \py!induction_proof! that is a \py!Theorem! with sides identical to $x$ and $y$; else \py!None!.
\item Remember the result under $(x,y)$.
\end{enumerate}
\end{algorithm}

\py!_offset_index(x, y)! is $(n, c)$ when $n$ is the only free multiplicity of the pair and every !-box multiplicity reading it is $n+c$ for one integer $c\ge1$; else \py!None!. \py!_contraction_induction(x, y)! proves $x=y$ by \py!_induction_rung! with no shift when \py!_offset_index! is \py!None!; otherwise it takes the first of \code!m0!, \code!m1!, \ldots{} not among the pair's symbols as $m$, substitutes the shift clause $n = m - c$ into both diagrams (any exception gives \py!None!), requires $m$ to be the only free multiplicity of the result, and proves the shifted pair by \py!_induction_rung!, returning the theorem and the shift $(n, m, c)$.

\section{Refutations and Symbolic-Sum Proofs}\label{sec:kernel-refute}

Two records stand beside \py!Theorem! and are exported by \py!lean! as well.

\begin{signature}
\texttt{@dataclass(\allowbreak{}frozen=True,\allowbreak{} eq=False) class Refutation}: \py!lhs!, \py!rhs!, \py!point: Mapping[str, object]!, \py!source: str = ""! --- the claim $\mathit{lhs}=\mathit{rhs}$ for every parameter value, refuted at \code!point!; \code!source! is its Lean export once written.

\texttt{refutation(\allowbreak{}decision,\allowbreak{} left,\allowbreak{} right) -> Refutation} --- a \py!Decision! with verdict \texttt{UNEQUAL} and a counterexample gives the refutation at that counterexample; a non-decision refuses with \code!cannot read a T!, a decision without counterexample with \texttt{the decision carries no counterexample}.

\texttt{@dataclass(\allowbreak{}frozen=True,\allowbreak{} eq=False) class Symbolic}: \py!lhs!, \py!rhs!, \py!source: str!, \texttt{domain: tuple[Clause,\allowbreak{} .\allowbreak{}.\allowbreak{}.\allowbreak{}] = (\allowbreak{})}, \texttt{assumptions: tuple[str,\allowbreak{} .\allowbreak{}.\allowbreak{}.\allowbreak{}] = (\allowbreak{})}, with \py!variables! and \py!statement()! as for \py!Theorem! --- $\mathit{lhs}=\mathit{rhs}$ for every value of the dimensions and phases, proved in Lean by the verified normalizer of symbolic sums (\cref{ch:lean-sym}); \code!source! is that Lean proof.

\texttt{symbolic(\allowbreak{}left,\allowbreak{} right,\allowbreak{} describe=None) -> Symbolic} --- \texttt{lean\_\allowbreak{}sym.\allowbreak{}export\_\allowbreak{}symbolic(\allowbreak{}left,\allowbreak{} right,\allowbreak{} describe=describe)}; a \py!NetUnsupported! or \py!ScalarUnsupported! refuses with \code!no symbolic proof: E!.
\end{signature}

\begin{algorithm}[\texttt{restricted\_\allowbreak{}symbolic}]\label{alg:kernel-restricted-symbolic}
The inner proof must be a \py!Symbolic! (\texttt{the restricted proof is not a normaliser proof}).
\begin{enumerate}
\item With \code!family!: put the header note \texttt{restricted to T: the family theorem,\allowbreak{} then its corollary at the clauses} on the source (\py!lean_restrict.header!, \cref{ch:lean-export}) and append the corollary at the clauses (\texttt{corollary(\allowbreak{}source,\allowbreak{} clauses,\allowbreak{} "main",\allowbreak{} \_\allowbreak{}var)}; a \py!RestrictUnsupported! refuses with \code!no corollary: E!). The domain is the restriction's clauses.
\item Otherwise the assumptions are the floors used. Try the shifted proof over the original diagrams (\py!_symbolic_shift!); with one, the header is the pair note followed by \code!; ! and the lifted note (\cref{sec:kernel-lean-restricted}). Without one, the header is the pair note followed by the \texttt{UNLINKED} note when a clause substitutes.
\end{enumerate}
\end{algorithm}

\begin{signature}
\texttt{@dataclass(\allowbreak{}frozen=True,\allowbreak{} eq=False) class FiniteProof}: \texttt{instances: tuple[tuple[Mapping[str,\allowbreak{} Fraction],\allowbreak{} RestrictedProof],\allowbreak{} .\allowbreak{}.\allowbreak{}.\allowbreak{}]}, \py!restriction!, \py!left!, \py!right! --- a proof of $\mathit{left}=\mathit{right}$ under a restriction with finitely many values: one \py!RestrictedProof! per value combination of its finite clauses.

\texttt{finite\_\allowbreak{}symbolic(\allowbreak{}proof,\allowbreak{} describe=None) -> Symbolic}
\end{signature}

\begin{algorithm}[\texttt{finite\_\allowbreak{}symbolic}]\label{alg:kernel-finite}
\begin{enumerate}
\item The restriction must have finite clauses (\texttt{the restriction has no clause with finitely many values}). The instance points must be exactly the value combinations of the finite clauses, in \py!itertools.product! order over the clauses' values (\texttt{the instances are not the value combinations of T,\allowbreak{} .\allowbreak{}.\allowbreak{}.\allowbreak{}}).
\item Each instance must be restricted to the restriction's \py!instance_at(point)! (\texttt{an instance is not restricted to P}) and be about the same pair, id for id (\texttt{an instance is about another pair}). Its source is \texttt{restricted\_\allowbreak{}symbolic(\allowbreak{}inner).\allowbreak{}source} for a normalizer proof, else \py!lean(check(inner))!.
\item A finite clause's values are \emph{integral} when it is a multiplicity or dimension clause and every value is a nonnegative integer. Combine the sources by a case split over the values (\texttt{lean\_\allowbreak{}finite.\allowbreak{}combine(\allowbreak{}sources,\allowbreak{} integral,\allowbreak{} restriction text,\allowbreak{} \_\allowbreak{}var,\allowbreak{} shared)}, \cref{ch:lean-export}), \code!shared! true when every instance is a family proof; a \py!FiniteUnsupported! refuses with \code!no combined export: E!.
\item Return \texttt{Symbolic(\allowbreak{}left,\allowbreak{} right,\allowbreak{} source,\allowbreak{} restriction.\allowbreak{}clauses)}.
\end{enumerate}
\end{algorithm}

\texttt{\_\allowbreak{}symbolic\_\allowbreak{}shift(\allowbreak{}proof,\allowbreak{} describe)} re-proves the claim over the original diagrams with the normalizer. Let the substituting clauses split into dimension equalities and multiplicity clauses. It gives \py!None! when there is no substituting clause, when a multiplicity clause is of another kind or is affine with $a\neq1$, or when the only multiplicity clause is a lower bound $n\ge c$ with $c<1$. Otherwise it starts from \texttt{dimension\_\allowbreak{}floors(\allowbreak{}left,\allowbreak{} right)} of \cref{ch:proofs}, raised by the restriction's floors that the proof used; applies the dimension equalities to both sides and drops their floors; and calls \py!lean_sym.export_symbolic! with those floors, the shift $(n,c)$ when the only multiplicity clause is a lower bound, the original diagrams with the dimension values when there are dimension equalities, and otherwise the list of multiplicity clauses $(n, c, \text{is equality})$ as \code!fixed!. With lower-bound clauses it appends \py!shift_corollary!. Any \py!NetUnsupported!, \py!ScalarUnsupported!, \py!RestrictUnsupported!, or \py!KernelError! gives \py!None!.

\section{Refusals}\label{sec:kernel-refusals}

\Cref{tab:kernel-refusals} lists every \py!KernelError! of the kernel with its condition. $T$ is a clause text, $E$ an exception text, $V$ a value, and quoted names are Python \py!repr!s.

\begin{small}
\begin{longtable}{@{}>{\raggedright\arraybackslash}p{0.55\linewidth}>{\raggedright\arraybackslash}p{0.4\linewidth}@{}}
\caption{Refusals of the proof kernel.}\label{tab:kernel-refusals}\\
\toprule
Message & Condition \\
\midrule
\endfirsthead
\toprule
Message & Condition \\
\midrule
\endhead
\bottomrule
\endfoot
\texttt{a Theorem is built only by the kernel's inference rules} & the constructor got another token \\
\texttt{the box multiplicity M is negative for large values} & a top-degree coefficient of a multiplicity is negative \\
\texttt{the box multiplicity M is negative and no floor of one symbol helps} & no floor of \cref{alg:kernel-mult-floor} \\
\texttt{W: a premise holds only under T} & a restricting clause names a quantified or substituted symbol ($W$ the rule) \\
\texttt{the pure rule R consumes a discard} & \textsc{Rule} of a pure rule at a discard \\
the replay failure of \texttt{\_replay\_one} & \textsc{Rule}: lookup, a missing node or wire, an unmatched match, an apply error, or a different step \\
\texttt{the diagrams have different comparison views} & \textsc{Iso} \\
\texttt{an algebraic meet of diagrams with bang boxes} & \textsc{Algebraic} \\
\texttt{a spider parameter at a symbolic dimension} & \textsc{Algebraic} \\
\texttt{no number-field certificate: E} & \textsc{Algebraic}: certificate construction, check, or key failed \\
\texttt{the diagrams differ beyond equal algebraic values} & \textsc{Algebraic}: keyed copies not isomorphic or scalars differ \\
\texttt{no symbolic contraction: E} & \textsc{Contraction}: unsupported, domain, or budget \\
\texttt{the symbolic contractions differ: R} & \textsc{Contraction} \\
\texttt{no path sum: E}; \texttt{the path sums are not equal: R} & \textsc{PathSum} \\
\texttt{repeel: E}; \texttt{peel: E} & the joint peel raised \\
\texttt{regroup: the boxes differ}; \texttt{regroup: the sides differ beyond box groups}; \texttt{regroup: the sides lay out their slots differently} & \textsc{Regroup} \\
\texttt{identity\_slide: \ldots}, \texttt{box\_split: \ldots}, \texttt{leg\_move: \ldots}, \texttt{power\_split: \ldots} & no candidate is identical to the stated result \\
\texttt{w\_fusion raised T: E} & \textsc{WFusion} \\
\texttt{align\_merge: \ldots} (six group messages); \texttt{align\_merge: the diagram is no aligned box merge of the other} & \textsc{AlignMerge} (group messages only reached through \texttt{\_align\_group}; \texttt{align\_merge} itself skips such a root) \\
\texttt{chain\_split: \ldots} (eight group messages); \texttt{chain\_split: the diagram is no chain box split of the other} & \textsc{ChainSplit} (likewise) \\
\texttt{'a' is not a phase parameter of the claim}; \texttt{'a' is read outside node phases}; \texttt{V is not a real value free of the copy and phase indices}; \texttt{the parameter 'a' or its value reads the induction symbols} & \textsc{InductionParam}, \textsc{PhaseInstance} \\
\texttt{phase\_instance raised T: E} & \textsc{PhaseInstance} \\
\texttt{instance: V is not a multiplicity}; \texttt{instance: 'x' is not a free multiplicity}; \texttt{instance: a symbol is a copy index}; \texttt{instance raised T: E} & \textsc{Instance} \\
\texttt{the region is not the hypothesis diagram through cut wires}; \texttt{the part off the region: E}; \texttt{the regions sit in different places of the two diagrams}; \texttt{the parts off the regions have different interfaces} & \textsc{LocalContraction} \\
\texttt{congr: \ldots} (five messages); \texttt{congr raised T: E} & \textsc{Congr} \\
\texttt{trans: the middle diagrams differ} & \textsc{Trans} \\
\texttt{'n' is not a free multiplicity or 'k' is not fresh} and the seven window messages; \texttt{rebuilding the induction raised T: E} & \textsc{Induction} \\
\texttt{the offset 's' is not fresh}; \texttt{a side at 0 depends on the offset} & \textsc{InductionOffset} \\
\texttt{'n' is not a free multiplicity or there is no base theorem}; \texttt{a base theorem holds under a restriction}; \texttt{base r does not start/end \ldots}; \texttt{periodic: E}; \texttt{rebuilding a base raised T: E}; \texttt{periodic: the sides scale by A and B} & \textsc{Periodic} \\
\texttt{subst: T substitutes nothing}; \texttt{subst: m is not fresh for the theorem}; \texttt{subst: T leaves the naturals}; \texttt{subst raised T: E}; \texttt{subst: a stated side is not the theorem's side at T} & \textsc{Subst} \\
\texttt{restricting raised T: E}; \texttt{the theorem is not about the restricted pair} & restricted reading \\
\texttt{shift: T is not n = a*m + c with a > 1 or c >= 1}; \texttt{'n' is not a free multiplicity or 'm' is not fresh}; \texttt{rebuilding the successor raised T: E} & \textsc{Shift} \\
\texttt{lift: T is not a multiplicity clause}; \texttt{'n' is not a free multiplicity}; \texttt{lift: T has no fresh symbol}; \texttt{lift: one clause per multiplicity}; \texttt{rebuilding the instance raised T: E} & \textsc{Lift} \\
\texttt{a family proof holds only under a restriction}; \texttt{an induction certificate carries no re-checkable step}; \texttt{the steps do not reach the recorded final diagram}; \texttt{a proof half does not begin at its end of the goal}; \texttt{cannot check a T} & \texttt{check} \\
\texttt{cannot read a T}; \texttt{a doubled verdict without its doubled pair and settling rung}; \texttt{the doubled pair is not the doubling of left and right}; \texttt{a saturation edge does not continue the path}; \texttt{a M verdict carries no re-checkable proof} & \texttt{from\_decision} \\
\texttt{the decision carries no counterexample} & \texttt{refutation} \\
\texttt{the restricted proof is not a normaliser proof}; \texttt{no corollary: E}; \texttt{no symbolic proof: E} & \texttt{restricted\_symbolic}, \texttt{symbolic}, \texttt{lean} \\
\texttt{no Lean refutation: E}; \texttt{lean takes a Theorem, got T} & \texttt{lean} \\
\texttt{the restriction has no clause with finitely many values}; \texttt{the instances are not the value combinations of T}; \texttt{an instance is not restricted to P}; \texttt{an instance is about another pair}; \texttt{no combined export: E} & \texttt{finite\_symbolic} \\
\texttt{no Lean natural for E}; \texttt{a shift has no opaque export}; \texttt{an K inference has no opaque export} & the opaque export \\
\end{longtable}
\end{small}

\section{Export to Lean}\label{sec:kernel-lean}

\begin{signature}
\texttt{lean(\allowbreak{}theorem,\allowbreak{} name="main",\allowbreak{} describe=None,\allowbreak{} *,\allowbreak{} semantic=True,\allowbreak{} guards=None) -> str}

Lean~4 source stating and proving \code!theorem! (a \py!Theorem!, \py!Refutation!, or \py!Symbolic!) under the declaration name \code!name!. \code!describe! gives each diagram's doc comment (default \py!_summary!: \texttt{a diagram of N nodes and M wires}); \code!semantic=False! forces the opaque export; \code!guards! (symbol to least value) adds the hypotheses \code!lo ≤ n! to the opaque statement, to the instantiation axioms of substitutions $s=n-c$, and to each lemma and step axiom whose theorem holds only from a floor the enclosing hypotheses give, each discharged by \lean!omega!.
\end{signature}

\begin{algorithm}[\texttt{lean} dispatch]\label{alg:kernel-lean}
\begin{enumerate}
\item A \py!Symbolic!: with \code!name="main"!, its source; otherwise \texttt{lean\_\allowbreak{}sym.\allowbreak{}export\_\allowbreak{}symbolic(\allowbreak{}lhs,\allowbreak{} rhs,\allowbreak{} name,\allowbreak{} describe)} (a \py!NetUnsupported! refuses with \code!no symbolic proof: E!). The re-export under another name proves the bare pair; the restriction header, corollary, and floors of a restricted \py!Symbolic! are not carried.
\item A \py!Refutation!: with a source and \code!name="main"!, the source; otherwise \texttt{lean\_\allowbreak{}refute.\allowbreak{}export\_\allowbreak{}refutation(\allowbreak{}lhs,\allowbreak{} rhs,\allowbreak{} point,\allowbreak{} name,\allowbreak{} describe)} (\cref{ch:lean-export}; \code!no Lean refutation: E! on \py!NetUnsupported!).
\item Anything but a \py!Theorem!: \code!lean takes a Theorem, got T!.
\item A theorem with a nonempty domain: the restricted export (\cref{sec:kernel-lean-restricted}).
\item With \code!semantic! and no guards: try the tensor-network export \texttt{lean\_\allowbreak{}net.\allowbreak{}export(\allowbreak{}theorem,\allowbreak{} name,\allowbreak{} describe)} (\cref{ch:lean-export-net}); a \py!NetUnsupported! leaves no source. A source without a line beginning \code!axiom unproved_! is returned. Otherwise try \texttt{lean\_\allowbreak{}sym.\allowbreak{}export\_\allowbreak{}symbolic(\allowbreak{}lhs,\allowbreak{} rhs,\allowbreak{} name,\allowbreak{} describe)} and return it unless it raises \py!NetUnsupported!. Otherwise return \texttt{lean\_\allowbreak{}adder.\allowbreak{}export\_\allowbreak{}family(\allowbreak{}lhs,\allowbreak{} rhs,\allowbreak{} variables,\allowbreak{} name)} when it is not \py!None! (\cref{sec:lean-export-adder}). Otherwise return the network source when there is one.
\item Otherwise the opaque export (\cref{sec:kernel-opaque}).
\end{enumerate}
\end{algorithm}

\subsection{The opaque export}\label{sec:kernel-opaque}

The opaque export writes every diagram as a constant of an opaque type and every kernel-checked step as an axiom, so that Lean checks the proof's composition. A builder (\py!_Lean!) accumulates declarations while it translates the proof tree; when a step has a semantic proof (with \code!semantic!), that proof is nested as a sub-module and linked to the step by one axiom.

\paragraph{Names.} \py!_var(name)! makes a Lean binder: every character outside \code![A-Za-z0-9_]! becomes \code!_!, and \code!v_! is prefixed when the result does not begin with a letter or underscore or is one of \code!at!, \code!axiom!, \code!by!, \code!def!, \code!do!, \code!else!, \code!end!, \code!from!, \code!fun!, \code!have!, \code!if!, \code!in!, \code!let!, \code!match!, \code!show!, \code!then!, \code!theorem!, \code!with! (\py!_LEAN_KEYWORDS!). \py!_nat_expr(e)! writes a multiplicity as a \lean!Nat! term: a nonnegative integer as itself, a symbol by \py!_var!, a sum or product as \code!(a + b)! or \code!(a * b)!, a power as \code!(b ^ e)!; anything else refuses with \code!no Lean natural for E!. Doc comments pass through \py!_comment!, which writes \code!-/! as \code!- /! and \code!/-! as \code!/ -!. Fresh declaration names are a stem followed by one counter shared by all stems (\code!ax1!, \code!iso2!, \code!hole3!, \ldots); diagram constants are \code!d0!, \code!d1!, \ldots{} in order of first use.

\paragraph{Diagram constants.} Each distinct diagram (identical id for id, with the same sorted free multiplicities) is declared once as \texttt{axiom d<i> : Nat $\to$ .\allowbreak{}.\allowbreak{}.\allowbreak{} $\to$ Diagram}, one \lean!Nat! per free multiplicity, with the doc comment \code!describe(d)!. A diagram is referenced as \code!d<i>! or \code!(d<i> a b)!, its arguments the binders of its multiplicities or, under a substitution, the given terms.

\paragraph{Step axioms.} \texttt{axiom(\allowbreak{}stem,\allowbreak{} note,\allowbreak{} binders,\allowbreak{} lhs,\allowbreak{} rhs)} returns \lean!rfl! when the two texts coincide; otherwise it declares (once per distinct statement text) \texttt{axiom <stem><i> : $\forall$ (\allowbreak{}b1 b2 : Nat),\allowbreak{} lo $\le$ b1 $\to$ lhs = rhs} with the doc comment \code!note!, where the hypotheses are the current step floors whose variables are among the binders, and returns \code!(<stem><i> b1 b2 (by omega))! with one \code!(by omega)! per hypothesis (the bare name without binders).

\begin{table}[H]
\centering\small
\begin{tabularx}{\linewidth}{@{}>{\raggedright\arraybackslash}p{0.3\linewidth}>{\raggedright\arraybackslash}X@{}}
\toprule
Inference & Opaque Lean term \\
\midrule
\texttt{refl} & \texttt{rfl} \\
\texttt{rule} & axiom \texttt{ax<i>} with note \texttt{rule R at n3, n5}; for a rule whose name begins \texttt{discard\_}, the stem \texttt{unproved\_DISCARD\_<rest>\_} \\
\texttt{iso} & axiom \texttt{iso<i>}, note \texttt{isomorphic comparison views} \\
\texttt{contraction} & with \texttt{semantic}: a nested semantic proof (\cref{sec:kernel-opaque-semantic}); else axiom \texttt{contraction<i>}, note \texttt{symbolic contractions agree, every symbol formal} \\
\texttt{path\_sum} & a nested semantic proof of the two sides; else axiom \texttt{path\_sum<i>}, note \texttt{path sums reduce to one form} \\
\texttt{peel}, \texttt{identity\_slide}, \texttt{align\_merge}, \texttt{power\_split} & with \texttt{semantic}: a nested linked-diagram proof (below); else an axiom of that stem \\
\texttt{repeel} & with \texttt{semantic}: \texttt{(pa).symm.trans (pb)} from two nested peel proofs of its source onto each side (when the source's scalar equals the right side's and both peels export); else an axiom \\
\texttt{box\_split}, \texttt{leg\_move}, \texttt{chain\_split}, \texttt{w\_fusion} & axiom of that stem \\
\texttt{algebraic} & axiom \texttt{algebraic<i>}, note \texttt{equal values in a checked number field} \\
\texttt{symm} & \texttt{(p).symm} (\texttt{rfl} stays \texttt{rfl}) \\
\texttt{trans} & \texttt{(p).trans (q)}, a \texttt{rfl} operand dropped \\
\texttt{subst}, \texttt{instance} & a lemma \texttt{fam<i>} over the premise, applied at the value, between instantiation axioms \\
\texttt{congr}, \texttt{regroup} & the premise under \texttt{congrArg} of a context axiom, between plug or layout axioms \\
\texttt{induction} & a theorem \texttt{ind<i>} by Lean's \texttt{induction} tactic \\
\texttt{shift} & refused: \texttt{a shift has no opaque export} \\
\texttt{induction\_offset}, \texttt{induction\_param}, \texttt{induction\_offset\_param}, \texttt{phase\_instance}, \texttt{periodic} & an axiom named after the kind over the two sides \\
\bottomrule
\end{tabularx}
\caption{Translation of inferences by the opaque export.}\label{tab:kernel-opaque}
\end{table}

The notes of the box axioms are: \code!repeel! \texttt{two joint peels of one box family}; \code!peel! \texttt{a joint peel of one box family}; \code!identity_slide! \texttt{a boxed identity spider moved onto another wire}; \code!box_split! \texttt{a boxed spider unfused in every copy}; \code!align_merge! \texttt{copy-aligned boxes of one count merged into one box}; \code!chain_split! \texttt{a chain box over independent chains split into aligned chains}; \code!leg_move! \texttt{legs moved between two wired spiders of one type}; \code!power_split! \texttt{a box of b\^{}(\allowbreak{}e + 1) copies split into b boxes of b\^{}e copies}; \code!w_fusion! \texttt{copies of a W node fed by one W output leg fused into it}.

\paragraph{Linked-diagram proofs.} With \texttt{semantic}, a \texttt{peel} goes through \py!peel_net! of its two sides; an \texttt{identity\_slide} through \py!lean_linked.export_slide! (equal scalars required), an \texttt{align\_merge} through \py!lean_merge.export_merge!, and a \texttt{power\_split} through \py!lean_split.export_split! (\cref{ch:lean-export}), each source handed to \py!linked_net(source, stem, tail, binders, lhs, rhs)! with the stem \code!S!, \code!M!, and \code!Q!. \py!linked_net! nests the source (\py!_nest!, namespace stem plus a running number), adds its imports, and declares
\begin{lstlisting}[style=lean]
axiom net<i> : (type_of% @NS.main) → ∀ (binders : Nat), LHS = RHS
\end{lstlisting}
with the note \texttt{the diagrams are the linked diagrams NS.L and NS.L' times one scalar, whose digit sums NS.main proves equal,} followed by the inference's tail: for a slide, the ghosts \texttt{NS.gsL}, \texttt{NS.gsL'}, the box groups, the flat boundary positions, and the zero phase index; for a merge, \texttt{LDiagram.evalP} with the pair boxes \texttt{NS.ps}, the box index read through \texttt{cph} and \texttt{cdim}, the aligned wires, the box groups, and the flat positions; for a split, \texttt{LDiagram.evalX} with the same readings and the copies from the first box's copy count on. The term is \texttt{(net<i> @NS.main binders)}. A refused export (\py!None!) falls back to the axiom of that stem.

\paragraph{Floors.} The floors of a theorem (\texttt{\_\allowbreak{}floors\_\allowbreak{}of(\allowbreak{}domain,\allowbreak{} variables)}) are, per variable, its least value above $0$: a multiplicity side clause's lower bound, or $-c$ (at least the fresh range's bound) for the fresh $m$ of a shift $s=m+c$ with $a=1$ and $c<0$. The builder keeps \code!guards! (the floors the enclosing hypotheses hold) and \code!step! (the floors the current step axioms carry). \py!term(t, floors=None)! sets \code!step! to the theorem's floors (or the given ones) that \code!guards! meets and translates $t$. \py!floors(t, values)! returns the theorem's floors when each holds at its value (a linear expression with nonnegative coefficients, evaluated at the guards' values, default $0$, by \py!at_least!), else none.

\paragraph{Substitution and instance.} For \code!subst! the premise is stated as \texttt{theorem fam<i> : $\forall$ (\allowbreak{}vars : Nat),\allowbreak{} hyps,\allowbreak{} lhs = rhs := fun vars h => body} and applied at the clause's value (\code!c! for an equality, \code!(a * m + c)! for an affine clause, \code!(n - c')! for a shift $s=n-c'$ under a guard on $n$). The term is \texttt{(\allowbreak{}inst\_\allowbreak{}l).\allowbreak{}symm.\allowbreak{}trans (\allowbreak{}(\allowbreak{}fam args).\allowbreak{}trans (\allowbreak{}inst\_\allowbreak{}r))}, with instantiation axioms \code!inst<i>! (note \code!instantiation at T!) from the premise's sides at the value to the conclusion's sides; under a guard the instantiation axioms carry the hypothesis \code!lo ≤ n! and receive \code!h_n! (\py!guarded!). \code!instance! is the same with the value written by \py!_nat_expr! and the note \code!instantiation at x = e!.

\paragraph{Context rules.} For \code!congr! the builder declares \texttt{axiom hole<i> : Nat $\to$ .\allowbreak{}.\allowbreak{}.\allowbreak{} $\to$ Diagram $\to$ Diagram} (doc \code!a context around a region!), and the plug axioms \code!lhs = (hole lhs')! (note \code!the region!) and \code!rhs = (hole rhs')! (note \code!the region replaced!); the term is \texttt{(\allowbreak{}plug\_\allowbreak{}l).\allowbreak{}trans (\allowbreak{}(\allowbreak{}congrArg ctx (\allowbreak{}p)).\allowbreak{}trans (\allowbreak{}plug\_\allowbreak{}r).\allowbreak{}symm)}. \code!regroup! declares \code!slots<i>! (doc \texttt{one permutation of the instantiated boundary slots,\allowbreak{} from the box groups}) and per side either a nested layout proof or an axiom \code!lay<i>! (note \texttt{the left side,\allowbreak{} its slots laid out by its box groups}), composed the same way.

\paragraph{Induction.} An \code!induction! inference becomes the theorem below (named \code!ind<i>!, or the export's \code!name! at top level), over the other variables \textit{others} and the index $n$, under the floors of the other variables; with the step symbol $k$:
\begin{lstlisting}[style=lean]
theorem ind7 : ∀ (others n : Nat), <floor hyps> lhs = rhs := by
  intro others n h_...
  induction n with
  | zero => exact (inst_s).trans ((BASE).trans (inst_g).symm)
  | succ k ih => exact (succ_own).trans ((EXPOSE).trans ((plug_own).trans
      ((congrArg ctx ih).trans ((plug_other).symm.trans
      ((STEP).trans (succ_other).symm)))))
\end{lstlisting}
\code!inst_s! and \code!inst_g! are instantiation axioms (note \code!n = 0!) from the sides at \code!0! to the base premise's sides; \code!succ_own! and \code!succ_other! state each side at \code!(k + 1)! equal to its peeled successor, by a nested linked-diagram peel proof when one exists and no step floor binds, else by an axiom \code!peel<i>! (note \code!n = k + 1!); \code!hole<i>! (doc \texttt{the exposed successor around a hole}) is the context \code!ctx!; \code!plug_own! states the exposed diagram equal to \code!ctx! of the own side at $k$ (note \code!the hypothesis region!) and \code!plug_other! the replaced diagram equal to \code!ctx! of the other side at $k$ (note \code!the region replaced!). For side \texttt{GOAL}, \code!ih! becomes \code!ih.symm! and the whole successor term is wrapped as \code!(...).symm!. \code!BASE! is a nested semantic proof of the base premise when \code!semantic! holds, the premise is not \code!refl!, and its sides have no free multiplicity; otherwise its term. The induction is referenced as \texttt{(\allowbreak{}ind7 others n (\allowbreak{}by omega).\allowbreak{}.\allowbreak{}.\allowbreak{})}.

\paragraph{Top level and header.} At top level, an \code!induction! theorem is emitted under \code!name!; any other theorem becomes \texttt{theorem name : $\forall$ (\allowbreak{}vars : Nat),\allowbreak{} lo $\le$ v $\to$ .\allowbreak{}.\allowbreak{}.\allowbreak{} lhs = rhs := fun vars h\_\allowbreak{}v => term} (without binders when it has no variable), the hypotheses from \code!guards!. The source is the header, a blank line, the declarations separated by blank lines, a blank line, and the closing line. Without nested proofs the header is
\begin{lstlisting}[style=lean]
/-
  ArchytasZX proof export. Each axiom states one rewrite-rule instance, isomorphism,
  algebraic meet, symbolic contraction, instantiation, peeling or context plug the
  ArchytasZX kernel checked; Lean checks how they compose.
  assumes <assumption>
-/

namespace ArchytasZX

axiom Diagram : Type
\end{lstlisting}
with one \code!assumes! line per assumption, and the closing line is \code!end ArchytasZX!. With nested proofs the comment reads \texttt{Each axiom states one rewrite-rule instance,\allowbreak{} isomorphism,\allowbreak{} algebraic meet,\allowbreak{} instantiation,\allowbreak{} peeling or context plug the ArchytasZX kernel checked; each symbolic contraction is a theorem C<i>.\allowbreak{}main over the networks of the two contracted diagrams,\allowbreak{} carried to the step by the axiom net<i> (\allowbreak{}the diagrams' translation into those networks); Lean checks how they compose.\allowbreak{}}, followed by the collected imports, the nested proofs, and the namespace \code!ArchytasZX.Opaque! (\py!OPAQUE_NS!), closed by \code!end ArchytasZX.Opaque!.

\begin{example}[An opaque export]
Two phaseless qubit Z spiders in sequence, $d$, fuse to one by \textsc{Rule}$(\sigma)$; \texttt{trans(\allowbreak{}rule(\allowbreak{}d,\allowbreak{} $\sigma$),\allowbreak{} symm(\allowbreak{}rule(\allowbreak{}d,\allowbreak{} $\sigma$)))} has size $4$, axioms \texttt{['spider\_\allowbreak{}fusion',\allowbreak{} 'spider\_\allowbreak{}fusion']}, and an empty statement. \py!lean(t, semantic=False)! prints:
\begin{lstlisting}[style=lean]
/-
  ArchytasZX proof export. Each axiom states one rewrite-rule instance, isomorphism,
  algebraic meet, symbolic contraction, instantiation, peeling or context plug the
  ArchytasZX kernel checked; Lean checks how they compose.
-/

namespace ArchytasZX

axiom Diagram : Type

/-- a diagram of 2 nodes and 1 wires -/
axiom d0 : Diagram

/-- a diagram of 1 nodes and 0 wires -/
axiom d1 : Diagram

/-- rule spider_fusion at n0, n1 -/
axiom ax1 : d0 = d1

theorem main : d0 = d0 :=
  (ax1).trans ((ax1).symm)

end ArchytasZX
\end{lstlisting}
\end{example}

\subsection{Nested semantic proofs}\label{sec:kernel-opaque-semantic}

\py!_semantic_pair(x, y)! is Lean source whose theorem \lean!main! proves $x=y$ over their networks: the first of \texttt{lean\_\allowbreak{}sym.\allowbreak{}export\_\allowbreak{}symbolic(\allowbreak{}x,\allowbreak{} y,\allowbreak{} "main")} (\cref{ch:lean-sym}), \texttt{lean\_\allowbreak{}pathsum.\allowbreak{}export\_\allowbreak{}net(\allowbreak{}x,\allowbreak{} y)}, and \texttt{lean\_\allowbreak{}pathsum.\allowbreak{}export\_\allowbreak{}qudit\_\allowbreak{}net(\allowbreak{}x,\allowbreak{} y)} (\cref{ch:lean-export}) that raises none of \py!NetUnsupported!, \py!ScalarUnsupported!, \py!PathSumError!, \py!ValueError!; \py!None! when all raise. \py!_nest(source, tag)! removes the source's \code!import! lines (returned separately), renames its first namespace \code!N! to \code!N.tag! in its \code!namespace! and \code!end! lines, and makes every \code!macro "! line \code!local!; it returns the body, the imports, and the new namespace.

A step with a nested proof is linked by one axiom \texttt{axiom net<i> : (\allowbreak{}type\_\allowbreak{}of\% @NS.\allowbreak{}main) $\to$ $\forall$ (\allowbreak{}vars : Nat),\allowbreak{} lhs = rhs}, with a doc comment that names what the translation asserts, and the term \code!(net<i> @NS.main vars)!. The tags are \code!C<j>! for a semantic pair, \code!P<j>! for a peel (\py!lean_linked.export_peel!), \code!S<j>! for an identity slide (\py!lean_linked.export_slide!), \code!Q<j>! for a power split (\py!lean_split.export_split!), and \code!R<j>! for a regroup layout (\py!lean_regroup.export_regroup!), $j$ one more than the number of nested proofs so far. A peel or slide proof requires equal scalars on its two sides.

\paragraph{Contraction steps over !-boxes.} For a contraction whose diagrams (the parts off the regions for a local contraction) carry !-boxes or free multiplicities, \py!inducted! re-proves the step by induction: when the nesting depth is below $2$ (\py!INDUCTED_DEPTH!), \py!_contraction_induction! gives a theorem (and possibly a shift $(n, m, c)$); its term becomes a lemma \code!fam<i>! over its variables under its floors; with a shift, the lemma is applied at \code!(n + c)! between instantiation axioms (note \code!instantiation at m = n + c!). For a local contraction, the result is carried to the step by an axiom \code!net<i>! whose doc comment states that the diagrams are the parts off the two regions, plugged into one context, which the lemma proves equal. Without a theorem the step falls back to the \code!contraction<i>! axiom.

\subsection{Restricted theorems}\label{sec:kernel-lean-restricted}

\begin{algorithm}[\texttt{\_\allowbreak{}lean\_\allowbreak{}restricted}]\label{alg:kernel-lean-restricted}
\begin{enumerate}
\item \emph{Guards.} \py!_guards(theorem)! is defined when the theorem has no origin, and every domain clause is either a multiplicity shift $s=m+c$ with $a=1$, $c<0$, and $m$ a variable, or a multiplicity side clause on a variable whose domain is the integers from a positive lower bound with nothing else (\py!_floor_only!); it then gives \py!_floors_of!. With guards, export the theorem without domain (no floors) by \py!lean(..., guards=guards)! and put the header note \texttt{for n >= 3: T by instantiation axioms under the hypotheses n >= 3}, or, when no clause's fresh symbol is guarded, \texttt{for n >= 3: the box floors,\allowbreak{} the hypotheses n >= 3 of the theorem and of its lemmas and steps that hold only from them}.
\item \emph{Ranged family.} Strip \code!subst! inferences until the theorem's sides are exactly (\py!is!) the origin pair, or no \code!subst! is left; call the result the family base. The \emph{ranged} clauses are the multiplicity shifts $s=m+c$ with $a=1$, $c<0$, and a fresh domain (\py!_guard_shift!); the \emph{own} clauses are the rest. When the family base lies on the origin, both groups are nonempty, and the family base restated with only the ranged clauses as its domain has guards, export that restated theorem, put the note \texttt{restricted to T\_\allowbreak{}own: the family theorem,\allowbreak{} then its corollary at the clauses}, and append the corollary at the own clauses (\code!no corollary: E!).
\item Otherwise strip every \code!subst! to the base theorem and export it without domain (\emph{plain}). The theorem is a \emph{family} theorem when the origin is exactly the base theorem's sides. When the plain export refuses, a family theorem with \code!semantic! falls back to the concrete export (below), and otherwise the refusal propagates. For a family theorem whose plain source, with \code!semantic!, has a line beginning \code!axiom ! (an opaque export), the concrete export is returned when it exists. Otherwise a family theorem gets the note \texttt{restricted to T: the family theorem,\allowbreak{} then its corollary at the clauses} and the appended corollary (\py!lean_restrict.corollary!; \code!no corollary: E!).
\item Otherwise, with \code!semantic!, try \py!_lean_shift!; with a result, the note is the pair note followed by \code!; ! and the lifted note. Else the note is the pair note followed by \code!; UNLINKED! text when a clause substitutes.
\end{enumerate}
\end{algorithm}

The \emph{pair note} (\py!_pair_note!) is \texttt{restricted to T: the theorem about the restricted pair}, followed by \code!, substituted x = e, ...! for the substituting clauses. The \emph{lifted note} (\py!_lifted_note!) joins \texttt{stated over the original diagrams by peeling (\allowbreak{}Equal.\allowbreak{}of\_\allowbreak{}split\_\allowbreak{}at)} (when a substituting clause is not a dimension equality) and \texttt{restated over the original diagrams with d free (\allowbreak{}main\_\allowbreak{}diagrams\_\allowbreak{}original)} (when there are dimension equalities). \py!UNLINKED! is \texttt{UNCERTIFIED for the original diagrams: the theorem is about the restricted diagrams,\allowbreak{} their link to the original diagrams not proved}.

\texttt{\_\allowbreak{}lean\_\allowbreak{}shift(\allowbreak{}plain,\allowbreak{} theorem,\allowbreak{} name,\allowbreak{} describe)} carries a restricted-pair theorem to the original diagrams: it requires an origin and a substituting clause, and every non-dimension substituting clause to be a multiplicity clause, and declines a single clause $n=m$. It applies the dimension equalities to the original pair, then states the original pair by \textsc{Shift} at a single non-equality multiplicity clause, by \textsc{Lift} at several multiplicity clauses or equality clauses, or by \py!_onto! with dimension equalities only; it exports that theorem by \py!lean_net.export! (with the original diagrams and dimension values when there are dimension equalities) and appends \py!shift_corollary! for the non-equality clauses. A \py!KernelError!, \py!NetUnsupported!, \py!ScalarUnsupported!, or \py!RestrictUnsupported! gives \py!None!.

\paragraph{Concrete export.} \texttt{\_\allowbreak{}concrete\_\allowbreak{}export(\allowbreak{}theorem,\allowbreak{} text)} proves a restricted theorem at its instantiated sides when the restriction makes them concrete. When the theorem has variables or !-boxes, its sides are replaced by \py!instantiated(theorem)! (an \py!_INDUCTION_ERRORS! exception gives \py!None!). It gives \py!None! when a side still has a free multiplicity or a !-box; when the two sides have more than $8192$ nodes together (\py!CONCRETE_NODES! $=2^{13}$; the network proof grows as about $N^2/2$ bytes); when a node holds a phase rule over a concrete dimension above $65536$ (\py!WIDE_RULE! $=2^{16}$, \py!_wide_rule!, the bound up to which the parser expands a rule into a table); when both sides are qubit path sums without amplitude atoms (\texttt{from\_\allowbreak{}diagram(\allowbreak{}.\allowbreak{},\allowbreak{} w\_\allowbreak{}nodes=True)}) and the reduction of either leaves no canonical form (\py!_stuck_path_sums!, inside the span \code!lean.concrete.path_sum_gate!; a \py!PathSumError! or \py!ValueError! counts as not stuck); or when \py!_semantic_pair! of the sides is \py!None!. Otherwise it returns that source with the header note \texttt{restricted to T: the instantiated pair over its networks}. \py!_concrete_sides(lhs, rhs, text)! runs these checks from the free-multiplicity test on given sides.

\paragraph{Family and instance.} When a restricted family export (the family branch of \py!_lean_restricted!, and \py!restricted_symbolic! for a family proof) is not diagram level after its corollary, \py!_with_concrete(family, left, right, clauses)! applies every substituting clause to the original pair (\py!theorem.origin!, or the restricted proof's \py!left!, \py!right!) as \py!instantiated! does, takes \py!_concrete_sides! of the result, and, when that source is diagram level, appends it after the family source, nested by \py!_nest! as namespace \code!ArchytasZX.Export.Inst! (imports dropped, macros local). The file then holds both the every-$n$ family theorem and, last, \code!Inst.main_diagrams!, a \lean!DEqual! over the instantiated original diagrams. Otherwise the family source stays as it is, and also, before any instantiation, when a substituting multiplicity clause's value times the pair's node count exceeds \py!INSTANCE_NODES! $=2^8$ (point claims at large $n$ keep one size in $n$). When that source is missing or not diagram level and a side holds a discard \code!G! (a node with no Lean kind), \py!_with_concrete! takes \py!_concrete_sides! of \py!mixed.double! of both sides instead, with the header note \py!DOUBLED_NOTE! (``the instantiated pair over the doubled (CPM) diagrams''), and \code!Inst.main_diagrams! is a \lean!DEqual! over the doubled instantiated diagrams. The \py!INSTANCE_NODES! bound is taken on the undoubled pair. At $d=2$, $n=2$, \code!repeater_chain/main! and \code!teleportation_chain/main! get this instance.
\par\endgroup
